% concatenated from CrazyCrystals.tex
\documentclass[11pt]{amsart}
\usepackage[a4paper,margin=1in]{geometry}
\usepackage{amsmath,amssymb,amsthm,mathrsfs,mathtools}
\usepackage[colorlinks=true,linkcolor=blue,citecolor=magenta,urlcolor=blue]{hyperref}
\usepackage{ bbold }
\usepackage{dirtytalk}
\usepackage{comment}

\newcommand{\eps}{\varepsilon}
\renewcommand{\phi}{\varphi}
\newcommand{\R}{\mathbb{R}}
\newcommand{\Z}{\mathbb{Z}}
\newcommand{\N}{\mathbb{N}}
\newcommand{\diam}{\text{diam}}

\theoremstyle{plain}
\newtheorem{theorem}{Theorem}[section]
\newtheorem{lemma}[theorem]{Lemma}
\newtheorem{proposition}[theorem]{Proposition}
\newtheorem{corollary}[theorem]{Corollary}
\newtheorem{question}[theorem]{Question}

\theoremstyle{definition}
\newtheorem{definition}[theorem]{Definition}
\newtheorem{remark}[theorem]{Remark}

\numberwithin{equation}{section}

\title[Arbitrarily Fast Quantum Dispersion in Long-Range Crystals]{Arbitrarily Fast Quantum Dispersion \\ in Long-Range Crystals}

\author{Ga\'etan Leclerc, Mostafa Sabri and Tuomas Sahlsten}

\thanks{This work was concluded during the MFO Oberwolfach event \textit{Novel Developments in Spectral Theory and the Dynamics of Quantum Systems with Applications}, 9 Aug - 14 Aug 2026, we kindly acknowledge the hospitality, inspiring discussions with the other participants and facilities that helped us to complete the work. G.L. and T.S. acknowledge support from the Research Council of Finland's Academy Research Fellowship \emph{``Quantum chaos of large and many body systems''}, grant Nos. 347365, 353738. AI tools were used for language editing, proofreading and literature searches. All original ideas and proofs are those of the authors, who independently checked the final manuscript and take full responsibility for its contents.}

\begin{document}

\begin{abstract}
We construct the first examples of long-range crystals exhibiting arbitrarily fast polynomial quantum dispersion. The Floquet functions of our Hamiltonians are highly oscillatory Weierstrass functions, whose rough autosimilar structure drives the fast dispersion. The proof develops a new Fourier decay theory for $C^\alpha$ images of Lebesgue measure, based on a Dolgopyat-inspired transfer operator method, and yields a van der Corput lemma for Weierstrass functions. As a consequence, the local time of classical Weierstrass functions of sufficiently large lacunarity exists and is $C^k$, answering a question raised by Geman and Horowitz in 1980.
\end{abstract}

\maketitle

\section{Introduction}

\subsection{Background and results on dispersion} Quantum transport and dispersion describe complementary aspects of the spreading of solutions $u(t)=e^{-itH}u_0$ of the Schr\"odinger equation $i\partial_tu=Hu$. Transport measures how rapidly the mass of a wave packet moves through space, while dispersion measures its flattening as $t\to\infty$. These phenomena have been extensively studied in settings such as discrete periodic and quasiperiodic systems on $\mathbb Z$. For Schr\"odinger operators $H$ on $\mathbb Z$, transport is called \textit{ballistic} if $\|Xe^{-itH}u_0\|_{\ell^2(\mathbb Z)}\asymp t$ as $t\to\infty$ for suitably regular $u_0$, where $Xu(n)=nu(n)$ is the position observable. The operator $H$ is called \textit{dispersive} if $\|e^{-itH}\|_{\ell^1(\mathbb Z)\to\ell^\infty(\mathbb Z)}\to0$ as $t\to\infty$. See e.g. \cite{DamanikMalinovitchYoung,BoutetdeMonvelSabri} for background on these.

For locally finite periodic systems, transport is quite well understood. Periodic Schr\"odinger operators exhibit ballistic motion under broad assumptions, see e.g. \cite{DamanikLukicYessen,FillmanPeriodic,BoutetdeMonvelSabri}. Dispersion is more delicate. In one dimension, dispersive estimates go back to Firsova \cite{Firsova} in the continuum, while in the discrete setting Mi and Zhao \cite{MiZhao1,MiZhao2} obtained estimates first for period two and then for general periods. More recently, Damanik, Fillman and Young \cite{DamanikFillmanYoung} established the bound $\|e^{-itH}\|_{\ell^1(\mathbb Z)\to\ell^\infty(\mathbb Z)}\lesssim |t|^{-1/3}$ for arbitrary periodic Schr\"odinger operators on $\mathbb Z$, matching the dispersion speed of the discrete Laplacian. For quasiperiodic operators, much less is known even in dimension one, with dispersion established for sufficiently small analytic potentials \cite{BZ20}. For the Fibonacci Hamiltonian, quantum dynamics has been related to the hyperbolic dynamics and fractal geometry of the trace map, see e.g. \cite{DamanikTcheremchantsev,DamanikGorodetski,DamanikGorodetskiYessen} and the recent work \cite{Leclerc2}.

A different direction is to relax local finiteness while retaining periodicity. In \cite{KernerPostSabriTaufer}, Kerner, Post, the second author and T\"aufer studied periodic crystals with summable long-range hopping, described by operators $(Hu)(n)=\sum_{k\in\mathbb Z}w_k u(n+k)$ with $(w_k)_{k\in\mathbb Z}\in\ell^1(\mathbb Z)$. They showed that allowing infinitely many hopping distances can lead to phenomena impossible in the locally finite setting, including purely singular continuous spectrum, flat bands whose eigenfunctions must have infinite support, ballistic transport without dispersion, and transitions between ballistic and super-ballistic transport. Long-range operators also arise naturally in the study of quasicrystals through \textit{Aubry duality}. Indeed, Aubry duality transforms a one-dimensional quasiperiodic Schr\"odinger operator $(H_{V,\alpha,\theta}u)(n)=u(n+1)+u(n-1)+V(\theta+n\alpha)u(n)$ into $(\widehat H u)(n)=\sum_{k\in\mathbb Z}\widehat V(k)u(n-k)+2\cos(2\pi(\theta+n\alpha))u(n)$, see \cite{AvilaJitomirskaya,HaroPuig}. Thus analytic potentials $V$ with infinitely many nonzero Fourier coefficients give non-locally finite long-range hopping on the dual side. These examples show that removing local finiteness can fundamentally change quantum dynamics while retaining strong structure such as periodicity.

The main purpose of this paper is to show that the same is true for quantum dispersion. We construct the first crystals exhibiting \emph{arbitrarily fast polynomial quantum dispersion}, showing that there is no universal polynomial upper bound on the speed of dispersion once local finiteness is removed. The mechanism is perhaps surprising. Previously known dispersive examples in crystals had sufficiently regular Floquet functions, allowing classical oscillatory integral estimates such as stationary phase and the standard van der Corput lemma to give rates such as $|t|^{-1/3}$ in the locally finite 1D setting of \cite{DamanikFillmanYoung} and $|t|^{-1/2}$ for the long-range example of \cite{KernerPostSabriTaufer}. Thus one might expect roughness of the Floquet function to obstruct dispersion. Our examples show the opposite: sufficiently \textit{autosimilar} roughness can produce arbitrarily fast polynomial decay.

Our Floquet functions are highly oscillatory \textit{Weierstrass functions}, with fractal oscillations at every scale. Such structures have appeared in random walks with self-similar long-range jumps and models of anomalous quantum transport \cite{HughesShlesingerMontroll,CaceresNizama}, and are closely connected to hyperbolic dynamics and solenoidal attractors through the study of the Hausdorff dimension of their graphs \cite{BaranskiBaranyRomanowska,Hunt,Ledrappier,Keller,Shen,Tsujii,Romanowska}. Related Weierstrass functions also arise as temporal distance functions for Axiom A systems \cite{TZ23,Leclerc2}. Here we exploit their exact autosimilarity to obtain oscillatory cancellation through a transfer-operator argument inspired by Dolgopyat's method \cite{Do98,Na05}. This leads to a new $C^\alpha$ van der Corput lemma for Weierstrass functions that may be of independent interest.

We first state our main result on dispersion. Given $\mathbf w=(w_k)_{k\in\mathbb Z}\in\ell^1(\mathbb Z)$ with $w_{-k}=w_k$ and $w_k\ge 0$, let
$$
(H_{\mathbf w}u)(n)=\sum_{k\in\mathbb Z}w_k u(n+k),\qquad n\in\mathbb Z.
$$
This is the adjacency operator of a translation-invariant lattice in which $n$ is connected to $n\pm k$ with edge weight $w_k$. After normalization by $\|\mathbf w\|_{\ell^1}$, the weights may be interpreted as hopping probabilities.

\begin{theorem}\label{thm:maincrystal}
For every $N\in\mathbb N$ there exists $\lambda_N\geq2$ such that, for every integer $\lambda\geq\lambda_N$, the Hamiltonian $H_{\mathbf w}$ with $w_{\pm\lambda^j}=2^{-j}$, $j\geq0$, and $w_k=0$ otherwise, satisfies
$$
\|e^{-itH_{\mathbf w}}\|_{\ell^1(\mathbb Z)\to\ell^\infty(\mathbb Z)}
\lesssim_N |t|^{-N},\qquad |t|\to\infty.
$$
\end{theorem}

The construction is based on hopping along the lacunary scales $1,\lambda,\lambda^2,\lambda^3,\ldots$, with geometrically decreasing amplitudes. Increasing the lacunarity makes the corresponding Floquet function increasingly irregular, but at the same time produces increasingly strong oscillatory cancellation. In this sense, roughness is not an obstruction to dispersion in our examples, but the main mechanism behind the argument. Regarding quantitative aspects, we note that a crude tracking of the constants in the proof gives a bound $\lambda_N\leq 2^{10^9N}$, which is far from optimal.

Theorem \ref{thm:maincrystal} naturally raises the question of how far this phenomenon can be pushed. In particular, can a long-range crystal disperse exponentially fast? We show that the answer is no, and in fact this obstruction holds for every crystal, independently of hoppings $\mathbf{w} = (w_k)_{k \in \mathbb Z}$.

\begin{proposition}\label{thm:noexponential-intro}
Let $H$ be any crystal on $\mathbb Z$. Then $\langle \delta_0,e^{-itH}\delta_0\rangle$ cannot decay exponentially as $|t|\to\infty$. In particular, exponential dispersion is impossible.
\end{proposition}

The reason is purely spectral. The diagonal matrix element $\langle\delta_0,e^{-itH}\delta_0\rangle$ is the Fourier transform of the spectral measure of $\delta_0$. Since $H$ is bounded, this measure has compact support. Exponential decay of its Fourier transform would, by a Paley-Wiener argument, force its density to extend analytically to a strip, which is impossible for a non-zero compactly supported density. 

\begin{remark}
\begin{itemize}
    \item[(1)] Our results show that arbitrarily high polynomial powers are attainable, whereas exponential dispersion is universally ruled out. Whether the intermediate regime of super-polynomial but sub-exponential dispersion can occur remains open. Beurling-Malliavin multiplier theorem \cite{BM} allows constructions of functions with such Fourier decay, so potentially with even more lacunary hopping, crystals in this faster dispersive regime could also be constructed by adapting our methods.

\item[(2)] If we do not have lacunary hopping, e.g. under an eventual monotonicity assumption on the Hamiltonian coefficients $w_k$, we prove the much stronger restriction that dispersion cannot be faster than $1/t$, see Corollary \ref{cor:monotone}. Thus long-range hopping by itself is not responsible for Theorem \ref{thm:maincrystal}: the lacunary arrangement of the interactions is essential to the mechanism.
\end{itemize}
\end{remark}

\subsection{Link to Fourier decay theory and a $C^\alpha$ van der Corput lemma}
We now explain the Fourier-analytic mechanism behind Theorem~\ref{thm:maincrystal}. By Floquet theory, a translation-invariant operator $H_{\mathbf w}$ on $\mathbb Z$ is diagonalised by the Fourier transform, with scalar Floquet function
$$
h(\theta)=\sum_{k\in\mathbb Z}w_ke^{2\pi ik\theta},\qquad \theta\in\mathbb T.
$$
The Schr\"odinger evolution becomes multiplication by $e^{-ith(\theta)}$ in Fourier space, giving
\begin{equation}\label{eq:osc-int-intro}
\|e^{-itH_{\mathbf w}}\|_{\ell^1(\mathbb Z)\to\ell^\infty(\mathbb Z)}
=
\sup_{n\in\mathbb Z}
\left|
\int_0^1 e^{2\pi in\theta}e^{-ith(\theta)}\,d\theta
\right|,
\end{equation}
see \cite[Section~2.2]{KernerPostSabriTaufer}. Thus dispersion is equivalent to quantitative decay of these oscillatory integrals, uniformly in $n$.

For $n=0$, the integral is the Fourier transform of the push-forward measure $h_*d\theta$, while for general $n$ it is the Fourier transform of the complex measure $h_*(e^{2\pi in\theta}\,d\theta)$. Thus dispersion becomes a problem about Fourier decay of images of Lebesgue measure, uniformly in the Fourier mode. This also clarifies its relation to spectral type. Fourier decay of a spectral measure rules out atoms, and sufficiently fast decay implies absolute continuity, whereas slower decay may occur for singular continuous measures \cite{SiSC7}. Conversely, absolute continuity gives decay of each fixed matrix element by the Riemann-Lebesgue lemma, but not the uniform decay required for dispersion; indeed, \cite{KernerPostSabriTaufer} gives long-range crystals with absolutely continuous spectrum and ballistic transport for which dispersion fails.

This places the problem in the broader theory of Fourier decay of dynamically defined measures, which has developed rapidly in recent years (see the survey \cite{Sahlsten}). One direction establishes Fourier decay directly for classes such as Patterson-Sullivan measures \cite{BD,BakerKhalilSahlsten,LequenSahlsten}, self-conformal measures \cite{AlgomHertzWang,AlgomHertzWangLog,BakerSahlsten,BakerKhalilSahlsten,JordanSahlsten,KaufmanCF,QueffelecRamare,SahlstenStevens}, self-affine and stationary measures \cite{Li,LiENS,LiSahlstenAffine,LiSahlsten,LindenstraussVarju,Rapaport}, and equilibrium states and other hyperbolic measures \cite{LeclercBunched,LeclercHyperbolic,LeclercJulia,BakerKhalilSahlsten}, using methods from additive combinatorics and renewal theory to random walks and transfer operators. Another direction, more relevant to us and going back to Kaufman \cite{Kaufman}, exploits smooth nonlinear images of measures, where the change of variables itself creates Fourier decay. Recent results of this type for nonlinear images of self-similar measures were developed in \cite{MosqueraShmerkin,ACWW,BakerBanaji,BanajiYu,BanajiYu2,AlgomBenOvadiaHertzShannon}.

However, the Floquet functions we consider are typically only H\"older, where classical oscillatory integral estimates using derivatives are unavailable and much less is known (see e.g. \cite{Leclerc}). Indeed, for the crystals we study in Theorem~\ref{thm:maincrystal}, with $w_{\pm\lambda^j}=2^{-j}$ and all other weights zero, the Floquet function is a \textit{Weierstrass function}. More generally, for non-constant real-analytic and $1$-periodic $\varphi:\mathbb R\to\mathbb R$, $0<\mu<1$ and $\lambda\geq2$, define
\begin{equation}\label{eq:weierstrass-intro}
W_{\mu,\lambda}(\theta)=W_{\mu,\lambda}^{(\varphi)}(\theta)
=
\sum_{j=0}^{\infty}\mu^j\varphi(\lambda^j\theta),\qquad \theta\in\mathbb T^1=\mathbb R/\mathbb Z.
\end{equation}
These are $\beta$-H\"older for every $\beta<\alpha:=|\log\mu|/\log\lambda$, but satisfy \textit{autosimilarity} relation:
$$
W_{\mu,\lambda}(\theta)=\varphi(\theta)+\mu W_{\mu,\lambda}(\lambda\theta),\qquad \theta\in\mathbb T^1.
$$
Thus increasing the lacunarity $\lambda$ makes the Floquet function increasingly rough while preserving its oscillatory structure across scales. This autosimilarity relation together with the non-constancy of $\phi$ provides a substitute for differentiability and, through a Dolgopyat-type argument \cite{Do98,Na05}, allows us to obtain Fourier decay uniformly in the additional Fourier mode in \eqref{eq:osc-int-intro}.

Our main ``$C^\alpha$ van der Corput lemma'' is the following.

\begin{theorem}\label{thm:main}
Let $\varphi:\mathbb R\to\mathbb R$ be a non-constant real-analytic $1$-periodic function, let $0<\mu<1$, and let $\lambda\geq2$ be an integer.
\begin{itemize}
\item[(1)] There exists $\beta=\beta(\lambda,\mu,\varphi)>0$ such that
$$
\sup_{n\in\mathbb Z}
\left|
\int_0^1 e^{2\pi in\theta}e^{itW_{\mu,\lambda}(\theta)}\,d\theta
\right|
\lesssim |t|^{-\beta},
\qquad |t|\to\infty.
$$
\item[(2)]There exist $c(\varphi)>0$ and $\lambda_*(\mu,\varphi)\ge2$ such that, for every integer $\lambda\ge\lambda_*(\mu,\varphi)$:
$$
\sup_{n\in\mathbb Z}
\left|
\int_0^1 e^{2\pi in\theta}e^{itW_{\mu,\lambda}(\theta)}\,d\theta
\right|
\lesssim
|t|^{-c(\varphi)\frac{\log\lambda}{|\log\mu|}},
\qquad |t|\to\infty.
$$
\end{itemize}
\end{theorem}

In terms of the maximal H\"older exponent $\alpha=|\log\mu|/\log\lambda$, the exponent in part~(2) is of order $1/\alpha$, and this order is optimal by Proposition~\ref{thm:maxdis}. Thus rougher autosimilar phases give stronger Fourier decay. By \eqref{eq:osc-int-intro}, this gives Theorem~\ref{thm:maincrystal}: for $w_{\pm\lambda^j}=2^{-j}$, taking $\mu=1/2$ and $\varphi(\theta)=2\cos(2\pi\theta)$ gives $W_{\mu,\lambda}=h$ and $\alpha=\log2/\log\lambda$, so the dispersive exponent is of order $\log\lambda/\log2$ and can be made arbitrarily large as the lacunarity $\lambda$ grows. Finally, when $W_{\mu,\lambda}\in C^{1+\alpha}$, equivalently $\lambda\mu<1$, the result can also be deduced from \cite{ACWW}, combined with the Tsujii-Zhang-inspired non-concentration argument for Weierstrass functions \cite{TZ23,Leclerc2} given in Appendix~\ref{ap:B} that appears to be new. However, this does not cover the $C^\alpha$ regime responsible for fast dispersion, our Dolgopyat-type argument applies in both regimes.

The $C^\alpha$ van der Corput lemma also has a regularity interpretation. For every $n\in\mathbb Z$, consider the complex push-forward measure $\nu_n=(W_{\mu,\lambda})_*(e^{2\pi in\theta}\,d\theta)$. Theorem~\ref{thm:main} gives polynomial decay of $\widehat{\nu_n}$ uniformly in $n$. Whenever the decay exponent is larger than $k+1$, Fourier inversion implies that $\nu_n$ has a $C^k$ density. Hence increasingly rough Weierstrass phases can produce increasingly smooth image densities. For $n=0$, the measure $\nu_0=(W_{\mu,\lambda})_*d\theta$ is precisely the occupation measure of the Weierstrass function, and its density (when it exists) is the \emph{occupation density} or \emph{local time} \cite{GemanHorowitz}. We therefore obtain:

\begin{corollary}\label{cor:localtime}
For every nonconstant real analytic $\varphi$ and $0<\mu<1$, the local time of $W_{\mu,\lambda}$ exists and is $C^k$ whenever $\lambda\geq\max\{\lambda_*(\mu,\varphi),\exp((k+1)|\log\mu|/c(\varphi))\}$. In particular, this includes the classical sine and cosine Weierstrass functions for all sufficiently large $\lambda$.
\end{corollary}

This addresses a classical question of Geman and Horowitz \cite{GemanHorowitz} for Weierstrass functions of sufficiently large lacunarity. They asked whether deterministic nowhere differentiable functions such as Weierstrass functions admit local times; see also e.g. \cite{FoucheMukeru,HarangLing,ImkellerEtAl,ImkellerPamen,Buc}, where this question is revisited. See also \cite{Rezakhanlou}. There has been some progress on the problem: in \cite{BKM}, the local time exists for almost every finite-dimensional perturbation of any fixed Weierstrass function (in sense of \textit{prevalence}). However, this does not give existence for a prescribed function such as the cosine Weierstrass functions arising from our crystals. \textit{Random} Weierstrass functions with independent random coefficients also have $L^2$ local times almost surely \cite{Romanowska}, and the case of piecewise linear $\varphi$ was studied in \cite{ImkellerPamen}. On the other hand, for $\varphi(\theta)=\sin(2\pi\theta)$ the local time of $W_{1/2,2}$ does \textit{not} exist \cite{Buc}. Thus existence can fail at a fixed lacunarity, whereas Corollary~\ref{cor:localtime} shows that for every fixed analytic $\varphi$ it holds with increasing smoothness once the lacunarity is sufficiently large. This is reminiscent of the theory of Bernoulli convolutions \cite{Solomyak,Hochman,Shmerkin,Varju}, where absolute continuity is expected for all sufficiently large contraction parameters. 

The regularity phenomenon behind Corollary~\ref{cor:localtime} also has a probabilistic analogue. Kahane's work on image measures of fractional Brownian motion shows that random functions with Hurst exponent $\alpha$ can produce Fourier decay with arbitrarily large exponent as $\alpha$ becomes small, and consequently increasingly regular local times, see \cite{Kahane}. His argument exploits the Gaussian structure of fractional Brownian motion. In our deterministic setting, cancellation instead comes from the rigid autosimilarity $W(\theta)=\varphi(\theta)+\mu W(\lambda\theta)$ and oscillatory transfer-operator estimates. In this sense, Theorem~\ref{thm:main} provides a pseudorandom substitute for randomness in producing strong Fourier decay and smoothness.

\subsection{Transport and spectrum of long-range crystals}

We now discuss how our rapid dispersion relates to transport and spectral properties, see Section~\ref{sec:further} for a more detailed discussion.

Long-range hopping can produce behaviour impossible for sufficiently local Hamiltonians. In \cite{KernerPostSabriTaufer}, fractional powers of the discrete Laplacian exhibit a transition between ballistic and super-ballistic transport. Using their transport criterion, we obtain super-ballistic transport for our Weierstrass crystals when $\alpha=|\log\mu|/\log\lambda<1/2$; see Corollary~\ref{cor:transport}. Thus increasing roughness can produce both faster dispersion and super-ballistic transport.

The spectrum is simple as a set. Since $W_{\mu,\lambda}$ is continuous, $\sigma(H)=W_{\mu,\lambda}(\mathbb T)$ is a compact interval. For $\varphi(\theta)=\cos(2\pi\theta)$, its upper endpoint is $1/(1-\mu)$, while for odd $\lambda$ one has $\sigma(H)=\left[-\frac{1}{1-\mu},\frac{1}{1-\mu}\right]$, up to normalization. For even $\lambda$, determining the lower spectral edge becomes a non-trivial extremal problem for the Weierstrass function. The spectral type is more subtle. A classical property of lacunary Fourier series \cite[Vol.~II, p.~267]{Bary} implies that our Weierstrass crystals have no point spectrum for any $\mu$ and $\lambda$. In the regime of Theorem~\ref{thm:main}(2), the spectrum is in fact purely absolutely continuous: $\nu_0=(W_{\mu,\lambda})_*d\theta$ is the spectral measure of $\delta_0$, and Corollary~\ref{cor:localtime} gives it a density. For general Weierstrass functions this can fail at smaller lacunarity: the local time of $W_{1/2,2}$ with $\varphi(\theta)=\sin(2\pi\theta)$ does not exist \cite{Buc}, despite the Fourier decay of Theorem~\ref{thm:main}(1). This example does not correspond to our crystals, and whether a similar phenomenon can occur for them remains open.



\subsection{Prospects for the theory}

We conclude with some questions concerning the long-range crystals considered above. The lacunary structure of the hopping suggests several natural extensions, in particular to random and higher-dimensional models.

(a) \emph{Long-range crystals without dispersion.}
How essential the regularity of the phase is for dispersion? Consider the $1$-periodic step function $H=1$ on $[0,1/2)$ and $H=-1$ on $[1/2,1)$, and set
$
W(\theta)=\sum_{n\geq0}\mu^nH(2^n\theta).
$
Since $H(2^n\theta)$ are independent fair signs under Lebesgue measure, $W_*d\theta$ is precisely the Bernoulli convolution $\nu_\mu$, the law of $\sum_{n\geq0}\pm\mu^n$. Thus the analogy between our Weierstrass image measures and Bernoulli convolutions becomes exact, and arithmetic obstructions to dispersion appear: for Pisot $\mu^{-1}$, the classical results of Erd\"os and Salem give $\widehat{\nu_\mu}(t)\not\to0$. This phase does not itself define a crystal of the class considered here, but it raises a natural question: can one perturb or smooth this construction to obtain a genuinely long-range crystal without dispersion while retaining the arithmetic obstruction? Interestingly, some regularity already changes: in Appendix \ref{ap:large-lacunarity} we show that if the phase is continuous and non-constant analytic on some interval, while possibly constant elsewhere, then its Weierstrass occupation measures have a Fourier-decaying weak limit as $\lambda\to\infty$. This suggests a link between regularity and arithmetic structure in determining if dispersion happens.

(b) \emph{Random and multiplicative hopping.}
Our construction is deterministic, with hopping supported on the lacunary sequence $\lambda^n$ and amplitudes $\mu^n$. A natural random analogue is $w_{\lambda^n}=\mu^nX_n,$
where $(X_n)$ are independent random variables, giving a random Weierstrass Floquet function. In view of the comparison with fractional Brownian motion discussed above, it is natural to ask whether arbitrarily fast polynomial dispersion holds almost surely when $\lambda$ is sufficiently large compared with $\mu^{-1}$. Another direction is to consider multiplicative amplitudes $w_{\lambda^n}=X_1\cdots X_n$, leading to Floquet functions with a multiplicative-cascade structure. Mandelbrot cascade measures are known to exhibit non-trivial polynomial Fourier decay \cite{ChenLiSuomala}; see also \cite{GarbanVargas}. It would be interesting to determine whether analogous random multiplicative constructions of the phase produce almost-sure fast dispersion, whether the dispersive exponent can be expressed in terms of Lyapunov exponents or moments of the multiplicative process, and whether corresponding phase transitions occur in the quantum dynamics. See e.g. \cite{Romanowska} for some works in this random setting for the local time of randomised Weierstrass functions.

(c) \emph{A $C^\alpha$ van der Corput principle for fractal measures.}
Our Fourier-decay theorem concerns images of Lebesgue measure under the autosimilar Weierstrass map. It is natural to ask whether the same mechanism survives for singular source measures: \textit{if $\nu$ is a Frostman measure on $\mathbb T$, does the image $(W_{\mu,\lambda})_*\nu$ have power Fourier decay for every $0<\mu<1$ and $\lambda\geq2$?} This question was also posed by Osama Khalil to the first author. A positive answer would complement the fractal van der Corput theorem of Algom, Chang, Wu and Wu \cite{ACWW}, where self-similar measures are mapped by sufficiently non-flat smooth functions. Here the situation is in a sense reversed: the source could be an arbitrary Frostman measure, while the map itself carries the rigid autosimilar structure. A particularly natural test case is a self-similar Cantor measure invariant under multiplication by $\lambda$, and so adapted to the transfer-operator structure used in our proof.

(d) \emph{Higher-dimensional crystals.}
Our construction is one-dimensional, but the basic Fourier-analytic formulation extends naturally to $\mathbb Z^d$. A translation-invariant long-range Hamiltonian on $\ell^2(\mathbb Z^d)$ has a Floquet function $h:\mathbb T^d\to\mathbb R$, and dispersion is governed by
$$
\sup_{k\in\mathbb Z^d}
\left|
\int_{\mathbb T^d}e^{2\pi i k\cdot\theta}e^{-ith(\theta)}\,d\theta
\right|.
$$
It is natural to ask whether there are long-range crystals on $\mathbb Z^d$ with arbitrarily fast polynomial dispersion. Separable Weierstrass phases reduce essentially to the one-dimensional theory, while genuinely higher-dimensional autosimilar phases should require new ideas. The relevant problem is to control, uniformly in $k$, the multiscale geometry of regions where $t\nabla h(\theta)$ is close to $2\pi k$. These resonant regions may have non-trivial geometry and interact across different directions and scales, leading to genuinely multidimensional non-concentration problems.

\subsection*{Organisation of the paper}

The paper is organised as follows. Section~\ref{sec:mainproof} proves the Fourier-decay theorem using twisted transfer operators and Dolgopyat-type estimates, and deduces the arbitrarily fast polynomial dispersion of Theorem~\ref{thm:maincrystal}. In Section~\ref{sec:maximal-dispersion} we study the maximal possible speed of dispersion, proving in particular that exponential dispersion is impossible for any crystal. Section~\ref{sec:further} discusses transport and spectral properties. Appendix~\ref{ap:A} gives a $C^1$ van der Corput lemma and an alternative proof of the smoother case of Theorem~\ref{thm:main}; Appendix~\ref{ap:B} proves non-concentration properties for Weierstrass functions and their consequences; and Appendix~\ref{ap:large-lacunarity} studies spectral measures in the large-lacunarity limit.


\section{Proof of Theorems \ref{thm:maincrystal} and Theorem \ref{thm:main}} \label{sec:mainproof}

Theorem \ref{thm:maincrystal} follows by establishing the more general Fourier decay statement Theorem \ref{thm:main} on the Fourier decay of complex images of Lebesgue measure under the Weierstrass functions $W := W_{\mu,\lambda}$. We then aim to prove Theorem \ref{thm:main}.

\subsection{Transfer operators}\label{sec:tran}

The core idea of the proof of Theorem \ref{thm:main} is to reduce Fourier decay to quantitative Dolgopyat's estimates on some twisted transfer operator, which is an estimate coming from the study of the mixing rate of some dynamical systems. For a good introduction to Dolgopyat estimates, the interested reader can consult Naud's work \cite{Na05} or the original article by Dolgopyat \cite{Do98}. In our setting, the autosimilarity relation 
$$W(\theta) = \varphi(\theta)+\mu W(\lambda \theta)$$ 
justifies that we put some focus on the dynamical system \say{multiplication by $\lambda$}, defined on the unit circle. This expanding map has $\lambda$ inverse branches, that we will denote $(g_a)_{a \in \{ 0,\dots, \lambda-1 \}}$. More precisely, for every $\lambda \geq 2$, we introduce the alphabet $\mathcal{A}(\lambda) := \{0,\dots,\lambda-1\}$ and the maps 
\[
g_a: \mathbb{S}^1 \sim [0,1] \rightarrow [a/\lambda,(a+1)/\lambda] \subset \mathbb{S}^1\,, \qquad  g_a(\theta) := \frac{a+\theta}{\lambda}\,.
\]
We will often identify $\mathbb{S}^1$ with $[0,1]/\sim$. The multiplication by $\lambda$ satisfies the following important property: for all $h_1,h_2:\mathbb{S}^1\rightarrow \mathbb{R}$, we have
$$ \int_{\mathbb{S}^1} h_1(\lambda \theta) h_2(\theta) d\theta = \int_{\mathbb{S}^1} h_1(\theta) \mathcal{L}_{\lambda}(h_2)(\theta) d\theta ,$$
where $\mathcal{L}_\lambda h(\theta) := \lambda^{-1} \sum_{a \in \mathcal{A}(\lambda)} h(g_a(\theta))$. The operator $\mathcal{L}_\lambda : L^1(\mathbb{S}^1) \rightarrow L^1(\mathbb{S}^1)$ is called a \textit{transfer operator} associated with the multiplication by $\lambda$. Iterating this transfer operator yields

$$
\mathcal{L}_\lambda^n(h)(\theta) = \frac{1}{\lambda^n}\sum_{\mathbf{a} \in \mathcal{A}(\lambda)^n} h(g_\mathbf{a}(\theta))
$$
where $\mathbf{a} = a_n\dots a_1 \in \mathcal{A}(\lambda)^n$, and $g_{\mathbf{a}}(\theta) := g_{a_n} \circ \dots \circ g_{a_1}(\theta) = x_\mathbf{a}+ \theta \lambda^{-n} $ for $x_\mathbf{a} = g_\mathbf{a}(0)$. Notice that, if $h$ is $L$-Lipchitz, then
$$ \mathcal{L}_\lambda^n(h) = \int_0^1 h d\theta + \mathcal{O}(L \lambda^{-n}), $$
which proves the \textit{exponential mixing} property of the multiplication by $\lambda$ for Lipschitz functions:
$$ \int h_1(\lambda^n \theta) h_2(\theta) d\theta = \int_0^1 h_1 d\theta \int_0^1 h_2 d\theta + \mathcal{O}(\lambda^{-n}). $$
This is also called a Perron-Frobenius-Ruelle Theorem. A crucial observation is that the larger $\lambda$ is, the quicker this dynamical system is decorrelating the observables $h_2$ and $h_1(\lambda^n \cdot)$. This plays an important role in the fast dispersion result when $\lambda$ grows. %\textcolor{magenta}{reformulate?} \textcolor{blue}{done (what should I reformulate ?)}\textcolor{magenta}{sorry, I meant maybe reforumlating ``this is essentially the reason'', it makes it feel like an implication, I was thinking maybe something along the lines of ``plays an important role'' or ``is at the heart of'' or smth like that, but I don't have a strong opinion about this either.}\\

To handle the proof of dispersion, we need to \say{twist} our transfer operator. For a choice of (analytic) phase $\varphi:\mathbb{S}^1 \rightarrow \mathbb{R} / 2\pi \mathbb{Z}$, we define the \say{twisted transfer operator} $\mathcal{L}_{\lambda,i \phi} : C^1([0,1],\mathbb{C}) \rightarrow C^1([0,1],\mathbb{C})$ by:
$$ \mathcal{L}_{\lambda,i\phi}(h)(\theta) := \frac{1}{\lambda}\sum_{a \in \mathcal{A}(\lambda)} e^{i \phi(g_a(\theta))} h(g_a(\theta)) = \mathcal{L}_\lambda(e^{i \phi} h). $$

We will need to iterate twisted transfer operators with changing phase at each iteration. For a sequence of phases $\Phi = (\phi_n)_{n \geq 0}$, let us then define
$$ \mathcal{L}_{\lambda,i \Phi}^{\langle n \rangle} = \mathcal{L}_{\lambda,i \phi_{n-1}} \circ \dots \circ \mathcal{L}_{\lambda,i \phi_0} .$$ A direct computation yield
$$ \mathcal{L}_{\lambda,i \Phi}^{\langle n \rangle} h = \lambda^{-n} \sum_{\mathbf{a} \in \mathcal{A}^{n}} e^{i S_n \Phi \circ g_\mathbf{a}} h \circ g_\mathbf{a}, $$
where $S_n \Phi(\theta) = \sum_{k=0}^{n-1} \phi_k(\lambda^k \theta)$, so that  $S_n\Phi\circ g_{\mathbf{a}}(\theta) =\sum_{k=0}^{n-1} \phi_k(g_{a_{n-k}\dots a_{1}}(\theta))=\sum_{k=1}^{n} \phi_{n-k} (g_{a_{k}\dots a_{1}} \theta)$. We used here that $\lambda g_a(\theta) = \theta \mod 1$, implying $\lambda^k g_{a_n\dots a_1}\theta\equiv g_{a_{n-k}\dots a_1}\theta$. The reversed indexing here reflects the order of composition of the transfer operators. \\

%\textcolor{magenta}{- There is a slightly annoying mismatch between the $\phi$ indexing which is $\{0,\dots,n-1\}$ and the $a$ indexing which is $\{1,\dots,n\}$.}%\\ - More importantly, computing $L^{<3>}h$ gives $\lambda^{-3}\sum_{a_1,a_2,a_3\in \mathcal{A}}e^{i\phi_0(g_{a_3}g_{a_2}g_{a_1}\theta)}e^{i\phi_1(g_{a_2}g_{a_1}\theta)}e^{i\phi_2g_{a_1}\theta}h(g_{a_3}g_{a_2}g_{a_1}\theta)$. This is equal to the last given expression $\lambda^{-3}\sum_{\mathbf{a}\in \mathcal{A}^3} e^{i\sum_{k=1}^3\phi_{3-k}(g_{a_k\dots a_1}\theta)}h\circ g_{\mathbf{a}}(\theta)$. In contrast, $S_n\Phi\circ g_{\mathbf{a}}(\theta) = \sum_{k=0}^2 \phi_k(\lambda^kg_{a_3a_2a_1}\theta) = \phi_0(g_{a_3a_2a_1}\theta)+\phi_1(\lambda g_{a_3a_2a_1}\theta)+\phi_2(\lambda^2g_{a_3a_2a_1}\theta)$. I don't see how the two phases are equal. For example $\lambda^2g_{a_3}g_{a_2}g_{a_1}\theta = \lambda^2(x_{a_3a_2a_1}+\theta\lambda^{-3})$ while $g_{a_1}\theta = x_{a_1}+\theta\lambda^{-1}$. }

Before heading to the proof of dispersion, let us state the \say{Lasota-Yorke} estimates, which states that most of the oscillations in $\mathcal{L}_{\lambda,i\Phi}^{\langle n \rangle}(h)$ comes from the phases.

\begin{lemma}[Lasota-Yorke]
Let $\Phi$ be a sequence of phases satisfying $\|\Phi\|_{C^1} := \sup_n \|\phi_n\|_{C^1} < \infty$. Let $h \in C^{1}([0,1],\mathbb{C})$. Then $ |\mathcal{L}_{\lambda,i\Phi}^{\langle n \rangle}(h)| \leq \|h\|_\infty $, and $$ |(\mathcal{L}_{\lambda,i\Phi}^{\langle n \rangle} h)'| \leq \frac{2}{\lambda} \|\Phi\|_{C^1} \|h\|_\infty + \lambda^{-n} \|h'\|_\infty. $$
\end{lemma}

\begin{proof}
First of all:
$$ |\mathcal{L}_{\lambda,i\Phi}^{\langle n \rangle}(h)| = \Big|\lambda^{-n} \sum_{\mathbf{a} \in \mathcal{A}^{n}} e^{i S_n \Phi \circ g_\mathbf{a}} h \circ g_\mathbf{a} \Big| \leq \lambda^{-n} \sum_{\mathbf{a} \in \mathcal{A}^{n}} |h\circ g_\mathbf{a}| \leq \|h\|_\infty. $$
The second inequality comes from:
$$ (\mathcal{L}_{\lambda,i\Phi}^{\langle n \rangle} h)'(\theta) = \lambda^{-n} \sum_{\mathbf{a} \in \mathcal{A}^{n}} i (S_n \Phi \circ g_\mathbf{a})' e^{i S_n \Phi \circ g_\mathbf{a}} h \circ g_\mathbf{a} + \lambda^{-n} \sum_{\mathbf{a} \in \mathcal{A}^{n}} e^{i S_n \Phi \circ g_\mathbf{a}} \lambda^{-n} h' \circ g_\mathbf{a}  $$
$$ \lesssim \sup_{\mathbf{a} \in \mathcal{A}^n} |(S_n \Phi \circ g_\mathbf{a})'| \cdot \|h\|_\infty + \lambda^{-n} \|h'\|_\infty. $$
This gives the desired estimate, since $\lambda g_a(\theta) = \theta \mod 1$, and so
$(S_n \Phi \circ g_\mathbf{a})' = \Big(\sum_{k=1}^{n} \phi_{n-k} \circ g_{a_{k} \dots a_1}\Big)' \leq \sum_{k=1}^{n} |\phi_{n-k}'|_\infty \lambda^{-k} \leq \|\Phi\|_{C^1} \frac{1}{\lambda-1} \leq \frac{2}{\lambda} \|\Phi\|_{C^1}$. 
\end{proof}

\subsection{Dispersion and Dolgopyat}

The dispersion bound will be reduced to what we call \say{Dolgopyat's estimates}, which is a contraction bound on some twisted transfer operator. Before writing the formal statement, let us see why this is natural and what are the natural phases appearing.  \\

Let $m \in \mathbb{Z}$ and consider the quantity $ \int_0^1 e^{i t W(\theta)} e^{2 i \pi m \theta} d\theta$. We have, by the autosimilarity relation $W(\theta) = \varphi(\theta)+\mu W(\lambda \theta)$:
$$ \int_0^1 e^{i t W(\theta)} e^{2 i \pi m \theta} d\theta = \int_0^1 e^{i t \big(\varphi(\theta) + \mu W(\lambda \theta)\big) } e^{2 i \pi m \theta} d\theta . $$
Then, write for some $\epsilon \in \{-1,1\}$, $m=\epsilon |m| $, and do the euclidean division of $|m|$ by $\lambda$, writing $|m|=|m_0|+\lambda |m^{(1)}|$ with $|m_0| \leq \lambda$. Then set $m_0 := \epsilon |m_0|$ and $m^{(1)} := \epsilon |m^{(1)}|$, so that $m=m_0+\lambda m^{(1)}$ with $|m_0|\leq \lambda$. We then have: 
$$ \int_0^1 e^{i t W(\theta)} e^{2 i \pi m \theta} d\theta = \int_0^1 e^{i \big( t \varphi(\theta) + 2 \pi m_0 \theta \big) } e^{i t \mu W(\lambda \theta)} e^{2 i \pi m^{(1)} \lambda \theta} d\theta   $$
$$ = \int_0^1 \mathcal{L}_\lambda \Big(e^{i ( t \varphi + 2 \pi m_0 \cdot ) } \Big)(\theta) e^{i t \mu W(\theta)} e^{2 i \pi m^{(1)} \theta} d\theta  $$
$$ = \int_0^1 \mathcal{L}_{\lambda,i \phi_{t,m_0}}(1)(\theta) e^{i t \mu W(\theta)} e^{2 i \pi m^{(1)} \theta} d\theta  $$
where $\phi_{t,m}(\theta) := t\varphi(\theta)+2 \pi m \theta$. Iterating the argument $n(t)$ times, for $\mu^{n} t \simeq 1$, yields
\begin{equation} \label{e:tranappear}
\int_0^1 e^{i t W(\theta)} e^{2 i \pi m \theta} d\theta = \int_0^1 \mathcal{L}_{\lambda,\Phi(t,m)}^{\langle n \rangle}(1) e^{i t \mu^n W(\theta)} e^{2 i \pi m^{(n)} \theta} d\theta,
\end{equation}
for the sequence of phases $\Phi(t,m)_k(\theta) := \phi_{t \mu^k, m_k}(\theta) = t \mu^k \varphi(\theta) + 2 \pi m_k \theta$, where $m=\sum_{j} m_j \lambda^j$ is the decomposition of $m$ in base $\lambda$. Notice in particular that, since $0 \leq |m_k| \leq \lambda-1$, we have $\|\Phi(t,m)\|_{C^1} \leq t \|\varphi\|_{C^1} + 2 \pi \lambda \leq t \cdot 4\pi \lambda \|\varphi\|_{C^1} $ for $t \geq \|\varphi\|_{C^1}^{-1}$.  \\

Cauchy-Schwarz then yields the key bound:
$$ \Big| \int_0^1 e^{i t W(\theta)} e^{2 i \pi m \theta} d\theta \Big| \leq \sqrt{\int_0^1 |\mathcal{L}_{\lambda,\Phi(t,m)}^{\langle n \rangle}(1)|^2 d\theta}. $$

Dispersion is then reduced to establishing uniform contraction for this family of twisted transfer operators. This is the so-called Dolgopyat's estimates.

\begin{theorem}[$L^2$-Dolgopyat estimates]
Fix $\varphi$, $\mu$, $\lambda$, and suppose that $W$ is not analytic. There exists $\varepsilon_0(\mu,\lambda,\varphi)>0$ such that the following holds. Denoting $n(t) := \lfloor  \frac{\ln(t)}{2|\ln(\mu)|} \rfloor$, for any $t\geq 1$ and $m \in \mathbb{Z}$:
$$ \forall h \in C^1, \ \Big\| \mathcal{L}_{\lambda,\Phi(t,m)}^{\langle n \rangle}(h) \Big\|_{L^2([0,1])} \lesssim \frac{\|h\|_{C^1}}{ t^{\varepsilon_0}} .$$
\end{theorem}

Let us check that $L^2$-Dolgopyat's estimates imply polynomial dispersion.

\begin{proof}[Proof (of dispersion assuming $L^2$-Dolgopyat).]
Let $\mu \in (0,1)$, $\lambda \in \mathbb{N} \setminus \{0,1\}$, and let $\varphi \in C^\omega(\mathbb{S}^1,\mathbb{R})$ be a non-constant analytic function. If $W$ is analytic, then $W'$ is itself analytic, $1$-periodic, and non-constant. It follows that it cannot concentrate, and our Van der Corput Lemma~\ref{thm:c1+alphavandercorput} applies, yielding
$$ \sup_{m \in \mathbb{Z}} \Big| \int_0^1 e^{ i t W(\theta)} e^{2 i \pi m \theta} d\theta \Big| \leq \frac{C}{t^{1/k}}, $$
where $k$ is the maximal order of vanishing of $W'$.

We can then suppose that $W$ is not analytic. In this case, by our previous computations, we can write
$$ \Big| \int_0^1 e^{ i t W(\theta)} e^{2 i \pi m \theta} d\theta \Big| \leq \Big(\int_0^1 |\mathcal{L}_{\lambda,i\Phi(t,m)}^{\langle n \rangle}(1)|^2 d\theta\Big)^{1/2}. $$
The $L^2$-Dolgopyat estimates then yield
$$ \Big(\int_0^1 |\mathcal{L}_{\lambda,i\Phi(t,m)}^{\langle n \rangle}(1)|^2 d\theta\Big)^{1/2} \leq \frac{C}{t^{\varepsilon_0}} $$
for some $\varepsilon_0>0$ depending on $\mu,\lambda,\varphi$, concluding the proof.
\end{proof}

Fast dispersion when $\lambda$ is large will need more work and is postponed to a bit later. Let us turn into the proof of the $L^2$ bound. The idea is to open up the modulus square in the integral, to find (writing $\Phi := \Phi(t,m)$):
$$ \int_0^1 |\mathcal{L}_{\lambda,i\Phi}^{\langle n \rangle}(h)(\theta)|^2 d\theta = \int_0^1 \Big| \lambda^{-n} \sum_{\mathbf{a}} e^{i S_n \Phi \circ g_{\mathbf{a}}(\theta)} h(g_\mathbf{a}(\theta)) \Big|^2 d\theta  $$ $$ = \lambda^{-2n} \sum_{\mathbf{a},\mathbf{b} \in \mathcal{A}^n(\lambda)} \int_0^1 e^{i \big( S_n \Phi \circ g_\mathbf{a}(\theta) - S_n \Phi \circ g_\mathbf{b}(\theta) \big)} h(g_\mathbf{a}(\theta)) \overline{h(g_\mathbf{b}(\theta))} d\theta. $$
The idea is then to apply Van Der Corput Lemma to the phases $S_n\Phi \circ g_\mathbf{a} - S_n\Phi \circ g_\mathbf{b}$. For reasons of technical simplicity that will be revealed later, we will study their $k$-th derivative for $k$ large enough. 
We begin by a separation result: there exists at least two words $\mathbf{a},\mathbf{b}$ such that $(S_n\Phi \circ g_\mathbf{a} - S_n\Phi \circ g_\mathbf{b})^{(k)}$ doesn't vanish. Recall here that $\Phi(t,m)_j(\theta)=t\mu^j \varphi(\theta)+2\pi m_j \theta$, so that $\Phi(t,m)_j^{(k)}(\theta) = t \mu^{j} \varphi^{(k)}(\theta)$, and $(S_n \Phi \circ g_\mathbf{a})^{(k)} = t \mu^n \sum_{j=1}^{n} { \varphi^{(k)}(g_{a_j \dots a_1}(\theta))}{(\mu \lambda^k)^{-j}}$.

\begin{lemma}[Separation]\label{lem:sepk}
Suppose that $W$ is not analytic. Let $k \geq 2$ be such that $\mu \lambda^k \geq 2$. There exists $c_0>0$ and infinitely many $N \geq 1$, and $\mathbf{a},\mathbf{b} \in \mathcal{A}^N$ such that $$ \forall \theta \in [0,1], \  \Big|\sum_{j=1}^{N} \frac{ \varphi^{(k)}(g_{a_j \dots a_1}(\theta))}{(\mu \lambda^k)^j} - \sum_{j=1}^{N} \frac{ \varphi^{(k)}(g_{b_j \dots b_1}(\theta))}{(\mu \lambda^k)^j}\Big| \geq c_0. $$
where $\varphi^{(k)}$ the $k$-th derivative of $\varphi$.
\end{lemma}

\begin{proof}
Fix $k \geq 2$ large enough. First of all, notice that the statements with $\forall \theta$ and with $\exists \theta$ are equivalent. Indeed, suppose that there exist $c_0>0$, $N \geq 1$, $\mathbf{a},\mathbf{b} \in \mathcal{A}^N$ and $\theta_0 \in [0,1]$ such that
$$ \Big|\sum_{j=1}^{N} \frac{\varphi^{(k)}(g_{a_j \dots a_1}(\theta_0))}{(\mu \lambda^k)^j} - \sum_{j=1}^{N} \frac{\varphi^{(k)}(g_{b_j \dots b_1}(\theta_0))}{(\mu \lambda^k)^j}\Big| \geq c_0. $$
Then, since the expression is continuous in $\theta$, there exists a small ball $B(\theta_0,\varepsilon)$ such that the same bound holds for all $\theta \in B(\theta_0,\varepsilon)$, replacing $c_0$ by $c_0/2$. Since the derivatives of the above expressions are uniformly bounded in $N$, $\mathbf a$ and $\mathbf b$ (as $\mu\lambda^k>1$), the radius $\varepsilon$ may be chosen depending only on $c_0,\mu,\lambda,\varphi$. Hence the length $N_0$ of the word $\mathbf c$ may also be chosen uniformly. Now choose some word $\mathbf{c} \in \mathcal{A}^{N_0}$ such that $g_\mathbf{c}([0,1]) \subset B(\theta_0,\varepsilon)$. Using the decomposition according to the word $\mathbf c$, for every $\theta\in[0,1]$ we have
\begin{align*}
&\left|\sum_{j=1}^{N+N_0}\frac{\varphi^{(k)}(g_{(ac)_j\dots(ac)_1}(\theta))}{(\mu\lambda^k)^j} - \sum_{j=1}^{N+N_0}\frac{\varphi^{(k)}(g_{(bc)_j\dots(bc)_1}(\theta))}{(\mu\lambda^k)^j}\right| \\
&\qquad = \frac{1}{(\mu\lambda^k)^{N_0}}\left|\sum_{j=1}^{N}\frac{\varphi^{(k)}(g_{a_j\dots a_1}(g_{\mathbf c}(\theta)))}{(\mu\lambda^k)^j} - \sum_{j=1}^{N}\frac{\varphi^{(k)}(g_{b_j\dots b_1}(g_{\mathbf c}(\theta)))}{(\mu\lambda^k)^j}\right| \geq \frac{c_0/2}{(\mu\lambda^k)^{N_0}}.
\end{align*}
Therefore, since $g_{\mathbf c}(\theta)\in B(\theta_0,\varepsilon)$, the absolute value of the left-hand side is at least $c_0/(2(\mu\lambda^k)^{N_0})$ for every $\theta\in[0,1]$. This is the statement of Lemma~11.5 for the words $\mathbf{ac}$ and $\mathbf{bc}$.

Let us then prove this statement with $\exists$ replacing $\forall$, assuming that $W$ is not analytic. We will prove the converse: let us assume that
$$ \forall \varepsilon>0, \forall N \text{ large enough}, \forall \mathbf{a},\mathbf{b} \in \mathcal{A}^N, \forall \theta \in [0,1],\ \Big|\sum_{j=1}^{N} \frac{ \varphi^{(k)}(g_{a_j \dots a_1}(\theta))}{(\mu \lambda^k)^j} - \sum_{j=1}^{N} \frac{ \varphi^{(k)}(g_{b_j \dots b_1}(\theta))}{(\mu \lambda^k)^j}\Big| \leq \varepsilon, $$
and let us show that this implies that $W$ is analytic. Notice first that, being given an infinite word $\mathbf{a} \in \mathcal{A}^\mathbb{N}$ written $\dots a_n \dots a_1$, the series 
\begin{equation}\label{e:Xadef}
X_\mathbf{a}(\theta) := \sum_{j=1}^\infty \frac{\varphi^{(k)}(g_{a_j \dots a_1}(\theta))}{(\mu \lambda^k)^j}
\end{equation}
converges and is analytic in $\theta \in [0,1]$. Indeed, since $\varphi$ is real analytic and periodic, $\varphi^{(k)}$ extends holomorphically to a complex neighbourhood of $\mathbb R$. The inverse branch $g_{a_j\dots a_1}$ contracts imaginary parts by $\lambda^{-j}$, so all the functions $\varphi^{(k)}\circ g_{a_j\dots a_1}$ are holomorphic on a common complex neighbourhood of $[0,1]$ and uniformly bounded there. Since $\mu\lambda^k>1$, the coefficients $(\mu\lambda^k)^{-j}$ are summable, and hence the series converges uniformly on this neighbourhood. Its sum $X_{\mathbf a}$ is therefore holomorphic there, and in particular real analytic on $[0,1]$.

The hypothesis that we made precisely translates to
$$ \forall \mathbf{a},\mathbf{b} \in \mathcal{A}^\mathbb{N}, \forall \theta, \ X_\mathbf{a}(\theta) = X_{\mathbf{b}}(\theta). $$
This allows us to prove a cohomology condition on $\varphi^{(k)}$. Indeed, if we define on the circle the map $\psi_k$ by $\psi_k(\theta) := X_\mathbf{a}(\theta)$, then, for every $b\in\mathcal A$ and every $\theta\in g_b([0,1])$, the definition \eqref{e:Xadef} gives
$$ \mu\lambda^k X_{\mathbf a b}(\lambda\theta-b)=\varphi^{(k)}(\theta)+X_\mathbf a(\theta). $$
Since $X_{\mathbf a b}=X_\mathbf a=\psi_k$ and $\psi_k$ is $1$-periodic, this means that
$$ \varphi^{(k)}(\theta) = \mu \lambda^k \psi_k(\lambda \theta) - \psi_k(\theta). $$
Notice that $\psi_k$ is (by definition) $1$-periodic and analytic on each interval $g_b((0,1))$. The cohomology condition
$$ \mu\lambda^k\psi_k(\lambda\theta)=\varphi^{(k)}(\theta)+\psi_k(\theta) $$
then propagates this analyticity across the finitely many junction points. Hence $\psi_k$ extends analytically across these points and is therefore analytic on $\mathbb S^1$.

Let us show that this cohomology condition can be integrated to $\varphi^{(k-1)}$. Since $\psi_k$ is analytic and one periodic, it admits an analytic 1-periodic antiderivative if and only if $\int_0^1  \psi_k(\theta)d\theta = 0$. But notice that
$$ 0 = \int_0^1 \varphi^{(k)}(\theta)d\theta = (\mu \lambda^{k} -1) \int_0^1\psi_k(\theta) d\theta. $$
Hence $\int_0^1\psi_k(\theta) d\theta = 0$ indeed and our cohomology condition can be integrated into
$$ \varphi^{(k-1)}(\theta)
= \mu \lambda^{k-1} \psi_{k-1}(\lambda\theta)
- \psi_{k-1}(\theta) $$
for some analytic and $1$-periodic $\psi_{k-1}$. We will then iterate this argument until we reach a cohomology condition on $\varphi$. The only possible difficulty would be that there exists some $j$ such that $\mu \lambda^j=1$. In this case, to integrate the relation $$\varphi^{(j)}(\theta) = \psi_j(\lambda \theta) - \psi_j(\theta),$$ just notice that the same relation holds replacing $\psi_j$ with $\psi_j - \int_0^1 \psi_j$ directly, and then we can still integrate. This iterated argument allows us to get the following cohomology condition: there exists some $1$-periodic and analytic map $\psi$ such that
$$ \varphi(\theta) = \mu \psi(\lambda \theta) - \psi(\theta). $$
Indeed, at each integration the coefficient $\mu\lambda^j$ becomes $\mu\lambda^{j-1}$, so after $k$ integrations the coefficient is $\mu$. This implies that $W$ is analytic, since then the series defining $W$ is telescopic and yields $W=-\psi$. The proof is done.
\end{proof}

\begin{lemma}[Tree Lemma]
Suppose that $W$ is not analytic. Fix $k \geq 2$ such that $\mu \lambda^k \geq 2$. Then, for every $\theta \in \mathbb{S}^1$, for all $n \geq 1$, and any interval $I$, we have
$$
\lambda^{-n} \#\Big\{ \mathbf{a} \in \mathcal{A}^n, \frac{1}{t \mu^n}(S_n \Phi(t,m) \circ g_\mathbf{a})^{(k)}(\theta) \in I \Big\} \leq C\Big( |I|^\gamma + \rho^n \Big),
$$
where $\gamma \in (0,1)$ is such that $(1-\lambda^{-N}) (\mu \lambda^k)^{N\gamma} = 1$, $\rho \in (0,1)$ is such that $(1-\lambda^{-N})\rho^{-N}=1$, $C := (c_0/8)^{-1} \rho^{-N}$, and where $N,c_0$ are given by the previous lemma, with $N$ large enough so that $2^{-N} \|\varphi^{(k)}\|_\infty \leq c_0/16$.
\end{lemma}

The idea of the \say{tree lemma} is that the map $\mathbf{a} \in \{0,1\}^{n} \mapsto(S_n \Phi(t,m) \circ g_\mathbf{a})^{(k)}(\theta)$ looks like the coordinate map of points in a Cantor set. The previous bound is then a sort of dimension estimate on this Cantor set. The natural alphabetical \say{tree structure} of $\{0,1\}^n$ is preserved in some sense and allows to control the non-concentration of these quantities. Taking $k$ large enough ensures that different scales clearly appear in the sum $\sum_j \varphi^{(k)}(\dots) (\lambda^k \mu)^{-j}$.

\begin{proof}
First of all, recall that
$$
\frac{1}{t \mu^n}(S_n \Phi(t,m) \circ g_\mathbf{a})^{(k)}(\theta)
= \sum_{j=1}^{n} \frac{\varphi^{(k)}(g_{a_j \dots a_1}(\theta))}{(\mu \lambda^k)^j}
=: \ell_\mathbf{a}(\theta).
$$
The desired bound on the cardinality is proved by strong induction on $n$. If $n \leq N$, the bound is trivial since
$$
\lambda^{-n} \#\{ \mathbf{a} \in \mathcal{A}^n, \ell_\mathbf{a}(\theta) \in I \}
\leq 1 \leq C\Big( |I|^\gamma + \rho^n \Big)
$$
as $C \rho^{N} \geq 1$. We can then assume that $n \geq N+1$ and suppose that the result holds for smaller $n$. If $|I| \geq c_0/4$, the bound is also trivial, since
$$
\lambda^{-n} \#\{ \mathbf{a} \in \mathcal{A}^n, \ell_\mathbf{a}(\theta) \in I \}
\leq 1 \leq C\Big( |I|^\gamma + \rho^n \Big),
$$
as $C (c_0/4)^{\gamma} \geq 1$. Let us then assume that $|I| \leq c_0/4$. We write
$$
\lambda^{-n} \#\{ \mathbf{a} \in \mathcal{A}^n, \ell_\mathbf{a}(\theta) \in I \}
= \lambda^{-N} \sum_{\widehat{\mathbf{a}} \in \mathcal{A}^N }
\lambda^{-(n-N)} \# \{\tilde{\mathbf{a}} \in \mathcal{A}^{n-N},
\ell_{\tilde{\mathbf{a}}\widehat{\mathbf{a}}}(\theta) \in I \}.
$$

Our separation condition ensures that, for the fixed $\theta$, there exists $\widehat{\mathbf a}_0\in\mathcal A^N$ such that
$$
d(\ell_{\widehat{\mathbf a}_0}(\theta),I)\geq c_0/4.
$$
Indeed, if $\mathbf a,\mathbf b$ are given by Lemma~\ref{lem:sepk}, then $|\ell_{\mathbf a}(\theta)-\ell_{\mathbf b}(\theta)|\geq c_0$, so at least one of them has distance at least $c_0/4$ from $I$. Furthermore, a direct computation yields
\begin{equation}\label{e:ell}
\forall \tilde{\mathbf{a}} \in \mathcal{A}^{n-N}, \qquad
\ell_{\tilde{\mathbf{a}}\widehat{\mathbf{a}}}(\theta)
=
\ell_{\widehat{\mathbf{a}}}(\theta)
+
\frac{\ell_{\tilde{\mathbf{a}}}(g_{\widehat{\mathbf{a}}}(\theta))}
{(\mu \lambda^k)^N}.
\end{equation}
Since $\mu\lambda^k\geq2$, uniformly in $\tilde{\mathbf a}$ and $\theta$,
$$
\left|
\frac{\ell_{\tilde{\mathbf{a}}}(g_{\widehat{\mathbf{a}}_0}(\theta))}
{(\mu \lambda^k)^N}
\right|
\leq \frac{\|\varphi^{(k)}\|_\infty}{2^N}
\leq \frac{c_0}{16}.
$$
Thus, by \eqref{e:ell},
$$
d(\ell_{\tilde{\mathbf a}\widehat{\mathbf a}_0}(\theta),I)
\geq \frac{c_0}{4}-\frac{c_0}{16}>0
$$
for every $\tilde{\mathbf a}\in\mathcal A^{n-N}$. Hence no word ending in $\widehat{\mathbf a}_0$ contributes to the cardinality above. Since the constants are independent of $\theta$, this argument is uniform in $\theta$.

This allows us to write, using the induction hypothesis,
\begin{align*}
&\lambda^{-n} \#\{ \mathbf{a} \in \mathcal{A}^n : \ell_\mathbf{a}(\theta) \in I \} \\
&\quad = \lambda^{-N}
\sum_{\widehat{\mathbf{a}} \in \mathcal{A}^N \setminus \{\widehat{\mathbf{a}}_0\}}
\lambda^{-(n-N)}
\#\{ \tilde{\mathbf{a}} \in \mathcal{A}^{n-N} :
\ell_{\tilde{\mathbf{a}}\widehat{\mathbf{a}}}(\theta) \in I \} \\
&\quad = \lambda^{-N}
\sum_{\widehat{\mathbf{a}} \in \mathcal{A}^N \setminus \{\widehat{\mathbf{a}}_0\}}
\lambda^{-(n-N)}
\#\Big\{ \tilde{\mathbf{a}} \in \mathcal{A}^{n-N} :
\ell_{\tilde{\mathbf{a}}}(g_{\widehat{\mathbf{a}}}(\theta))
\in (\mu\lambda^k)^N
\big(I-\ell_{\widehat{\mathbf{a}}}(\theta)\big) \Big\} \\
&\quad \leq \lambda^{-N}
\sum_{\widehat{\mathbf{a}} \in \mathcal{A}^N \setminus \{\widehat{\mathbf{a}}_0\}}
C\Big( |(\mu\lambda^k)^N I|^\gamma + \rho^{n-N} \Big) \\
&\quad = C\big(|I|^\gamma+\rho^n\big),
\end{align*}
since $(1-\lambda^{-N})(\mu\lambda^k)^{N\gamma}
=(1-\lambda^{-N})\rho^{-N}=1$. The induction is done.
\end{proof}

We are ready to prove our $L^2$-Dolgopyat estimates. This is essentially a use of Van Der Corput's lemma.


\begin{proof}[Proof (of $L^2$-Dolgopyat estimates)]
We choose $n$ such that $\mu^n t \sim t^{1/2}$. We start by writing, denoting $\Phi=\Phi(t,m)$ for clarity,
\begin{align*}
\Big\| \mathcal{L}_{\lambda,\Phi(t,m)}^{\langle n \rangle}(h) \Big\|_{L^2([0,1])}^2 &= \int_0^1 |\mathcal{L}_{\lambda,i\Phi(t,m)}^{\langle n \rangle}(h)(\theta)|^2 d\theta \\
&= \int_0^1 \Big| \lambda^{-n} \sum_{\mathbf{a}} e^{i S_n \Phi \circ g_{\mathbf{a}}(\theta)} h(g_\mathbf{a}(\theta)) \Big|^2 d\theta \\
&= \lambda^{-2n} \sum_{\mathbf{a},\mathbf{b} \in \mathcal{A}^n(\lambda)} \int_0^1 e^{i \big( S_n \Phi \circ g_\mathbf{a}(\theta) - S_n \Phi \circ g_\mathbf{b}(\theta) \big)} h(g_\mathbf{a}(\theta)) \overline{h(g_\mathbf{b}(\theta))} d\theta.
\end{align*}

Let us then iterate a little more our transfer operator. Recall that $\int h d\theta=\int \mathcal{L}_\lambda^{\mathfrak n}(h)d\theta$ for any $\mathfrak n$. We choose $\mathfrak n$ such that $\lambda^{2\mathfrak n}\simeq t^{1/(4k)}$ and write
$$
\Big\| \mathcal{L}_{\lambda,\Phi(t,m)}^{\langle n \rangle}(h) \Big\|_{L^2([0,1])}^2 = \frac{1}{\lambda^{2n+2\mathfrak n}} \sum_{\substack{\mathbf a,\mathbf b\in\mathcal A(\lambda)^n\\ \mathbf c\in\mathcal A(\lambda)^{2\mathfrak n}}} \int_0^1 e^{i \big( S_n \Phi \circ g_\mathbf{a}(g_\mathbf{c}(\theta)) - S_n \Phi \circ g_\mathbf{b}(g_\mathbf{c}(\theta)) \big)} h(g_\mathbf{ac}(\theta)) \overline{h(g_\mathbf{bc}(\theta))} d\theta.
$$
Since $h$ is $C^1$ and $g_{\mathbf{ac}}'=\lambda^{-n-2\mathfrak n}\leq\lambda^{-2\mathfrak n}\lesssim t^{-1/(4k)}$, we have $h(g_{\mathbf{ac}}(\theta))=h(x_{\mathbf{ac}})+\mathcal{O}(\|h\|_{C^1}t^{-1/(4k)})$, and therefore
$$
\Big\| \mathcal{L}_{\lambda,\Phi(t,m)}^{\langle n \rangle}(h) \Big\|_{L^2([0,1])}^2 \leq \frac{\|h\|_\infty^2}{\lambda^{2n+2\mathfrak n}} \sum_{\substack{\mathbf a,\mathbf b\in\mathcal A(\lambda)^n\\ \mathbf c\in\mathcal A(\lambda)^{2\mathfrak n}}} \left| \int_0^1 e^{i \big( S_n \Phi \circ g_\mathbf{a}(g_\mathbf{c}(\theta)) - S_n \Phi \circ g_\mathbf{b}(g_\mathbf{c}(\theta)) \big)} d\theta \right| + \frac{\|h\|_{C^1}^2}{t^{1/(4k)}}.
$$

Let us concentrate on the integral $I_{\mathbf{a},\mathbf{b},\mathbf{c}}(t):=\int_0^1 e^{i \big( S_n \Phi \circ g_\mathbf{a}(g_\mathbf{c}(\theta)) - S_n \Phi \circ g_\mathbf{b}(g_\mathbf{c}(\theta)) \big)}d\theta$. We write
$$
I_{\mathbf{a},\mathbf{b},\mathbf{c}}(t)=\int_0^1 e^{iT\psi_{\mathbf{a},\mathbf{b},\mathbf{c}}(\theta)}d\theta,
$$
where $T=t\mu^n\lambda^{-(2k+1)\mathfrak n}\simeq t^{(2k-1)/(8k)}$ and
$$
\psi_{\mathbf{a},\mathbf{b},\mathbf{c}}(\theta)=t^{-1}\mu^{-n}\lambda^{(2k+1)\mathfrak n}\big(S_n\Phi\circ g_{\mathbf a}(g_{\mathbf c}(\theta))-S_n\Phi\circ g_{\mathbf b}(g_{\mathbf c}(\theta))\big).
$$
The classical Van Der Corput lemma gives $|I_{\mathbf{a},\mathbf{b},\mathbf{c}}(t)|\leq c_kT^{-1/k}\lesssim t^{-(2k-1)/(8k^2)}$ whenever $|\psi_{\mathbf{a},\mathbf{b},\mathbf{c}}^{(k)}(\theta)|\geq1$ on $[0,1]$. Hence
\begin{align*}
& \Big\| \mathcal{L}_{\lambda,\Phi(t,m)}^{\langle n \rangle}(h) \Big\|_{L^2([0,1])}^2 \\
& \lesssim \|h\|_{C^1}^2 \Big(t^{-(2k-1)/(8k^2)}+t^{-1/(4k)}+\lambda^{-2n-2\mathfrak n}\#\{(\mathbf a,\mathbf b,\mathbf c)\in\mathcal A^{2n+2\mathfrak n}:\exists\theta,\ |\psi_{\mathbf a,\mathbf b,\mathbf c}^{(k)}(\theta)|\leq1\}\Big).
\end{align*}

To conclude, we bound this cardinality. We first compute
\begin{align*}
\psi_{\mathbf{a},\mathbf{b},\mathbf{c}}^{(k)}(\theta) &= \lambda^{\mathfrak n}\big(\ell_{\mathbf a}(g_\mathbf c(\theta))-\ell_{\mathbf b}(g_\mathbf c(\theta))\big) \\
&= \lambda^{\mathfrak n}\big(\ell_{\mathbf a}(g_\mathbf c(0))-\ell_{\mathbf b}(g_\mathbf c(0))\big)+\mathcal O(\lambda^{-\mathfrak n}).
\end{align*}
It follows that, for $t$ large enough, using the tree lemma,
\begin{align*}
&\lambda^{-2n-2\mathfrak n}\#\{(\mathbf a,\mathbf b,\mathbf c)\in\mathcal A^{2n+2\mathfrak n}:\exists\theta,\ |\psi_{\mathbf a,\mathbf b,\mathbf c}^{(k)}(\theta)|\leq1\} \\
&\quad\leq\lambda^{-2n-2\mathfrak n}\#\{(\mathbf a,\mathbf b,\mathbf c)\in\mathcal A^{2n+2\mathfrak n}:\lambda^{\mathfrak n}|\ell_{\mathbf a}(g_\mathbf c(0))-\ell_{\mathbf b}(g_\mathbf c(0))|\leq2\} \\
&\quad\leq\lambda^{-n-2\mathfrak n}\sum_{(\mathbf b,\mathbf c)\in\mathcal A^{n+2\mathfrak n}}\lambda^{-n}\#\{\mathbf a\in\mathcal A^n:\ell_\mathbf a(g_\mathbf c(0))\in[\ell_\mathbf b(g_\mathbf c(0))-2\lambda^{-\mathfrak n},\ell_\mathbf b(g_\mathbf c(0))+2\lambda^{-\mathfrak n}]\} \\
&\quad\leq2C\big(\lambda^{-\gamma\mathfrak n}+\rho^n\big)\lesssim t^{-\varepsilon_0},
\end{align*}
where one may take
$$
\varepsilon_0=\min\left\{\frac{2k-1}{8k^2},\frac{1}{4k},\frac{\gamma}{8k},\frac{|\log\rho|}{2|\log\mu|}\right\}>0.
$$
Thus
$$
\Big\| \mathcal{L}_{\lambda,\Phi(t,m)}^{\langle n \rangle}(h) \Big\|_{L^2}^2
\lesssim t^{-\varepsilon_0}\|h\|_{C^1}^2.
$$
Taking square roots and replacing $\varepsilon_0$ by $\varepsilon_0/2$ proves the desired $L^2$-Dolgopyat estimate.
\end{proof}

\subsection{Fast dispersion for large $\lambda$}

We proved that dispersion always occurs for Weierstrass crystals, but the rate of dispersion is not very explicit in general. We will now turn to the case where $\mu$ and $\varphi$ are fixed, but $\lambda$ is allowed to be chosen large. We will show that in this case, Dolgopyat's estimates can be improved and the contraction of twisted transfer operators occurs faster. Let us state our enhanced Dolgopyat's estimate.


\begin{theorem}[Dolgopyat's estimates]\label{thm:trueDolgopyat}
Let $\varphi$ be analytic, $1$-periodic, and non-constant. Let $\mu \in (0,1)$. Then there exists $\rho \in (0,1)$ such that the following holds. For every $\beta\in(0,1]$ there exists $\lambda_*(\beta)\geq1$ such that, for every $\lambda\geq\lambda_*(\beta)$, there exists $t_*(\beta,\lambda)\geq1$ such that for all $t \geq t_*(\beta,\lambda)$, we have, uniformly in $m$:
$$ \forall h \in C^{1}, \ \Big\| \mathcal{L}_{\lambda,i\Phi(t,m)}^{\langle n_\beta(t) \rangle}(h) \Big\|_{1,{t \mu^{n_\beta(t)}}} \leq \frac{\|h\|_{1,t}}{t^\rho} $$
where $n_\beta(t) := \Big\lfloor \frac{\beta \ln(t)}{|\ln(\mu)|} \Big\rfloor $, so that $\mu^{n_\beta(t)} \simeq t^{-\beta}$, and where $\|h\|_{C^1_t} = \|h\|_{1,t} := \|h\|_\infty+\frac{\|h'\|_{\infty}}{t}$. The constant $\lambda_*(\beta)$ can be chosen of the form $\lambda_*(\beta) := \max(C_1(\varphi),\mu^{-C_2(\varphi)/\beta})$.
\end{theorem}

First of all, let us prove that this implies fast dispersion when $\lambda$ is large.

\begin{proof}[Proof of Theorem \ref{thm:main}(3) assuming Theorem \ref{thm:trueDolgopyat}]
Let $\lambda$ be a large integer, and let $N$ be the maximal integer such that $\lambda \geq \lambda_*(1/N)$. We have $N \gtrsim_{\varphi} \frac{\ln(\lambda)}{|\ln(\mu)|}$. Then, choose for all $0 \leq j \leq N-1$ the parameters $\beta_j := \frac{1}{N-j}$ (so that $(1-\frac{j}{N})(1-\beta_j) = 1-\frac{j+1}{N}$). Notice that $\lambda \geq \max_j \lambda_*(\beta_j)$, since $\lambda_*(\beta)=\max(C_1(\varphi),\mu^{-C_2(\varphi)/\beta})$ is decreasing in $\beta$ and $\beta_j\geq 1/N$. Suppose that $t \geq \max_j t_*(\beta_j,\lambda)^N$. Recall that \eqref{e:tranappear}, choosing $n := N \lfloor \frac{\ln t}{N|\ln(\mu)|} \rfloor $ (so that $t \simeq \mu^{-n}$), we have
$$ \int_{\mathbb{S}^1} e^{i t W(\theta)} e^{2 i \pi m \theta} d\theta = \int_{\mathbb{S}^1} \mathcal{L}_{\lambda,i \Phi(t,m)}^{\langle n \rangle}(1)(\theta) e^{i t \mu^n W(\theta)} e^{2 i \pi m^{(n)}\theta} d\theta. $$
We then cut the composition in $N$ equal parts, introducing the times $t_j \simeq t \mu^{jn/N} \simeq t^{(1-\frac{j}{N})}$ (satisfying $t_{j} \mu^{n_{\beta_j}(t_{j})}\simeq t_j^{1-\beta_j}=t_{j+1}$) and $t_j \geq t^{1-j/N} \geq t_*(\beta_j,\lambda)$:
$$ \Big| \int_{\mathbb{S}^1} e^{i t W(\theta)} e^{2 i \pi m \theta} d\theta  \Big| \leq \prod_{j=0}^{N-1} \sup_{\tilde m} \| \mathcal{L}_{\lambda,i\Phi(t_j,\tilde{m})}^{n_{\beta_j}(t_j)} \|_{\mathcal{L}({C}^{1}_{t_{j}} , {C^{1}_{t_{j+1}}})}  \leq \prod_{j=0}^{N-1} t_j^{-\rho} \leq t^{-\rho \sum_{j=0}^{N-1} (1-j/N)} \leq t^{-\rho N/2}. $$
This proves dispersion at speed $t^{-C(\varphi) \frac{\ln(\lambda)}{|\ln(\mu)|}}$.
\end{proof}
Let us then turn into the proof of the enhanced Dolgopyat's estimates. We first show that they can be reduced to their $L^2$-version.



\begin{lemma}
Suppose that there exists $\rho \in (0,1)$ that depends only on $\varphi$ such that, for all $\beta \in (0,1/2]$, there exists $\lambda_*(\beta)$ of the form $\max(C_1(\varphi),\mu^{-C_2(\varphi)/\beta})$ such that for all $\lambda \geq \lambda_*(\beta)$, and for all $t$ large enough depending on $\beta,\lambda$, we have $$ \forall h \in C^1, \  \int_{\mathbb{S}^1} |\mathcal{L}_{\lambda,i\Phi(t,m)}^{\langle n_\beta(t) \rangle} h|^2 d\theta \leq \frac{\|h\|_{1,t}^2}{t^{\rho}} .$$
Then Dolgopyat estimates holds.
\end{lemma}

\begin{proof}

Let $\beta \in (0,1]$. Let $\varepsilon := \rho/8$, and fix $\beta' := \beta(1-\varepsilon)$ and $\beta'':=\beta'/2$. For each $\beta''$, let $t_*(\beta'')$ be such that the $L^2$ estimate in the hypothesis holds for all $t\geq t_*(\beta'')$. Now let $\lambda \geq \max( \mu^{-\frac{100}{\beta \rho}} ,\lambda_*(\beta''))$ and $t \geq t_*(\beta'')$. Notice that, since $\varepsilon=\rho/8$ and $\beta'=\beta(1-\varepsilon)$, this choice of $\lambda$ is in particular large enough to satisfy the lower bounds $\lambda\geq\mu^{-3/(\varepsilon\beta)}$ and $\lambda\geq\mu^{-10/\beta'}$ used below. Moreover, since $\beta''$ is a fixed multiple of $\beta$, this lower bound on $\lambda$ is still of the form $\max(C_1(\varphi),\mu^{-C_2(\varphi)/\beta})$. \\
First of all, recall that (for $t \geq \|\varphi\|_{C^1}^{-1}$) $$ \|\Phi(t,m)\|_{C^1} \leq 4 \pi \lambda t \|\varphi\|_{C^1}. $$
In particular, Lasota-Yorke gives, for any $n$:
$$ \|(\mathcal{L}_{\lambda,i\Phi(t,m)}^{\langle n \rangle}h)'\|_\infty \leq 8 \pi t \| \varphi \|_{C^1} \|h\|_\infty + \lambda^{-n}\| h' \|_\infty$$
and  $\|\mathcal{L}_{\lambda,i\Phi(t,m)}^{\langle n \rangle}(h)\|_\infty \leq \|h\|_\infty$, so that the $\|\cdot \|_{1,t}$ norm can be bounded by
\begin{align*}
\|\mathcal{L}_{\lambda,i\Phi(t,m)}^{\langle n \rangle}(h)\|_{1,t \mu^{n}} &=  \|\mathcal{L}_{\lambda,i\Phi(t,m)}^{\langle n \rangle}(h)\|_{\infty}  + \frac{\|(\mathcal{L}_{\lambda,i\Phi(t,m)}^{\langle n \rangle} h)'\|_\infty}{t \mu^{n}}   \\
& \leq \|h\|_{\infty}  + t^{-1} \mu^{-n} \Big( 8 \pi t \| \varphi \|_{C^1} \|h\|_\infty + \lambda^{-n}\| h' \|_\infty \Big)  \\
& \leq \big(1+8\pi \|\varphi\|_{C^1} \mu^{-n} \big) \|h\|_\infty + \frac{\|h'\|_\infty}{(t \mu^n) \lambda^n} .
\end{align*}
We apply this by writing $\mathcal{L}_{\lambda,i\Phi(t,m)}^{\langle n_\beta(t) \rangle }(h) = \mathcal{L}_{\lambda,i\Phi(t',m')}^{\langle n_{\beta}(t)-n_{\beta'}(t) \rangle}\big( \mathcal{L}_{\lambda,i\Phi(t,m)}^{\langle n_{\beta'}(t) \rangle}(h) \big) $ where $\beta' := \beta(1-\varepsilon)$, $m'$ is some integer, and $t' := t \mu^{n_{\beta(1-\varepsilon)}(t)} \sim t^{1-\beta'}$. We recall here that  $n_\beta(t) := \Big\lfloor \frac{\beta \ln(t)}{|\ln(\mu)|} \Big\rfloor $, so that $\mu^{n_\beta(t)} \simeq t^{-\beta}$. Notice also that $\mu^{n_{\beta}(t)-n_{\beta'}(t)} \simeq t^{-\varepsilon \beta}$. Hence we find:
$$ \|\mathcal{L}_{\lambda,i\Phi(t,m)}^{\langle n_\beta(t) \rangle}(h)\|_{1,t \mu^{n_\beta(t)}} \lesssim t^{\varepsilon \beta} \| \mathcal{L}_{\lambda,i\Phi(t,m)}^{\langle n_{\beta'}(t) \rangle}(h) \|_\infty + \frac{\|(\mathcal{L}_{\lambda,i\Phi(t,m)}^{\langle n_{\beta'}(t) \rangle}h)' \|_\infty}{t^{1-\beta} \lambda^{n_\beta(t) - n_{\beta'}(t)}}. $$
Applying Lasota-Yorke another time on the second term yields:
\begin{align*}\|\mathcal{L}_{\lambda,i\Phi(t,m)}^{{\langle n_\beta(t) \rangle}}(h)\|_{1,t \mu^{n_\beta(t)}} &\lesssim t^{\varepsilon \beta} \| \mathcal{L}_{\lambda,i\Phi(t,m)}^{\langle n_{\beta'}(t) \rangle}(h) \|_\infty + \frac{8 \pi t \| \varphi \|_{C^1} \|h\|_\infty + \lambda^{-n_{\beta'}(t)}\| h' \|_\infty}{t^{1-\beta} \lambda^{n_\beta(t) - n_{\beta'}(t)}} \\
 &\leq t^{\varepsilon \beta} \| \mathcal{L}_{\lambda,i\Phi(t,m)}^{\langle n_{\beta'}(t) \rangle}(h) \|_\infty + \frac{8 \pi t^\beta \| \varphi \|_{C^1} \|h\|_\infty }{ \lambda^{n_\beta(t) - n_{\beta'}(t)}} + t^{\beta} \lambda^{- n_\beta(t)} \frac{ \| h' \|_\infty}{t} .
\end{align*}
Since $n_\beta(t) - n_{\beta'}(t) \simeq \varepsilon \beta \ln(t)/|\ln(\mu)| $, the fact that $\lambda \geq  \mu^{-\frac{3}{\varepsilon \beta} }$ ensures that $\lambda^{n_\beta(t) - n_{\beta'}(t)} \geq t^{2}$ when $t$ is large enough. It follows that, in that case, for $t$ large enough:
$$ \|\mathcal{L}_{\lambda,i\Phi(t,m)}^{\langle n_\beta(t) \rangle}(h)\|_{1,t \mu^{n_\beta(t)}} \lesssim t^{\varepsilon \beta} \| \mathcal{L}_{\lambda,i\Phi(t,m)}^{\langle n_{\beta'}(t) \rangle}(h) \|_\infty + \frac{\|h\|_{1,t}}{t}.  \quad \quad (*) $$
To reduce to the $L^2$ version of Dolgopyat's estimates, we use Perron-Frobenius-Ruelle in the following way. We have, by Cauchy-Schwarz, denoting $\beta'':=\beta'/2$ (for some time $t''$ and integer $m''$):
\begin{align*} |\mathcal{L}_{\lambda,i\Phi(t,m)}^{\langle n_{\beta'}(t) \rangle}(h)|^2 &= |\mathcal{L}_{\lambda,i\Phi(t'',m'')}^{\langle n_{\beta'}(t)-n_{\beta''}(t) \rangle}( \mathcal{L}_{\lambda,i\Phi(t,m)}^{\langle n_{\beta''}(t) \rangle}(h) )|^2 \\
&\leq \mathcal{L}_\lambda^{n_{\beta'}(t)-n_{\beta''}(t)}\Big( \big| \mathcal{L}_{\lambda,i\Phi(t,m)}^{\langle n_{\beta''}(t) \rangle}(h) \big|^2 \Big).
\end{align*}
Now, Perron-Frobenius-Ruelle and Lasota-Yorke yield 
$$ \mathcal{L}_\lambda^{n_{\beta'}(t)-n_{\beta''}(t)}\Big( \big| \mathcal{L}_{\lambda,i\Phi(t,m)}^{\langle n_{\beta''}(t) \rangle}(h) \big|^2 \Big) \lesssim \int_0^1 \big| \mathcal{L}_{\lambda,i\Phi(t,m)}^{\langle n_{\beta''}(t) \rangle}(h) \big|^2 d\theta + \Big\| \big| \mathcal{L}_{\lambda,i\Phi(t,m)}^{\langle n_{\beta''}(t) \rangle}(h) \big|^2 \Big\|_{C^1} \lambda^{-(n_{\beta'}(t)-n_{\beta''}(t))}  $$
$$ \leq \int_0^1 \big| \mathcal{L}_{\lambda,i\Phi(t,m)}^{\langle n_{\beta''}(t) \rangle}(h) \big|^2 d\theta + \|h\|_\infty\Big( \|h\|_\infty+ 2\Big\| \Big(\mathcal{L}_{\lambda,i\Phi(t,m)}^{\langle n_{\beta''}(t) \rangle}(h)\Big)'\Big\|_\infty \Big) \lambda^{-(n_{\beta'}(t)-n_{\beta''}(t))} $$
$$ \leq \int_0^1 \big| \mathcal{L}_{\lambda,i\Phi(t,m)}^{\langle n_{\beta''}(t) \rangle}(h) \big|^2 d\theta + \lambda^{-(n_{\beta'}(t)-n_{\beta''}(t))} \|h\|_\infty \Big( (1+16 \pi t\|\varphi\|_{C^1})\|h\|_\infty + 2\lambda^{-n_{\beta''}(t)}\| h' \|_\infty \Big) $$
$$ \lesssim \int_0^1 \big| \mathcal{L}_{\lambda,i\Phi(t,m)}^{\langle n_{\beta''}(t) \rangle}(h) \big|^2 d\theta + t \lambda^{-(n_{\beta'}(t)-n_{\beta''}(t))} \|h\|_{1,t}^2 $$
$$ \lesssim  \int_0^1 \big| \mathcal{L}_{\lambda,i\Phi(t,m)}^{\langle n_{\beta''}(t) \rangle}(h) \big|^2 d\theta + \frac{\|h\|_{1,t}^2}{t^2} $$
since $\lambda^{-(n_{\beta'}(t)-n_{\beta''}(t))} \sim t^{- \ln(\lambda) \beta'/|\ln(\mu)| } \leq t^{-3}$ for $t$ large enough, as $\lambda \geq \mu^{-10/\beta'}$. Plugging this inequality in $(*)$ yields
$$ \|\mathcal{L}_{\lambda,i\Phi(t,m)}^{\langle n_\beta(t) \rangle}(h)\|_{1,t \mu^{n_\beta(t)}} \lesssim t^{\varepsilon \beta} \Bigg( \Big(\int_0^1 \big| \mathcal{L}_{\lambda,i\Phi(t,m)}^{\langle n_{\beta''}(t) \rangle}(h) \big|^2 d\theta\Big)^{1/2} +  \frac{\|h\|_{1,t}}{t} \Bigg). $$
We conclude by using our $L^2$-Dolgopyat hypothesis: since $\lambda \geq \lambda_*(\beta'')$, we have, for $t$ large enough:
$$ \|\mathcal{L}_{\lambda,i\Phi(t,m)}^{\langle n_\beta(t) \rangle}(h)\|_{1,t \mu^{n_\beta(t)}} \lesssim t^{\varepsilon \beta} \Big( \frac{ \|h\|_{1,t}}{t^{\rho/2}}  +  \frac{\|h\|_{1,t}}{t} \Big) \leq \frac{\|h\|_{1,t}}{t^{\rho/4}}, $$
since $\varepsilon \beta \leq \rho/8$. The choice $\varepsilon=\rho/8$ is made precisely so that the loss $t^{\varepsilon\beta}$ can be absorbed by the gain $t^{-\rho/2}$.
\end{proof}



Now we turn to the verification of the $L^2$ bound. The strategy is adapted from earlier, but critically the separation lemma and the tree lemma needs to be adapted to a \say{large lambda} version, since the previous bound gets very poor as $\lambda$ grows large. This is the purpose of the next three lemmas, that culminate in a large lambda friendly version of the tree lemma. Our primary goal is to prove a non-concentration condition on the second derivative of $S_n \Phi(t,m) \circ g_\mathbf{a}$. We thus define, as earlier (taking $k=2)$: 
$$ \forall \mathbf{a} \in \mathcal{A}^*, \ \ell_\mathbf{a}(\theta) := \sum_{j=1}^{n} \frac{ \varphi^{(2)}(g_{a_j \dots a_1}(\theta))}{(\mu \lambda^2)^{j-1}} ,$$
so that 
$$\frac{\mu\lambda^2}{t\mu^n}(S_n \Phi(t,m) \circ g_\mathbf{a})'' = \ell_{\mathbf{a}}(\theta) = \varphi^{(2)}(g_{a_1}(\theta)) + \mathcal{O}(\lambda^{-1}) = \ell_{a_1}(\theta)+\mathcal{O}(\lambda^{-1}).$$ 
Recall that $\mathcal{A}^* = \bigcup_n \mathcal{A}^n.$
Since $\lambda$ is chosen large with respect to $\mu$ and $\varphi$, we can assume that $\mu \lambda \geq 2$, so that the previous series converges uniformly in $\mathbf{a},\theta$. We study the distribution of $(\ell_\mathbf{a}(\theta))_{\mathbf{a} \in \mathcal{A}^*}$.
\begin{lemma}[non-concentration of periodic analytic maps]
Let $f:\mathbb{R} \rightarrow \mathbb{R}$ be real analytic, non-constant, and $1$-periodic. Then there exists $\gamma_f \in (0,1)$ and $L_f\in(0,1)$ such that for any segment $I \subset \mathbb{S}^1$ of diameter $\leq L_f$: 
$$ |\{ \theta \in \mathbb{S}^1, \ f(\theta) \in I \}| \leq \frac{1}{10} {|I|^{\gamma_f}}. $$
\end{lemma}

\begin{proof}
For every $\theta_0 \in \mathbb{S}^1$, one can write on some small enough ball $B(\theta_0)$ centered at $\theta_0$: $f(\theta) = f(\theta_0)+F_{\theta_0}(\theta-\theta_0)^{k}$ where $F_{\theta_0}$ is a local diffeomorphism. The number $k$ is bounded from above by some $k_{max}$ since $f$ is analytic and not constant. By compactness, we can cover $\mathbb{S}^1$ into a finite number of these balls $B_j$. We then bound:
$$ |\{ \theta \in \mathbb{S}^1, \ f(\theta) \in I \}| = \sum_j |\{ \theta \in B_j, \ f(\theta) \in I \}|$$
$$ \leq \sum_j C_j \ |\{ x \in F_j^{-1}(B_j), \ x^{k_j} \in F_j^{-1}(I-f(\theta_j)) \}| \leq C|I|^{1/k_{max}}. $$
For some constant $C$ that only depends on $f$. Now, if $|I|\leq L_f := (10 C)^{-2k_{max}}$, then
$$ |\{ \theta \in \mathbb{S}^1, \ f(\theta) \in I \}|
\leq C |I|^{1/(2k_{max})} |I|^{1/(2k_{max})}
\leq \frac{1}{10}|I|^{1/(2k_{max})}, $$
which is the desired bound for $\gamma_f := 1/(2k_{max})$.
\end{proof}

\begin{lemma}[Separation Lemma]
Suppose that $\varphi$ is analytic, 1-periodic, not constant. There exists $\gamma_{sep}(\varphi) \in (0,1)$ such that, for $\lambda$ large enough $(\geq \max(4(\|\varphi\|_{C^3} L_{\varphi^{(2)}})^{2}, \mu^{-1}) $ is enough, where $L_{\varphi^{(2)}}$ is given by the previous lemma), for any target function $t:\mathbb{S}^1 \rightarrow \mathbb{R}$, for any $\theta \in \mathbb{S}^1$, there exists a set $D_\theta(t)  \subset \mathcal{A}(\lambda)$ with $\# D_{\theta}(t) \geq (1-\lambda^{-\gamma_{sep}}) \#\mathcal{A}(\lambda)$ such that: 
$$ \forall a \in D_\theta(t), \ \forall \mathbf{b} \in \mathcal{A}^*, \ |\ell_{\mathbf{b}a}(\theta) -t(\theta)| \geq \frac{1}{\sqrt{\lambda}}. $$
\end{lemma}

\begin{proof}
Define 
$$D_\theta(t) := \Big\{a \in \mathcal{A}(\lambda), \ |\ell_a(\theta)-t(\theta)| > \frac{2}{\sqrt{\lambda}}\Big\}.$$
Then indeed, for $\lambda$ large enough, we have for all $a \in D_\theta(t)$ and for all $\mathbf{b} \in \mathcal{A}^*$, $|\ell_{\mathbf{b}a}(\theta)-t(\theta)| \geq \frac{1}{\sqrt{\lambda}}$ since  $ |\ell_{\mathbf{b}a}(\theta) - \ell_a(\theta)| \leq  {\|\varphi^{(2)}\|_{\infty}} \sum_{j=1}^\infty (\mu \lambda^2)^{-j} \leq  \frac{\|\varphi^{(2)}\|_{\infty}}{\lambda}\leq \frac{1}{\sqrt{\lambda}} $ if $\lambda \geq \|\varphi^{(2)}\|_{\infty}^2$, since $\mu \lambda \geq 2$. 
Moreover, notice that since $\varphi$ is Lipschitz, we have when $|\nu-x_a|\leq \frac{1}{\lambda}$, for $\lambda$ large enough ($\geq \|\varphi\|_{C^3}^2 $ is enough):
$$ \mathbb{1}_{[t(\theta)-2\lambda^{-1/2},t(\theta)+2\lambda^{-1/2}]}(\ell_a(\theta)) \leq \mathbb{1}_{[t(\theta)-3\lambda^{-1/2},t(\theta)+3\lambda^{-1/2}]}(\varphi^{(2)}(\nu))  .$$
Integrating in $\nu$ yields 
$$ \lambda^{-1} \mathbb{1}_{[t(\theta)-2\lambda^{-1/2},t(\theta)+2\lambda^{-1/2}]}(\ell_a(\theta)) \leq |\{ \nu \in [x_a,x_a+\lambda^{-1}], \ \varphi^{(2)}(\nu) \in [t(\theta)-3\lambda^{-1/2},t(\theta)+3\lambda^{-1/2}] \}| . $$
We can then bound the cardinal of $\mathcal{A}(\lambda) \setminus D_\theta(t)$ as follows, summing over $a \in \mathcal{A}(\lambda)$ the previous equation:
$$ \lambda^{-1} \#\Big\{a \in \mathcal{A}(\lambda), \ |\ell_a(\theta)-t(\theta)| \leq \frac{2}{\sqrt{\lambda}}\Big\}  $$ $$ \leq |\{ \nu \in \mathbb{S}^1, \ \varphi^{(2)}(\nu) \in [t(\theta)-3\lambda^{-1/2},t(\theta)+3\lambda^{ -1/2}] \}| \leq \lambda^{-\gamma_{sep}} $$
for some $\gamma_{sep} \in (0,1)$ that depends only on $\varphi$, by the previous lemma, taking $\lambda$ sufficiently large, where we used the fact that $\varphi^{(2)}$ is analytic, 1-periodic, and non-constant. Hence $\#D_\theta(t)\geq (1-\lambda^{-\gamma_{sep}})\#\mathcal{A}(\lambda)$.
\end{proof}


\begin{lemma}[Tree lemma]
Let $t:\mathbb{S}^1 \rightarrow \mathbb{R}$ be any target function.
We have, for any $\lambda$ large enough (so that the previous lemma applies), for all $n \geq 1$ and for all $\sigma>0$:
$$ \lambda^{-n} \#\{ \mathbf{a} \in \mathcal{A}(\lambda)^n, \ |\ell_\mathbf{a}(\theta) - t(\theta)|<\sigma \} \leq \lambda (\sigma^{\gamma_{sep}/2}+\lambda^{-\gamma_{sep} n}). $$
where $\gamma_{sep}$ was introduced in the previous lemma and depends only on $\varphi$.
\end{lemma}

Here, as before, the \say{tree structure} should be understood as the \say{self-similarity} relation 
$$\ell_{\mathbf{ab}}(\theta) = \ell_\mathbf{b}(\theta)+ \frac{\ell_\mathbf{a}(g_\mathbf{b}(\theta))}{(\mu \lambda^2)^{|\mathbf{b}|}},$$ which allows to think of the set $(\ell_{\mathbf{a}}(\theta))_{\mathbf{a} \in \mathcal{A}^\infty}$ as a (family of) self similar sets in $\mathbb{R}$. The target $t$ will be later chosen of the form $\ell_\mathbf{b}(\theta)$ for some fixed word $\mathbf{b}$.

\begin{proof}
Intuitively, the proof is an induction on the scales $\sigma$. Formally, this is actually an induction on $n$, and then a separation on cases depending on the values of $\sigma.$ \\
When $n=0$ or $n=1$, the bound is clear, since we have the trivial bound $$\lambda^{-n} \#\{ \mathbf{a} \in \mathcal{A}(\lambda)^n, \ |\ell_\mathbf{a}(\theta) - t(\theta)|<\sigma \} \leq 1 \leq \lambda\left(\sigma^{\gamma_{sep}/2}+\lambda^{-n\gamma_{sep}}\right).$$


Now suppose by induction that the result holds for $n$, and let us prove it for $n+1$. If $\sigma \in (\lambda^{-1/2},1)$, the bound is also trivial. So let us assume that $\sigma \leq \lambda^{-1/2}$. Now, notice, using the set $D_\theta(t)$ defined in the previous lemma:
\begin{align*}
&\lambda^{-(n+1)} \#\{ \mathbf{a}b \in \mathcal{A}(\lambda)^{n+1}, \ |\ell_{\mathbf{a}b}(\theta) - t(\theta)|<\sigma \} \\
&\quad = \lambda^{-1} \sum_{b \in \mathcal{A}\setminus D_\theta(t)} \lambda^{-n} \# \{ \mathbf{a} \in \mathcal{A}(\lambda)^{n}, \ |\ell_{\mathbf{a}b}(\theta) - t(\theta)|<\sigma \} \\
&\quad = \lambda^{-1} \sum_{b \in \mathcal{A}\setminus D_\theta(t)} \lambda^{-n} \# \Big\{ \mathbf{a} \in \mathcal{A}(\lambda)^{n}, \ \Big|\frac{\ell_{\mathbf{a}}(g_b(\theta))}{\mu \lambda^2} + \ell_b(\theta) - t(\theta)\Big|<\sigma \Big\} \\
&\quad = \lambda^{-1} \sum_{b \in \mathcal{A}\setminus D_\theta(t)} \lambda^{-n} \# \Big\{ \mathbf{a} \in \mathcal{A}(\lambda)^{n}, \ |\ell_{\mathbf{a}}(g_b(\theta)) + (\ell_b(\theta)-t(\theta))\mu\lambda^2|<\sigma\mu\lambda^2 \Big\} \\
&\quad \leq \lambda^{-1} \sum_{b \in \mathcal{A}\setminus D_\theta(t)} \lambda \left( (\sigma\mu\lambda^2)^{\gamma_{sep}/2} + \lambda^{-n\gamma_{sep}} \right) \\
&\quad \leq \lambda^{-\gamma_{sep}} \cdot \lambda \left( (\sigma\mu\lambda^2)^{\gamma_{sep}/2} + \lambda^{-n\gamma_{sep}} \right) \\
&\quad \leq \lambda \left( \sigma^{\gamma_{sep}/2} + \lambda^{-(n+1)\gamma_{sep}} \right).
\end{align*}
\end{proof}

We are then ready to prove the $L^2$ estimates, using the same strategy as earlier.

\begin{theorem}
There exists $\rho \in (0,1)$ that depends only on $\varphi$ such that, for all $\beta \in (0,1/2]$, for all $\lambda$ sufficiently large and all $t$ sufficiently large depending on $\beta,\lambda$,, we have, uniformly in $m$,
$$
\forall h \in C^1, \qquad
\int_{\mathbb{S}^1} |\mathcal{L}_{\lambda,i\Phi(t,m)}^{\langle n_\beta(t) \rangle} h|^2 d\theta
\leq \frac{\|h\|_{1,t}^2}{t^{\rho}}.
$$
The condition on $\lambda$ is of the form $\lambda \geq \max(C_1(\varphi),\mu^{-C_2(\varphi)/\beta}).$
\end{theorem}


\begin{proof}[Proof (of $L^2$-Dolgopyat estimates)]
Let $\beta \in (0,1/2]$ and recall that $n := n_\beta(t)$ is such that $\mu^{n_\beta(t)} \sim t^{-\beta} \geq t^{-1/2}$. We suppose that $\lambda$ is large enough so that the previous tree lemma applies, and that $\lambda \geq \mu^{-10/\beta}.$ We start by writing (denoting $\Phi = \Phi(t,m)$ for clarity):
\begin{align*} 
\Big\| \mathcal{L}_{\lambda,\Phi(t,m)}^{\langle n \rangle}(h) \Big\|_{L^2([0,1])}^2 &= \int_0^1 |\mathcal{L}_{\lambda,i\Phi(t,m)}^{\langle n \rangle}(h)(\theta)|^2 d\theta \\
& = \int_0^1 \Big| \lambda^{-n} \sum_{\mathbf{a}} e^{i S_n \Phi \circ g_{\mathbf{a}}(\theta)} h(g_\mathbf{a})(\theta) \Big|^2 d\theta \\
& = \lambda^{-2n} \sum_{\mathbf{a},\mathbf{b} \in \mathcal{A}^n(\lambda)} \int_0^1 e^{i \big( S_n \Phi \circ g_\mathbf{a}(\theta) - S_n \Phi \circ g_\mathbf{b}(\theta) \big)} h(g_\mathbf{a}(\theta)) \overline{h(g_\mathbf{b}(\theta))} d\theta.
\end{align*}
Let us then iterate a little more our transfer operator: recall that we have the relation $\int h d\theta = \int \mathcal{L}_\lambda^{\mathfrak{n}}(h) d\theta$ for any $\mathfrak{n}$. We then choose $\mathfrak{n}$ such that $\lambda^{2\mathfrak{n}} \simeq t^{1/8}$ and we write:
$$ \Big\| \mathcal{L}_{\lambda,\Phi(t,m)}^{\langle n \rangle}(h) \Big\|_{L^2([0,1])}^2 = \frac{1}{\lambda^{2n+2 \mathfrak{n}}} \underset{\mathbf{c} \in \mathcal{A}(\lambda)^{2\mathfrak{n}}}{\sum_{\mathbf{a},\mathbf{b} \in \mathcal{A}(\lambda)^n}} \int_0^1 e^{i \big( S_n \Phi \circ g_\mathbf{a}(g_\mathbf{c}(\theta)) - S_n \Phi \circ g_\mathbf{b}(g_\mathbf{c}(\theta)) \big)} h(g_\mathbf{ac}(\theta)) \overline{h(g_\mathbf{bc}(\theta))} d\theta. $$
Since $h$ is $C^1$ and $g_{\mathbf{ac}}' = \lambda^{-n-2\mathfrak{n}} \leq t^{-2}$ for $t$ large enough (as $\lambda \geq \mu^{-10/\beta}$),, this gives, writing $h(g_{\mathbf{ac}}(\theta)) = h(x_{\mathbf{ac}})+\mathcal{O}(\|h\|_{C^1} t^{-2})$:
$$ \Big\| \mathcal{L}_{\lambda,\Phi(t,m)}^{\langle n \rangle}(h) \Big\|_{L^2([0,1])}^2 \leq \frac{\|h\|_\infty^2}{\lambda^{2n+2\mathfrak{n}}} \underset{\mathbf{c} \in \mathcal{A}(\lambda)^{2\mathfrak{n}}}{\sum_{\mathbf{a},\mathbf{b} \in \mathcal{A}(\lambda)^n}} \Big| \int_0^1 e^{i \big( S_n \Phi \circ g_\mathbf{a}(g_\mathbf{c}(\theta)) - S_n \Phi \circ g_\mathbf{b}(g_\mathbf{c}(\theta)) \big)}  d\theta \Big| + \frac{\|h\|_{{1,t}}^2}{t} .$$
Let us then concentrate on the integral $I_{\mathbf{a},\mathbf{b},\mathbf{c}}(t) := \int_0^1 e^{i \big( S_n \Phi \circ g_\mathbf{a}(g_\mathbf{c}(\theta)) - S_n \Phi \circ g_\mathbf{b}(g_\mathbf{c}(\theta)) \big)}  d\theta$. We have
$$ I_{\mathbf{a},\mathbf{b},\mathbf{c}}(t) =  \int_0^1 e^{i T \psi_{\mathbf{a},\mathbf{b},\mathbf{c}}(\theta)}  d\theta $$
with $T=t \mu^{n-1} \lambda^{-5\mathfrak{n}-2} \geq t^{1/8}$ and 
$$\psi_{\mathbf{a},\mathbf{b},\mathbf{c}}(\theta) = t^{-1} \mu^{1-n} \lambda^{5\mathfrak{n}+2} \big( S_n \Phi \circ g_\mathbf{a}(g_\mathbf{c}(\theta)) - S_n \Phi \circ g_\mathbf{b}(g_\mathbf{c}(\theta)) \big).$$

Again the classical van Der Corput lemma ensures that $|I_{\mathbf{a},\mathbf{b},\mathbf{c}}(t)| \leq c T^{-1/2}$ (for some universal constant $c$) if $|\psi_{\mathbf{a},\mathbf{b},\mathbf{c}}^{(2)}(\theta)| \geq 1$ on $[0,1]$. We hence get the bound:
$$ \Big\| \mathcal{L}_{\lambda,\Phi(t,m)}^{\langle n \rangle}(h) \Big\|_{L^2([0,1])}^2 \lesssim \|h\|_{1,t}^2 \Big( \frac{1}{t^{1/16}} + \lambda^{-2n-2\mathfrak{n}} \#\{ (\mathbf{a},\mathbf{b},\mathbf{c}) \in \mathcal{A}^{2n+2\mathfrak{n}} \ | \ \exists \theta, |\psi_{\mathbf{a},\mathbf{b},\mathbf{c}}^{(2)}(\theta)| \leq 1\} \Big). $$
To conclude, we need to bound this cardinality. We first compute:
\begin{align*}
\psi_{\mathbf{a},\mathbf{b},\mathbf{c}}^{(2)}(\theta) &=: \lambda^{\mathfrak{n}} \big( \ell_{\mathbf{a}}(g_\mathbf{c}(\theta)) - \ell_{\mathbf{b}}(g_\mathbf{c}(\theta)) \big)  \\
& = \lambda^{\mathfrak{n}} \big( \ell_{\mathbf{a}}(g_\mathbf{c}(0)) - \ell_{\mathbf{b}}(g_\mathbf{c}(0)) \big) + \mathcal{O}(\lambda^{-\mathfrak{n}}) .
\end{align*}
It follows that, for $t$ large enough, we can write, using the tree lemma:
\begin{align*}
& \lambda^{-2n-2\mathfrak{n}} \#\{ (\mathbf{a},\mathbf{b},\mathbf{c}) \in \mathcal{A}^{2n+2\mathfrak{n}} \ | \ \exists \theta, |\psi_{\mathbf{a},\mathbf{b},\mathbf{c}}^{(2)}(\theta)| \leq 1\}\\
&\quad \leq \lambda^{-2n-2\mathfrak{n}} \#\{ (\mathbf{a},\mathbf{b},\mathbf{c}) \in \mathcal{A}^{2n+2\mathfrak{n}} \ , \ \lambda^{\mathfrak{n}} \big| \ell_{\mathbf{a}}(g_\mathbf{c}(0)) - \ell_{\mathbf{b}}(g_\mathbf{c}(0)) \big| \leq 2\} \\
&\quad \leq \lambda^{-n-2\mathfrak{n}} \sum_{(\mathbf{b},\mathbf{c}) \in \mathcal{A}^{n+2\mathfrak{n}}} \lambda^{-n} \#\{ \mathbf{a} \in \mathcal{A}^n, \ \ell_\mathbf{a}(g_\mathbf{c}(0)) \in [\ell_\mathbf{b}(g_\mathbf{c}(0))-2\lambda^{-\mathfrak{n}},\ell_\mathbf{b}(g_\mathbf{c}(0))+2\lambda^{-\mathfrak{n}} ] \} \\
&\quad \leq 2 \lambda \Big( \lambda^{- \gamma_{sep} \mathfrak{n}/2} + \lambda^{-n \gamma_{sep}} \Big) \lesssim t^{-\gamma_{sep}/40}.
\end{align*}
Since also the van der Corput contribution is $\lesssim t^{-1/16}$, we obtain
$$
\int_{\mathbb S^1}
\left|\mathcal L_{\lambda,i\Phi(t,m)}^{\langle n_\beta(t)\rangle}h\right|^2\,d\theta
\lesssim
t^{-\rho}\|h\|_{1,t}^2,
\qquad
\rho:=\min\left\{\frac1{16},\frac{\gamma_{sep}}{40}\right\}>0.
$$
Thus $\rho$ depends only on $\varphi$.
\end{proof}

This establishes the desired  Dolgopyat's estimates, and fast dispersion for strongly irregular Weierstrass crystals.




\section{Limits on the rate of dispersion}  \label{sec:maximal-dispersion}

A natural question is how rapidly quantum dispersion can occur in long-range crystals. Also what are the conditions on the hoppings that guarantee some minimal rate of dispersion? Our main results provide arbitrarily fast polynomial decay in suitable regimes. One may therefore ask whether exponential decay is possible. The first result of this section shows that the answer is negative.

\subsection{Maximal rate of dispersion}

We now prove Proposition \ref{thm:noexponential-intro}, recall it here:

\begin{proposition}\label{thm:exponential}
Let $H$ be any crystal on $\mathbb Z$. Then $\langle \delta_0,e^{-itH}\delta_0\rangle$ cannot decay exponentially as $|t|\to\infty$. In particular, exponential dispersion is impossible.
\end{proposition}
\begin{proof}
The proof is a combination of the following steps:
    \begin{enumerate}
        \item If $\mu$ is the spectral measure of the crystal $H$, then by the spectral theorem, $\langle \delta_0,e^{-itH}\delta_0\rangle = \int_{\sigma(H)} e^{-itx}\,d\mu(x) = \widehat{\mu}(t)$.
        \item Since exponential decay implies $\widehat{\mu}\in L^1(\mathbb R)$, Fourier inversion shows that $\mu$ is absolutely continuous with a continuous density $g$.
        \item So assume the spectrum is AC with spectral density $g \in L^2(\mathbb R)$. Then $\langle \delta_0,e^{-itH}\delta_0\rangle = \widehat{g}(t)$.
        \item If $\widehat{g}(t)$ decays exponentially, say $|\widehat{g}(t)|\leq Ce^{-c|t|}$, then the Paley-Wiener theorem implies that $g$ is the restriction to $\mathbb R$ of a function holomorphic in every strip $\{z\in\mathbb C:|\operatorname{Im}z|<b\}$ with $b<c$, see \cite[Chapter~4]{SteinShakarchiComplex}.
        \item This is impossible because, since $H$ is bounded, the spectral density $g$ has compact support. Hence $g$ vanishes on some open interval of $\mathbb R$. Its holomorphic extension therefore vanishes on a set with an accumulation point in the strip, and hence vanishes identically by holomorphicity. This contradicts with $\int g\,dx=\widehat{\mu}(0)=\|\delta_0\|_2^2=1$.
    \end{enumerate}
    
    \end{proof}

No $L^2$ version of the Paley-Wiener theorem is needed in the proof above. Since $\widehat{g}$ decays exponentially, $\widehat{g}\in L^1(\mathbb R)$, while $g\in L^1(\mathbb R)$ as a spectral density. Fourier inversion therefore applies, and the exponential decay makes the inversion integral converge locally uniformly in a complex strip, giving directly the required holomorphic extension of $g$.


\begin{comment}
\begin{lemma}\label{lem:fdimpliesac}
Let $\mu \in \text{Prob}(\mathbb{R})$ be a compactly supported Borel probability measure. Suppose that $\widehat{\mu} \in L^1(\mathbb{R})$. Then there exists a continuous and bounded function $g:\mathbb{R} \rightarrow \mathbb{R}$ such that $d\mu(x) = g(x) dx$.
\end{lemma}

\begin{remark}
It is easily seen that $g$ can then be computed as $g(x) = \lim_{\varepsilon \rightarrow 0} \mu((x-\varepsilon,x+\varepsilon))/(2 \varepsilon)$.
It follows that the function $g$ has to be nonnegative, with support equal to the support of $\mu$, and in $L^1(\mathbb{R})$.  In particular, $\mu$ is absolutely continuous with respect to the Lebesgue measure.
\end{remark}

\begin{proof}
Let us define the continuous and bounded function $g$ by the formula $g(x) := \int_{\mathbb{R}} \widehat{\mu}(\xi) e^{2 i \pi x \xi} d\xi$. We will check that $\mu$ and $d\nu(x) := g(x)dx$ are the same tempered distributions. To do so, we will compute their Fourier transform in $S'(\mathbb{R})$ and check that they are the same. By density of the Schwartz space $\mathcal{S}(\mathbb{R})$ in $C^0(\mathbb{R})$, and by Riesz representation theorem for compactly supported measures in $\mathbb{R}$, this will imply that they are the same measures. \\

Let $f \in \mathcal{S}(\mathbb{R})$ be in the Schwartz space. We have, by definition of the Fourier transform in $\mathcal{S}'$ and by Fubini:
$$ \langle \widehat{\mu},f \rangle_{\mathcal{S}',\mathcal{S}} = \langle \mu, \widehat{f} \rangle_{\mathcal{S}',\mathcal{S}} = \int_\mathbb{R} \widehat{f}(x) d\mu(x)  = \int_\mathbb{R} \int_{\mathbb{R}} f(\xi) e^{-2 i \pi x \xi} d\xi d\mu(x) $$ $$ = \int_\mathbb{R} f(\xi) \int_{\mathbb{R}}  e^{-2 i \pi x \xi}  d\mu(x) d\xi = \int_\mathbb{R} f(\xi) \widehat{\mu}(\xi) d\xi. $$
This is just saying that the Fourier transform in $\mathcal{S}'$ of $\mu$ can be identified with the $L^1$ function $\widehat{\mu}$.
Similarly, by Fubini and the formula for the inverse Fourier transform, we find:
$$ \langle \widehat{g(x)dx},f \rangle_{\mathcal{S}',\mathcal{S}} = \langle g(x)dx, \widehat{f} \rangle_{\mathcal{S}',\mathcal{S}}= \int_\mathbb{R} \widehat{f}(x) g(x) dx$$
$$ = \int_\mathbb{R} \widehat{f}(x) \int_\mathbb{R} \widehat{\mu}(\xi) e^{2 i \pi x \xi} d\xi dx = \int_\mathbb{R} \widehat{\mu}(\xi) \int_\mathbb{R} \widehat{f}(x) e^{2 i \pi x \xi} dx d\xi =  \int_\mathbb{R} \widehat{\mu}(\xi) f(\xi) d\xi.$$
The Fourier transforms are the same in $\mathcal{S}'(\mathbb{R})$, and the proof is done.
%\end{proof}

\begin{remark}
It is a very general fact that any compactly supported distribution $T$ have a Fourier transform that can be identified with an analytic function, with formula $\widehat{T} = \langle T, e^{2 i \pi \xi \cdot} \rangle.$ Furthermore, for compactly supported measures, we always have, in the sense of distributions:
$$ d\mu(x) = \int_\mathbb{R} e^{2 i \pi \xi x} \widehat{\mu}(\xi)  d\xi$$
where the integral is a potentially diverging oscillatory integral, which should be understood as the distribution $$ \langle T, \varphi \rangle := \int_\mathbb{R} \int_\mathbb{R} e^{2 i \pi \xi x} \varphi(x) \widehat{\mu}(\xi) dx d\xi .$$
For example, $\delta_0 = \int_\mathbb{R} e^{2 i \pi \xi x} d\xi$ in this sense (this is a rewriting of the inverse Fourier transform formula actually). The spirit of the previous proof is just that the oscillatory integral for $d\mu(x)$ is converging when $\widehat{\mu}$ is in $L^1$ and defines naturally a continuous function.
\end{remark}
\end{comment}

The argument leaves open the possibility of dispersion faster than every polynomial but slower than exponential. There is no analogous Fourier-analytic obstruction in this regime: if $g\in C_c^\infty(\mathbb R)$, then repeated integration by parts gives $|\widehat g(t)|\leq C_N(1+|t|)^{-N}$ for every $N\geq1$. More strongly, the Beurling-Malliavin multiplier theorem \cite{BM} gives compactly supported functions with sub-exponential Fourier decay, for instance of the form
$$
\exp\left(-c\frac{|t|}{(\log |t|)^2}\right).
$$
By Favard's theorem, compactly supported probability measures with infinite support can be realized as spectral measures of bounded Jacobi matrices; see \cite{KillipSimon}. For crystals, however, the inverse problem is more restrictive: in the one-vertex translation-invariant case the spectral measure must have the form $h_*\mathcal L$ for a periodic Floquet function $h$, and dispersion moreover requires estimates uniform in the off-diagonal Fourier mode. This leads to the following question.

\begin{question}
Can one construct a long-range crystal exhibiting super-polynomial but sub-exponential quantum dispersion?
\end{question}

Our construction suggests one possible approach. Let $\lambda_n\to\infty$ be a strictly increasing sequence of integers and consider the non-stationary Weierstrass phase
$$
W(\theta)=\sum_{n=0}^\infty \mu^n\varphi(\lambda^{\langle n\rangle}\theta),
\qquad
\lambda^{\langle n\rangle}:=\lambda_1\cdots\lambda_n.
$$
Increasing the lacunarity with the scale could produce progressively stronger oscillatory cancellation and hence super-polynomial dispersion. Alternatively, one could ask whether ideas from the Beurling-Malliavin construction can be adapted to produce Floquet functions of long-range crystals with such decay.

\subsection{Regularity obstructions to fast dispersion}

So far we have studied how increasingly irregular Floquet functions can produce fast dispersion. We now turn to the converse question: to what extent does regularity of the Floquet function limit the possible rate of dispersion? We first show that even local H\"older regularity gives an obstruction. More precisely, if the Floquet function is $C^\alpha$ on any open interval, then the dispersive norm cannot be $o(t^{-1/\alpha})$. Thus faster dispersion requires the Floquet function to be correspondingly irregular everywhere.

\begin{proposition}\label{thm:maxdis}
Let $h : \mathbb{R}\rightarrow \mathbb{R}$ be $1$-periodic and continuous. Suppose that there exists some open set $U \subset \mathbb{S}^1$ such that $h_{|U} \in C^\alpha(U,\mathbb{R})$ for $\alpha \in (0,1]$. Then
$$
\limsup_{t \rightarrow \infty} \Big(t^{1/\alpha}\sup_{m \in \mathbb{Z}} \Big| \int_0^1 e^{i t h(\theta)} e^{-2 i \pi m \theta} d\theta \Big| \Big)> 0.
$$
\end{proposition}

In particular, a crystal whose Floquet function is locally $C^\alpha$ cannot disperse strictly faster than $t^{-1/\alpha}$. Notice that the assumption is only local: a single interval on which $h$ has this regularity already obstructs faster decay. This complements Theorem~\ref{thm:main}, where increasingly rough Weierstrass Floquet functions produce increasingly fast polynomial dispersion.

\begin{proof}[Proof of Proposition \ref{thm:maxdis}]
Suppose, towards a contradiction, that there exists $\varepsilon_0(t) \rightarrow 0$ as $t \rightarrow \infty$ such that $$ \sup_{m \in \mathbb{Z}} \Big| \int_0^1 e^{i t h(\theta)} e^{-2 i \pi m \theta} d\theta \Big| \leq \varepsilon_0(t) t^{-1/\alpha}. $$
Let $\rho \in C^\infty_c(\mathbb{R},\mathbb{R}^+)$ be a smooth nonnegative bump function supported on $[-1/4,1/4]$, with integral one. Choose $a \in U$ and denote the associated approximation of $\delta_a$: $\rho_{a,\varepsilon} := \varepsilon^{-1} \rho\big((\cdot-a)/\varepsilon\big)$. First, we have the bound
$$ \Big|\int_{\mathbb{S}^1} e^{it h(\theta)} \rho_{a,\varepsilon}(\theta) d\theta - e^{it h(a)} \Big| \leq  \int_{\mathbb{S}^1} |e^{it h(\theta)} - e^{i t h(a)} | \rho_{a,\varepsilon}(\theta) d\theta \leq t \sup_{\theta \in [a-\varepsilon,a+\varepsilon]}|h(\theta)-h(a)| \lesssim t \varepsilon^\alpha. $$
In particular, $$ \Big|\int_{0}^{1} e^{it h(\theta)} \rho_{a,\varepsilon}(\theta) d\theta \Big| \geq 1-C t\varepsilon^\alpha.$$
Second, we have:
$$ \Big|\int_{0}^{1} e^{it h(\theta)} \rho_{a,\varepsilon}(\theta) d\theta \Big|= \Big|\sum_{m \in \mathbb{Z}} c_m(\rho_{a,\varepsilon}) \int_{-1/2}^{1/2} e^{i t h(\theta)} e^{2i\pi m \theta} d\theta \Big|$$
$$ \lesssim \varepsilon_0(t) t^{-1/\alpha} \|c(\rho_{a,\varepsilon})\|_{\ell^1(\mathbb{Z})} $$
with $$ \|c(\rho_{a,\varepsilon})\|_{\ell^1(\mathbb{Z})} = \sum_{m \in \mathbb{Z}} |\widehat{\rho_{a,\varepsilon}}(m)| = \sum_{\nu \in \varepsilon \mathbb{Z}} |\widehat{\rho}(\nu)| \lesssim \varepsilon^{-1} \|\widehat{\rho}\|_{L^1(\mathbb{R})}. $$
Hence $$ \Big|\int_{-1/2}^{1/2} e^{it h(\theta)} \rho_{a,\varepsilon}(\theta) d\theta \Big| \lesssim \varepsilon_0(t) t^{-1/\alpha} \varepsilon^{-1} .$$
Choosing $\varepsilon :=  (2Ct)^{-1/\alpha}$ and letting $t \rightarrow +\infty$ yields a contradiction.
\end{proof}

We next apply this observation to crystals with monotone weights in $H$. Suppose that the weights are eventually decreasing, $\omega(n+1)\leq \omega(n)$ for all sufficiently large $n$. We show that this mild assumption already forces the Floquet function to be $C^1$ away from the integers, and hence Proposition~\ref{thm:maxdis} gives a universal $1/t$ obstruction to dispersion.

\begin{lemma}
If $u_n \geq 0$ is eventually decreasing and satisfies $\sum_n u_n < \infty$, then $u_n = o(1/n)$. Moreover, the sequence $\varepsilon_n := n u_n$ satisfies $\sum_{n} |\varepsilon_n-\varepsilon_{n+1}|<\infty$.
\end{lemma}

\begin{proof}
We have $2n u_{2n} \leq u_n+\dots+u_{2n} \leq \sum_{k=n}^{\infty} u_k \rightarrow 0$. Then $n u_n \rightarrow 0$ by monotonicity of $u_n$. Now, letting $\varepsilon_n := n u_n$, We have $$ \sum_{n=0}^N |\varepsilon_n - \varepsilon_{n+1}| = \sum_{n=0}^N |n u_n - (n+1) u_{n+1}| \leq \sum_{n=0}^N n ( u_n -  u_{n+1}) + \sum_{n=0}^N u_{n+1} $$
which converge since, by summation by parts, $$ \sum_{n=0}^N n ( u_n -  u_{n+1}) \leq n u_n + \sum_{n} {u_n} \leq \varepsilon_n + \|u\|_{\ell^1}  $$
is bounded.
\end{proof}

\begin{theorem}\label{thm:monotone-regularity}
Let $(\omega(n))_{n\in\mathbb Z}$ be crystal weights and suppose that $(\omega(n))_{n\geq0}$ is eventually decreasing. Then the associated Floquet function $h$ satisfies $h\in C^0(\mathbb S^1,\mathbb R)\cap C^1((0,1),\mathbb R)$. In particular, by symmetry, $h'(1/2)=0$.
\end{theorem}

\begin{proof}
We have $h(\theta) := \sum_{n \in \mathbb{Z}} \omega(n) e^{2 i \pi n \theta}$ for $\theta \in \mathbb{R}$. Since $\sum_n \omega(n) = 1$ and $\omega(n) \geq 0$, this function is continuous, bounded, periodic on $\mathbb{R}$. To see that it is differentiable on $\mathbb{R} \setminus \mathbb{Z}$, we do an Abel transform. Rewrite $h$ as $$ h(\theta) = \omega(0) + 2\sum_{n=1}^\infty \omega(n) (-1)^n + 2 \sum_{n=1}^\infty \omega(n) (\cos(2 \pi n \theta) - (-1)^n). $$ Then, notice that, for all $N \geq 1$, we have (denoting $\varepsilon_n = n \omega(n)$)
$$  \sum_{n=1}^N \omega(n) (\cos(2 \pi n \theta)-(-1)^n) = \sum_{n=1}^N \varepsilon(n) \frac{\cos(2 \pi n \theta)-(-1)^n}{n} $$ $$= \varepsilon(N) S_N(\theta) - \sum_{n=1}^{N-1} (\varepsilon(n+1)-\varepsilon(n)) S_n(\theta) $$
where $S_n(\theta) = \sum_{j=1}^n \frac{\cos(2 \pi j \theta)-(-1)^j}{j} $. Notice that, for $\theta \in (0,1)$, we have $$S_n(\theta) = - \sum_{j=1}^n \frac{1}{2 \pi }\int_{1/2}^\theta \sin(2 \pi j \nu) d\nu = -\frac{1}{2 \pi }\int_{1/2}^\theta  \sum_{j=1}^n \sin(2 \pi j \nu) d\nu = -\int_{1/2}^\theta  s_n(\nu) d\nu, $$
where $$s_n(\nu) := \frac{1}{2\pi} \sum_{j=1}^n \sin(2 \pi j \nu) =   \frac{ \sin(\pi n \nu) \sin(\pi(n+1)\nu)}{2 \pi  \sin(\pi \nu)} .$$
We have $s_n(\theta) \leq \frac{C}{\theta} + \frac{C}{1-\theta} $ for some $C \leq 1$ independent on $n$ and $\theta \in (0,1)$. Hence the same bound holds for $S_n(\theta)$. Added with the fact that $\varepsilon(n)-\varepsilon(n+1)$ is summable, we get convergence locally uniformly of the series:
$$ \forall \theta \in (0,1), \quad h(\theta) = \omega(0) + 2\sum_{n=1}^\infty \omega(n) (-1)^n + \sum_{n=1}^\infty \Big(\varepsilon(n)-\varepsilon(n+1) \Big) S_n(\theta) .$$
The bound on $s_n(\theta)$ allows us to exchange the sum and the integral and we find
$$ \forall \theta \in (0,1), \quad h(\theta) =  \omega(0) + 2\sum_{n=1}^\infty \omega(n) (-1)^n + 2 \int_{1/2}^\theta \sum_{n=1}^\infty \Big(\varepsilon(n)-\varepsilon(n+1) \Big) s_n(\nu) d\nu . $$
In particular, $h \in C^{1}((0,1),\mathbb{R})$, and for any $\theta \in (0,1)$, we have
$$ h'(\theta) = 2 \sum_{n=1}^\infty (\varepsilon_n-\varepsilon_{n+1}) s_n(\theta) .$$\end{proof}

Combining Theorem~\ref{thm:monotone-regularity} with Proposition~\ref{thm:maxdis} for $\alpha=1$ immediately gives:

\begin{corollary}\label{cor:monotone}
If the crystal weights $(\omega(n))_{n\geq0}$ are eventually decreasing, then dispersion cannot be faster than $1/t$.
\end{corollary}

Moreover, for every nontrivial crystal in this class, the spectral measure has a nonzero absolutely continuous component. Indeed, since $h$ is nonconstant and $C^1$ on $(0,1)$, there exists an interval on which $h'\neq0$, and the push-forward of Lebesgue measure restricted to this interval is absolutely continuous.

\section{Transport and spectrum} \label{sec:further}

We conclude with further consequences of the roughness of the Weierstrass Floquet functions to the transport and spectrum of the corresponding long-range crystals.

\subsection{Speed of transport}

In addition to dispersion, \textit{transport} is another widely studied notion in quantum dynamics. Here we are interested in two natural questions:
\begin{itemize}
\item What is the speed of transport for Weierstrass crystals, in particular in regimes where the spectrum is purely singular continuous?
\item Can a crystal exhibit genuinely sub-ballistic transport?
\end{itemize}
The transport is said to be \textit{ballistic} if
$$
\lim_{t\to\infty}\frac{\|Xe^{itH}\psi\|}{t}=c>0
$$
for suitably localized initial states $\psi$, for instance compactly supported $\psi$. Here $X$ denotes the position operator,
$$
(X\psi)(n):=n\psi(n),\qquad n\in\mathbb Z.
$$
More generally, one may consider initial states satisfying $\|X\psi\|<\infty$.

For non-locally finite crystals, both ballistic motion and much faster transport can occur. In \cite{KernerPostSabriTaufer}, the following two phenomena were established:
\begin{enumerate}
\item Under suitable regularity assumptions on the Floquet eigenvalues, the transport is ballistic. In particular, in the one-vertex setting considered here, if the Floquet function $h$ belongs to $W^{1,2}(\mathbb T^d)$, then the transport of $\delta_0$ is ballistic.
\item For the fractional Laplacian on $\mathbb Z$, sufficiently rough Floquet functions can instead lead to instantaneous divergence of the second moment. More precisely, for $H=(-\Delta)^\alpha$ with $\alpha\leq1/4$,
$$
\|Xe^{itH}\delta_0\|=+\infty
$$
for every $t\neq0$.
\end{enumerate}

The same Sobolev mechanism can be applied directly to the Weierstrass crystals considered in this paper. We record the short argument below, as it gives a simple relation between the regularity of the Floquet function and the speed of transport.

\textbf{The proof of (2)} relies on the following observation: if $\|x e^{itH}\delta_0\|$ is finite, let $\psi_t = e^{itH}\delta_0$. Then by hypothesis, $\sum_n n^2 |\psi_t(n)|^2<\infty$. But $\psi_t(n) = \langle \delta_n, e^{itH}\delta_0\rangle = \langle F^{-1}\delta_n,F^{-1}e^{itH}\delta_0\rangle=\int_0^1 e^{-2\pi i n\theta} e^{ith(\theta)} d\theta=\widehat{\varphi}_t(n)$, where $\varphi_t(\theta)=e^{ith(\theta)}$.

In other words, 
\[
\|x e^{itH}\delta_0\|<\infty \iff \varphi_t(\theta)=e^{ith(\theta)}\in W^{1,2}(\mathbb{T}^d) \,.
\] 
We show that for the fractional Laplacian, $\varphi_t(\theta)$ is in $W^{1,1}(\mathbb{T})$ but not in $W^{1,2}(\mathbb{T})$.

This argument will work for any crystal having the same property. This suggests that the Weierstrass crystal with singular continuous spectrum explodes in finite time as well.

\textbf{The proof of (1)} is a bit lengthy in the paper because it follows from a more general theorem. For the special case of (1), one can argue as follows. Let's simplify even further and take $d=1$ (this is unnecessary). We are assuming that $h\in W^{1,2}$. Since exponentiation is smooth, then $\varphi_t=e^{ith}\in W^{1,2}(\mathbb{T})$ as well for any $t$. So it has a weak derivative which is by definition $(\partial_w \varphi_t)(\theta)=2\pi i \sum_n n\widehat{\varphi_t}(n)e^{2\pi i n\theta}$, and $\partial_w \varphi_t\in L^2$.

On the other hand, 
\[
\sum_n e^{2\pi i n\theta} n \widehat{\varphi_t}(n) = \sum_n e^{2\pi i n\theta} n \int_0^1 e^{-2\pi i nx}e^{ith(x)}dx = \sum_n e^{2\pi i n\theta} n (e^{itH}\delta_0)(n) = \widehat{xe^{it H}\delta_0}(\theta) \,.
\]
So we get $\widehat{xe^{itH}\delta_0}\in L^2$ and
\[
\widehat{xe^{itH}\delta_0}(\theta)= \frac{1}{2\pi i}(\partial_w\varphi_t)(\theta)=\frac{1}{2\pi i}t(\partial_w h)(\theta)e^{ith(\theta)} \,.
\]
So $\|xe^{itH}\delta_0\|<\infty$ for any $t$ and
\begin{equation}\label{e:conspeed}
\frac{\|xe^{itH}\delta_0\|}{t} = \frac{\|\widehat{xe^{itH}\delta_0}\|}{t} = \frac{\|\partial_w h\|}{2\pi}
\end{equation}
for any $t$, in particular for the limit. (I didn't realize before that this is constant in $t$).

When you start from $\psi\neq \delta_0$, you need the limit, because lower order terms (in $t$) appear in the chain rule, involving derivatives of $\widehat{\psi}(\theta)$.

Note that from the previous calculation, as long as $\|xe^{itH}\delta_0\|<\infty$, then $\varphi_t=e^{ith}\in W^{1,2}$ and $\frac{\| x e^{itH}\delta_0\|}{t}=\frac{\|\partial_w \varphi_t\|}{2\pi t}$. When $h$ is $W^{1,2}$, we get \eqref{e:conspeed}.


\begin{lemma}
For $\varphi$ fixed and $\mu$ fixed, for $\lambda$ large enough, $W_{\mu,\lambda} \notin C^{1/2}(\mathbb{S}^1)$.
\end{lemma}

\begin{proof}
The local study of $W$ done before proves that $W$ has oscillations like $x^\alpha$ for $\alpha:=|\ln(\mu)|/\ln(\lambda)$ everywhere and at every scale. If $\alpha<1/2$, $W$ can not be $1/2$-Hölder. See appendix \ref{ap:B} for details.
\end{proof}

\begin{lemma}
$W^{1,2}(\mathbb{S}^1) \subset C^{1/2}(\mathbb{S}^1)$.
\end{lemma}



\begin{proof}
Let $h \in W^{1,2}(\mathbb{S}^1)$; by definition, its derivative in the sense of distributions is in $L^2(\mathbb{S}^1)$. By the fundamental theorem of calculus in $W^{1,2}$, we can write:
$$ h(x)-h(y) = \int_x^y h'(t)dt.$$
Then, by Cauchy-Schwarz:
$ |h(x)-h(y)| \leq \|h'\|_{L^2(\mathbb{S}^1)} |x-y|^{1/2} $.
\end{proof}

\begin{lemma}
If $\lambda$ is large enough depending on $\mu,\varphi$, for all $t \neq 0$, $e^{i t W_{\mu,\lambda}} \notin W^{1,2}$.
\end{lemma}

\begin{proof}
The function $e^{it W}$ can not be $1/2$-Hölder, since $e^{it \theta}$ is a local analytic diffeomorphism (for $t \neq 0$) so that $W$ would be itself $1/2$-Holder. Hence $e^{it W} \notin W^{1,2}(\mathbb{S}^1)$ for $t \neq 0$.
\end{proof}

\begin{corollary}\label{cor:transport}
The Weierstrass crystal blows up in finite time as soon as $|\ln(\mu)|/\ln(\lambda) < 1/2$.
More generally, any crystal with Floquet function that is not $1/2$-Hölder blows up in finite time.
\end{corollary}

%\subsection{Further directions and open problems}

%\textcolor{red}{This needs integrating to introduction and adding discussion}
%Then include our findings:
%\begin{itemize}
%\item If the Floquet function $h$ is $C^1$, then the spectrum of $H$ cannot be purely singularly continuous. This uses Sard's theorem.
%\item Can we have some singularly continuous spectrum with a regular $h$ ? Currently the only way to exclude it is if $h$ is analytic.
%\item What if the coefficients $w(k)$ are decreasing ? Can we say something about the spectra or dynamics ?
%\item Check if we can use results on random Fourier series to say something about crystals constructed at random.
%\item Review the proof of Pitt et.al. regarding singular continuous spectrum. See if it encompasses the Weierstrass function, and maybe some random functions.
%\end{itemize}

%Additionnal Questions:

%\begin{itemize}
%\item What is the fastest speed possible ?
%\item What happens if arrows are spaced out ? If the serie is lacunary ? Does it speed up ?
%\item Is there always some absolutely continuous spectrum ? Is it possible for the spectrum to be purely singular ? Purely atomic ?

%\textcolor{magenta}{The answer to all these questions is known from the crystal paper: the spectrum can be purely singularly continuous (Thm 6.1) and the top of the spectrum cannot be an eigenvalue, in particular the spectrum cannot be purely atomic. Here the positivity of the weights is important. If you allow complex hermitian weights on your crystal, it can be pure atomic.} \textbf{That's very interesting !!}
%\item Is it possible to have dense atoms in the spectrum ? (Comparing with Anderson localization...)
%\item \textcolor{magenta}{What is the speed of transport for the Weierstrass crystal with purely singularly continuous spectrum ?}

%\textcolor{magenta}{Update: We now know that blow up occurs in finite time in the large $\lambda$ regime.}
%\item \textcolor{magenta}{What is the range of a Weierstrass function? It is bounded in $[\frac{\|\varphi\|_\infty}{1-\mu},\frac{\|\varphi\|_{\infty}}{1-\mu}]$, but is there equality, particularly in case of $\varphi = \cos 2\pi \cdot$ ? This would give the spectrum of the corresponding crystal as a set.}
%\item \textcolor{magenta}{}
%\end{itemize}



\subsection{Spectrum}

We conclude with some observations on the spectral type of the Weierstrass crystals. Although our main results concern dispersion, the spectral question is closely related: the spectral measure of $\delta_0$ is precisely the occupation measure $(W_{\mu,\lambda})_*d\theta$ of the Floquet function
$$
W_{\mu,\lambda}(x)=\sum_{k\geq0}\mu^k\cos(2\pi\lambda^k x).
$$
There are two regimes in which the spectral type is understood. When $\lambda\mu<1$, the spectrum is purely absolutely continuous, while for $\lambda$ sufficiently large relative to $\mu^{-1}$ our Fourier-decay theorem again gives purely absolutely continuous spectrum, despite the much lower regularity of $W_{\mu,\lambda}$. The most interesting case is the critical parameter $\lambda\mu=1$, where we expect purely singular continuous spectrum.

Indeed, if $\lambda\mu<1$, then $W_{\mu,\lambda}\in C^1$ with
$$
W'_{\mu,\lambda}(x)=-2\pi\sum_{k\geq0}(\lambda\mu)^k\sin(2\pi\lambda^k x).
$$
By the classical uniqueness theorem for lacunary series \cite[Theorem 6]{Zyg32}, $W'_{\mu,\lambda}$ cannot vanish on a set of positive measure unless it vanishes identically. Hence $W'_{\mu,\lambda}\neq0$ almost everywhere, and \cite[Theorem 6.3]{KernerPostSabriTaufer} implies that the spectrum is purely absolutely continuous. At the other extreme, Theorem~\ref{thm:main} gives sufficiently fast Fourier decay whenever $\lambda$ is sufficiently large relative to $\mu^{-1}$. Once the decay exponent exceeds $1$, Fourier inversion implies that $(W_{\mu,\lambda})_*d\theta$ has a continuous density, and hence the spectrum is purely absolutely continuous.

At the critical case $\lambda\mu=1$ the situation changes. Recall that $\Lambda^\ast$ denotes the Zygmund class, and let $\Lambda^\ast_0\subset\Lambda^\ast$ denote the small Zygmund class, consisting of functions satisfying $
f(x+t)+f(x-t)-2f(x)=o(|t|)$, $t\to0.$ A classical criterion of Pitt \cite{AndersonHousworthPitt} states that if a real-valued periodic function $h$ belongs to the small Zygmund class $\Lambda^\ast_0$ and is almost everywhere non-differentiable, then $M_h$ has purely singular continuous spectrum. This criterion gives that the related examples
$$
h(x)=\sum_n a_n2^{-n}\cos(2\pi2^nx),\qquad
a_n\to0,\qquad \sum_n a_n^2=\infty,
$$
considered in \cite{KernerPostSabriTaufer}. It does not apply directly to the classical Weierstrass function. One has
$$
W_{\mu,\lambda}\in
\begin{cases}
\Lambda^\ast_0, & \lambda\mu<1,\\
\Lambda^\ast\setminus\Lambda^\ast_0, & \lambda\mu=1.
\end{cases}
$$
whereas for $\lambda\mu<1$ the function is already $C^1$. Thus Pitt's criterion narrowly misses precisely the critical case.

There is nevertheless evidence for singular continuity at criticality. As we mentioned in the introduction, the local time of $W_{1/2,2}$ defined with $\varphi(\theta)=\sin(2\pi\theta)$ does not exist \cite{Buc}, that is, the image of the Lebesgue measure under $W_{1/2,2}$ is singular. Also almost every level set is finite. On the other hand, Theorem~\ref{thm:main} gives dispersion for the corresponding multiplication operator, showing in particular that dispersion does not by itself imply absolute continuity of the spectral measure, although this operator is not a crystal of the class considered here. Moreover, the arguments behind Pitt-type results suggest studying the derivatives of the partial sums: sufficiently high recurrence of these derivatives can be used to construct a full-measure set whose image has zero Lebesgue measure. See Appendix \ref{ap:B}, in particular Proposition \ref{prop:density}, for related properties. For the critical Weierstrass function these derivatives are governed by lacunary ergodic sums, suggesting a possible route to proving singularity. We thus have the following diagram for the spectrum of long-range crystals arising from $W_{\mu,\lambda}$:
\begin{itemize}
\item if $\lambda<\mu^{-1}$, the spectrum is purely absolutely continuous;
\item if $\lambda=\mu^{-1}$, we conjecture that the spectrum is purely singular continuous;
\item if $\lambda$ is sufficiently large relative to $\mu^{-1}$, the spectrum is purely absolutely continuous.
\end{itemize}
It is natural to ask more generally whether the spectrum is purely absolutely continuous throughout the supercritical regime $\lambda>\mu^{-1}$. If the critical conjecture holds, $\lambda\mu=1$ would give a particularly interesting example of a long-range crystal exhibiting both singular continuous spectrum and quantum dispersion.




\appendix

\section{A $C^1$ Van Der Corput Lemma}\label{ap:A}


It is useful for us to have at our disposition a general version of Van Der Corput lemma. We prove a version here, which might be of independent interest. A similar result (although only controlling $m = 0$) was proved by Algom, Chang, Wu, Wu \cite{ACWW} in a more general setting. We prove our version for completeness and because the proof is quite elementary in the case of the Lebesgue measure. 


\begin{theorem}[$C^{1}$ Van Der Corput]\label{thm:c1+alphavandercorput}
Let $h:[0,1] \rightarrow \mathbb{R}$ be a $C^{1}$ function. Suppose that $h$ is sufficiently nonlinear, in the following sense:
$$ \forall y \in \mathbb{R}, \ \lambda( \theta \in [0,1], \ h'(\theta)= y ) =0. $$
Then, for any continuous function $\chi:[0,1] \rightarrow \mathbb{C}$, we have:
$$   \sup_{m \in \mathbb{R}} \Big| \int_{0}^1 e^{2 i \pi (\xi h(\theta) - m \theta)} \chi(\theta) d\theta \Big| \underset{|\xi| \rightarrow \infty}{\longrightarrow} 0. $$
\end{theorem}




\begin{proof}[Proof of Theorem \ref{thm:c1+alphavandercorput}]
Choose $\omega:(0,1] \rightarrow (0,\infty)$ a concave modulus of continuity for $\chi$ and $h'$ (one can choose for example the concave shell of any common modulus of continuity). By concavity, and since $\omega(0^+)=0^+$, $\omega$ is increasing, continuous and bijective near zero. Moreover, concavity yields $\omega(x/2) \geq \omega(x)/2$. For all $|\xi| \geq 1$, set $n(\xi) := \max \{ n \geq 1 \ | \ \xi^2 4^{-n} \omega(2^{-n}) \geq 1 \}$. Obviously we have $n(\xi) \underset{|\xi| \rightarrow \infty}{\longrightarrow} \infty$. Moreover, since $$ 1 \geq \xi^{2} 4^{-(n(\xi)+1)} \omega(2^{-(n(\xi)+1)}) \geq \xi^{2} 4^{-n(\xi)} \omega(2^{-n(\xi)})/8 \geq 1/8 ,$$
we find the growth relation
$$ 2^{-n(\xi)} \sqrt{\omega(2^{-n(\xi)})} \simeq 1/\xi .$$
Now that this is fixed, let us start the proof. We let $|\xi| \geq 1$ and we will write $n=n(\xi)$.
We will consider the numbers with finite binary digits: $$\forall \mathbf{a} \in \{0,1\}^{n}, \ x_\mathbf{a} := \sum_{k=1}^n a_k 2^{-k}.$$ We have:
$$\int_{0}^1 e^{2 i \pi (\xi h(\theta) - m \theta)} \chi(\theta) d\theta = \sum_{\mathbf{a} \in \{0,1\}^n} \int_{x_\mathbf{a}}^{x_{\mathbf{a}}+2^{-n}} e^{2 i \pi (\xi h(\theta) - m \theta)} \chi(\theta) d\theta  $$
$$ = 2^{-n} \sum_{\mathbf{a} \in \{0,1\}^n} \int_{0}^{1} e^{2 i \pi (\xi h(x_\mathbf{a}+\theta 2^{-n}) - m (x_\mathbf{a} + \theta 2^{-n}))} \chi(x_\mathbf{a}+\theta 2^{-n}) d\theta $$
$$ = 2^{-n} \sum_{\mathbf{a} \in \{0,1\}^n}  \chi(x_\mathbf{a}) e^{-2 i \pi m x_\mathbf{a}}\int_{0}^{1} e^{2 i \pi (\xi h(x_\mathbf{a}+\theta 2^{-n}) - m \theta 2^{-n})}  d\theta + \mathcal{O}(\omega(2^{-n})).$$
Now, since $h$ is $C^{1}$, we have by the mean value theorem, for all $\varepsilon$,
$$ h(x_{\mathbf{a}} + \varepsilon) = h(x_\mathbf{a})+h'(y(x_{\mathbf{a}},\varepsilon)) \varepsilon $$ for some $y \in [x_\mathbf{a},x_{\mathbf{a}}+\varepsilon]$, which gives the asymptotic expansion
$$ h(x_{\mathbf{a}} + \varepsilon) = h(x_\mathbf{a})+h'(x_\mathbf{a}) \varepsilon + \mathcal{O}(\varepsilon \  \omega(\varepsilon)), $$
from which we find
$$ \int_{0}^1 e^{2 i \pi (\xi h(\theta) - m \theta)} d\theta =  2^{-n} \sum_{\mathbf{a} \in \{0,1\}^n} \chi(x_\mathbf{a}) e^{-2 i \pi m x_\mathbf{a}}\int_{0}^{1} e^{2 i \pi (\xi h(x_\mathbf{a})+\xi h'(x_\mathbf{a}) \theta 2^{-n}  - m \theta 2^{-n})} e^{2 i \pi \xi \cdot \mathcal{O}(2^{-n} \omega(2^{-n}))} d\theta $$ $$ = 2^{-n} \sum_{\mathbf{a} \in \{0,1\}^n} \chi(x_\mathbf{a}) e^{2 i \pi (\xi h(x_\mathbf{a})-m x_\mathbf{a})} \int_{0}^{1} e^{2 i \pi \xi 2^{-n} ( h'(x_\mathbf{a}) - \frac{m}{\xi}) \theta} d\theta + \mathcal{O}\big(\sqrt{\omega(2^{-n})}\big) $$
Now, the inner integral is explicit. Denoting $\sigma^{-1} := \xi 2^{-n} \pi \simeq \omega(2^{-n})^{-1/2}$ and $\mu := m/\xi$, we compute:
$$ \Big|\int_{0}^{1} e^{2 i \sigma^{-1} ( h'(x_\mathbf{a}) - \mu) \theta} d\theta \Big| = \Big|\text{sinc}\Big(\frac{h'(x_\mathbf{a})-\mu}{\sigma} \Big)\Big| \leq \min\Big( 1 \ , \ \frac{\sigma}{|h'(x_\mathbf{a})-\mu|}  \Big). $$
We can then bound these integrals by either $1$ if $|h'(x_\mathbf{a})-\mu| \leq \sqrt{\sigma}$, or by $\sqrt{\sigma}$ if $|h'(x_\mathbf{a})-\mu| \geq \sqrt{\sigma}$.
Hence, we find the bound:
$$ \Big|\int_{0}^1 e^{2 i \pi (\xi h(\theta) - \mu \theta)} \chi(\theta) d\theta \Big| \leq  2^{-n} \sum_{\mathbf{a} \in \{0,1\}^n} \|\chi\|_\infty \Big| \int_{0}^{1} e^{2i  \sigma^{-1} ( h'(x_\mathbf{a}) - \mu) \theta} d\theta \Big| + \mathcal{O}\big(\sqrt{\omega(2^{-n})}\big) $$
$$ \lesssim \frac{1}{2^n}  \sum_{\mathbf{a} \in \{0,1\}^n} \min\Big( 1 \ , \ \frac{\sigma}{|h'(x_\mathbf{a})-\mu|}  \Big) + \mathcal{O}\big(\sqrt{\omega(2^{-n})}\big) $$
$$ \leq \frac{1}{2^n} \# \{ \mathbf{a} \in \{0,1\}^n, \ |h'(x_\mathbf{a}) - \mu| \leq \sqrt{\sigma} \} +   \mathcal{O}\big(\sqrt{\omega(2^{-n})}\big) .$$
To conclude, we use the regularity of $h'$, in the following way. We have, for all $\theta \in [0, 2^{-n}]$, the bound 
$$ |h'(x_\mathbf{a}) - h'(x_\mathbf{a}+\theta)| \leq \omega(\theta). $$
In particular, for $\mathbf{a} \in \{0,1\}$, we have:
$$ \Big(|h'(x_\mathbf{a}) - \mu| \leq \sqrt{\sigma} \Big) \Rightarrow \Big( \forall \theta \in [0,2^{-n}], \ |h'(x_\mathbf{a}+\theta)-\mu| \leq \sqrt{\sigma} + \omega(2^{-n}) \Big). $$
Which allows us to bound:
$$ \frac{1}{2^n} \# \{ \mathbf{a} \in \{0,1\}^n, \ |h'(x_\mathbf{a}) - \mu| \leq \sqrt{\sigma} \} = \lambda\Big( \underset{|h'(x_\mathbf{a}) - m|  \leq \sqrt{\sigma}}{\bigcup_{\mathbf{a} \in \{0,1\}^n}} [x_\mathbf{a},x_\mathbf{a}+2^{-n}] \Big) $$
$$ \leq \lambda\Big( \theta \in [0,1], \ |h'(\theta)-m| \leq C \omega(2^{-n})^{1/4}  \Big). $$
Now the pushforward measure $(h')_* \lambda$ is non-atomic by hypothesis, and is of compact support. It follows that it's cumulative distribution function is uniformly continuous, with modulus of continuity given by some $\tilde{\omega}$. We can then bound
$$ \lambda\Big( \theta \in [0,1], \ |h'(\theta)-m| \leq C \omega(2^{-n})^{1/4}  \Big) \leq \tilde{\omega}(C \omega(2^{-n})^{1/4}). $$
by our nonlinearity hypothesis on $h'$.
In the end, our oscillatory integral is bounded by:
$$ \Big|\int_{0}^1 e^{2 i \pi (\xi h(\theta) - m \theta)} \chi(\theta) d\theta \Big| \lesssim \tilde{\omega}(C \omega(2^{-n})^{1/4}) + \sqrt{\omega(2^{-n})}.$$
Which proves that $\Big|\int_{0}^1 e^{2 i \pi (\xi h(\theta) - m \theta)} \chi(\theta) d\theta \Big| \underset{|\xi| \rightarrow \infty} \longrightarrow 0$, and this uniformly in $m$.
\end{proof}

In fact, the previous proof gives an explicit rate of decay depending on the Frostman exponent of $(h')_* \lambda$ and the modulus of continuity of $\chi$ and $h'$. In the setting where $h \in C^{1+\alpha}$, we get polynomial decay under natural hypothesis.



\begin{theorem}[$C^{1+\alpha}$ Van Der Corput]\label{thm:c1+alphavandercorput}
Let $\alpha \in (0,1)$ and let $h:[0,1] \rightarrow \mathbb{R}$ be a $C^{1+\alpha}$ function. Suppose that $h$ is sufficiently nonlinear, in the following sense:
$$ \exists C_0, \ \exists \gamma \in (0,1), \ \forall I \subset \mathbb{R} \text{ segment}, \  \lambda( \theta \in [0,1], \ h'(\theta) \in I ) \leq C_0 |I|^\gamma .$$
Then there exists a constant $C \geq 1$ that depends only on $C_0$ such that, for any $\alpha$-Hölder function $\chi:[0,1] \rightarrow \mathbb{C}$, we have:
$$  \forall |\xi| \geq 1, \ \sup_{m \in \mathbb{R}} \Big| \int_{0}^1 e^{2 i \pi (\xi h(\theta) - m \theta)} \chi(\theta) d\theta \Big| \leq C \|h\|_{C^{1+\alpha}}  \|\chi\|_{C^\alpha} |\xi|^{-\alpha \gamma/6}. $$
\end{theorem}

\begin{proof}
Recall that the previous proof gives us a rate of decay of the form
$$ \sup_m \Big| \int_0^1 e^{2 i \pi \xi h} \chi d\theta \Big| \lesssim \tilde{\omega}\Big(C \omega(2^{-n(\xi)})^{1/4} \Big) + \omega(2^{-n(\xi)})^{1/2}$$
where $n(\xi)$ is fixed so that $\xi^2 4^{-n(\xi)} \omega(2^{-n(\xi)}) \simeq 1$.
In our current setting, we have $\tilde{\omega}(t) \simeq t^\gamma$ and $\omega(t) \simeq t^\alpha$, so that
$$ \xi^2 4^{-n(\xi)} \omega(2^{-n(\xi)}) \simeq \xi^{2} 2^{(2+\alpha)(n(\xi))} \simeq 1$$
gives $ 2^{-n(\xi)} \simeq \xi^{-\frac{2}{2+\alpha}} $,
and the bound becomes
$$ \tilde{\omega}\Big(C \omega(2^{-n(\xi)})^{1/4} \Big) + \omega(2^{-n(\xi)})^{1/2} \lesssim \xi^{-\frac{2\alpha \gamma}{4(2+\alpha)}} \lesssim \xi^{- \alpha \gamma/6}, $$
which is the announced bound.
\end{proof}

\begin{comment}
\begin{theorem}[$C^{1}$ van den Corput]\label{thm:c1+alphavandercorput}
Let $\alpha \in (0,1)$ and let $h:[0,1] \rightarrow \mathbb{R}$ be a $C^{1+\alpha}$ function. Suppose that $h$ is sufficiently nonlinear, in the following sense:
$$ \exists C_0, \ \exists \gamma \in (0,1), \ \forall I \subset \mathbb{R} \text{ segment}, \  \lambda( \theta \in [0,1], \ h'(\theta) \in I ) \leq C_0 |I|^\gamma .$$
Then there exists $C \geq 1$ \textbf{that only depends on...} such that, for any $\alpha$-Hölder function $\chi:[0,1] \rightarrow \mathbb{C}$, we have:
$$  \forall |\xi| \geq 1, \ \sup_{m \in \mathbb{R}} \Big| \int_{0}^1 e^{2 i \pi (\xi h(\theta) - m \theta)} \chi(\theta) d\theta \Big| \leq C \|\chi\|_{C^\alpha} |\xi|^{-\alpha \gamma/5}. $$
\end{theorem}




\begin{proof}[Proof of Theorem \ref{thm:c1+alphavandercorput}]
We will consider the numbers with finite binary digits: $$\forall \mathbf{a} \in \{0,1\}^{n}, \ x_\mathbf{a} := \sum_{k=1}^n a_k 2^{-k}.$$ We have:
$$\int_{0}^1 e^{2 i \pi (\xi h(\theta) - m \theta)} \chi(\theta) d\theta = \sum_{\mathbf{a} \in \{0,1\}^n} \int_{x_\mathbf{a}}^{x_{\mathbf{a}}+2^{-n}} e^{2 i \pi (\xi h(\theta) - m \theta)} \chi(\theta) d\theta  $$
$$ = 2^{-n} \sum_{\mathbf{a} \in \{0,1\}^n} \int_{0}^{1} e^{2 i \pi (\xi h(x_\mathbf{a}+\theta 2^{-n}) - m (x_\mathbf{a} + \theta 2^{-n}))} \chi(x_\mathbf{a}+\theta 2^{-n}) d\theta $$
$$ = 2^{-n} \sum_{\mathbf{a} \in \{0,1\}^n}  \chi(x_\mathbf{a}) e^{-2 i \pi m x_\mathbf{a}}\int_{0}^{1} e^{2 i \pi (\xi h(x_\mathbf{a}+\theta 2^{-n}) - m \theta 2^{-n})}  d\theta + \mathcal{O}(\|\chi\|_{C^\alpha} 2^{-n \alpha}).$$
\textbf{todo: adapt the proof with the bump and clean the $m/\xi$ thing.}Now, since $h$ is $C^{1+\alpha}$, we have the asymptotic expansion
$$ h(x_{\mathbf{a}} + \varepsilon) = h(x_\mathbf{a})+h'(x_\mathbf{a}) \varepsilon + \mathcal{O}(\varepsilon^{1+\alpha}), $$
from which we find
$$ \int_{0}^1 e^{2 i \pi (\xi h(\theta) - m \theta)} d\theta =  2^{-n} \sum_{\mathbf{a} \in \{0,1\}^n}  e^{-2 i \pi m x_\mathbf{a}}\int_{0}^{1} e^{2 i \pi (\xi h(x_\mathbf{a})+h'(x_\mathbf{a}) \theta 2^{-n}  - m \theta 2^{-n})} e^{2 i \pi \xi \cdot \mathcal{O}(2^{-(1+\alpha)n})} d\theta $$ $$ = 2^{-n} \sum_{\mathbf{a} \in \{0,1\}^n}  e^{2 i \pi (\xi h(x_\mathbf{a})-m x_\mathbf{a})} \int_{0}^{1} e^{2 i \pi \xi 2^{-n} ( h'(x_\mathbf{a}) - m) \theta} d\theta + \mathcal{O}(\xi 2^{-(1+\alpha)n}) .$$
\textbf{computation mistake here, should be $m/\xi$, but I don't think this matter much. }Now, the inner integral is explicit. Denoting $\sigma^{-1} := \xi 2^{-n} \pi$, we compute:
$$ \Big|\int_{0}^{1} e^{2 i \sigma^{-1} ( h'(x_\mathbf{a}) - m) \theta} d\theta \Big| = \Big|\text{sinc}\Big(\frac{h'(x_\mathbf{a})-m}{\sigma} \Big)\Big| \leq \min\Big( 1 \ , \ \frac{\sigma}{|h'(x_\mathbf{a})-m|}  \Big). $$
Fix now some $\rho \in (0,1)$. We can then bound these integrals by either $1$ if $|h'(x_\mathbf{a})-m| \leq \sigma^{1-\rho}$, or by $\sigma^{\rho}$ if $|h'(x_\mathbf{a})-m| \geq \sigma^{1-\rho}$.
Hence, we find the bound:
$$ \Big|\int_{0}^1 e^{2 i \pi (\xi h(\theta) - m \theta)} d\theta \Big| \leq  2^{-n} \sum_{\mathbf{a} \in \{0,1\}^n} \Big| \int_{0}^{1} e^{2i  \sigma^{-1} ( h'(x_\mathbf{a}) - m) \theta} d\theta \Big| + \mathcal{O}(\xi 2^{-(1+\alpha)n}) $$
$$ \leq \frac{1}{2^n}  \sum_{\mathbf{a} \in \{0,1\}^n} \min\Big( 1 \ , \ \frac{\sigma}{|h'(x_\mathbf{a})-m|}  \Big) + \mathcal{O}(\xi 2^{-(1+\alpha)n}) $$
$$ \leq \frac{1}{2^n} \# \{ \mathbf{a} \in \{0,1\}^n, \ |h'(x_\mathbf{a}) - m| \leq \sigma^{1-\rho} \} +   \mathcal{O}(\sigma^{\rho} + \xi 2^{-(1+\alpha)n}) .$$
To conclude, we use the $\alpha$-Hölder regularity of $h'$, in the following way. We have, for all $\theta \in [0, 2^{-n}]$, the bound \textcolor{magenta}{I suppose there is a typo and it should be $h'$ instead of $h$ in all the following.}
$$ |h(x_\mathbf{a}) - h(x_\mathbf{a}+\theta)| \leq C|\theta|^\alpha. $$
In particular, for $\mathbf{a} \in \{0,1\}$, we have:
$$ \Big(|h(x_\mathbf{a}) - m| \leq \sigma^{1-\rho} \Big) \Rightarrow \Big( \forall \theta \in [0,2^{-n}], \ |h(x_\mathbf{a}+\theta)-m| \leq \sigma^{1-\rho} + C 2^{-n \alpha} \Big). $$
Which allows us to bound:
$$ \frac{1}{2^n} \# \{ \mathbf{a} \in \{0,1\}^n, \ |h'(x_\mathbf{a}) - m| \leq \sigma^{1-\rho} \} = \lambda\Big( \underset{|h(x_\mathbf{a}) - m|  \leq \sigma^{1-\rho}}{\bigcup_{\mathbf{a} \in \{0,1\}^n}} [x_\mathbf{a},x_\mathbf{a}+2^{-n}] \Big) $$
$$ \leq \lambda\Big( \theta \in [0,1], \ |h'(\theta)-m| \leq \sigma^{1-\rho} + C 2^{-n \alpha}  \Big). $$
Now, by our nonlinearity hypothesis, this is bounded by
$$ \leq C_0 \Big( \sigma^{1-\rho} + C 2^{-n \alpha}  \Big)^\gamma = \mathcal{O}( \sigma^{(1-\rho)\gamma} + C 2^{-n \alpha \gamma}).   $$
In the end, our oscillatory integral is bounded by:
$$ \Big|\int_{0}^1 e^{2 i \pi (\xi h(\theta) - m \theta)} d\theta \Big| = \mathcal{O}\Big(  \sigma^{(1-\rho)\gamma} +  2^{-n \alpha \gamma} + \sigma^{\rho} + \xi 2^{-(1+\alpha)n} \Big) $$
$$ = \mathcal{O}\Big(  \sigma^{(1-\rho)\gamma} +  2^{-n \alpha \gamma} + \sigma^{\rho} +\sigma^{-1} 2^{-\alpha n} \Big).$$
We then choose $\rho := \frac{\gamma}{1+\gamma} \in (0,1/2)$, to find
$$ = \mathcal{O}\Big(  \sigma^{\frac{\gamma}{1+\gamma}} +  2^{-n \alpha \gamma} +\sigma^{-1} 2^{-\alpha n} \Big) .$$
Finally, we optimize by choosing $n$ such that $\sigma^{\frac{\gamma}{1+\gamma}}$ and $\sigma^{-1} 2^{-\alpha n}$ decays at the same speed: $\sigma \simeq 2^{- n \alpha \frac{1+\gamma}{1+2\gamma}}$ (that is, $\xi \simeq 2^{n} \sigma^{-1} \simeq 2^{n(1+\alpha \frac{1+\gamma}{1+2\gamma})}$ ). With this choice, the bound becomes
$$ = \mathcal{O}( 2^{- \frac{\alpha \gamma}{1+2\gamma} n} ). $$
Rewriting this in term of $\xi$, we finally find
$$ \Big|\int_{0}^1 e^{2 i \pi (\xi h(\theta) - m \theta)} d\theta \Big| \lesssim |\xi|^{- \frac{\alpha \gamma}{(1+2\gamma)(1+\alpha \frac{1+\gamma}{1+2\gamma})}} \lesssim |\xi|^{-\alpha \gamma/5}. $$
\textcolor{magenta}{Seems good, we just need to correct the $m/\xi$ and see if all goes well.}
\end{proof}
\end{comment}

Our proof gives a decay rate of the same order of magnitude as the usual Van Der Corput Lemma when $h$ is smooth, but with an additional factor of $1/6$. Notice that any periodic analytic non-constant function $h : \mathbb{S}^1 \rightarrow \mathbb{R}$ always satisfy that $(h')_*\lambda$ is $\gamma$-Frostman, for $\gamma^{-1} < \infty$ being the higher order of vanishing of $h'$. Consequently (this is also a direct consequence of the usual Van Der Corput Lemma):

\begin{corollary}
Any crystal associated to an analytic non-constant Floquet function exhibit polynomial dispersion.
\end{corollary}

\begin{corollary}
    Any locally finite crystal exhibits dispersion.
\end{corollary}

\section{On the local regularity of Weierstrass functions}\label{ap:B}

\subsection{Fractal oscillations of Weierstrass functions} 

The goal of this section is to do a detailled analysis of the oscillations of (non-analytic) Weierstrass maps. This appendix join some well known result about Weierstrass maps, and adapt some less known results adapted from \cite{Leclerc2,TZ23} (the detailed analysis on the oscillations of $X$ is very close to the core of the arguments used to study the temporal distance function in those two articles). We will prove the following regularity result; actually stronger statement will be described at the end of the section.

\begin{theorem}[Cohomology condition]{\label{thm:regWeierstrass}}
The Weierstrass map $W$ is analytic if and only if there exists $\psi \in C^\omega(\mathbb{S}^1,\mathbb{R})$ such that $\varphi(\theta)=\mu \psi(\lambda \theta)-\psi(\theta)$. Denote $\alpha \in (0,\infty)$ so that $\mu \lambda^\alpha=1$. If $W$ is not analytic, then: \begin{itemize}
\item When $\alpha \in \mathbb{N}$,  $W \in (\bigcup_{\beta<\alpha} C^\beta) \setminus C^\alpha$.
\item When $\alpha \notin \mathbb{N}$, $X \in C^\alpha \setminus (\bigcup_{\gamma>\alpha} C^{\gamma})$.
\end{itemize}
\end{theorem}

\begin{remark}
Let us do three examples. First, consider the case of $W(\theta) = \sum_{n=0}^\infty 2^{-n} \cos(2\pi 3^n \theta)$. There can not exists any analytic $\psi \in C^\omega(\mathbb{S}^1,\mathbb{R})$ such that $\cos(2 \pi \theta) = {\psi(3\theta)}/2-\psi(\theta)$. Indeed, if it was the case, then we would find  $2\pi \sin(2 \pi  \theta) = (3/2) \psi'(3\theta)-\psi'(\theta)$. Computing the first Fourier coefficient of this equation yields $c_1(\psi')=- i \pi$, and computing the $3^k$-Fourier coefficient, $k \geq 1$, gives the relation $c_{3^k}(\psi') = (3/2) c_{3^{k-1}}(\psi')$. It follows that $ c_{3^k}(\psi')  = - i \pi (3/2)^k$ which is not bounded, a contradiction. Hence $W \in C^{\ln(2)/\ln(3)} \setminus C^{\ln(2)/\ln(3)+}$. \\

As a second example, consider $W(\theta) = \sum_{n=0}^\infty 4^{-n} \cos(2\pi 2^n \theta)$. There can not exists any analytic $\psi \in C^\omega(\mathbb{S}^1,\mathbb{R})$ such that $\cos(2 \pi \theta) = {\psi(2\theta)}/4-\psi(\theta)$. The previous proof applies again by taking the third derivative of the equation and studying the Fourier coefficients of $\psi''$. Hence $\psi \in C^{2-} \setminus C^2$. \\

Last, consider $\sum_{n=0}^\infty 2^{-n} \sin(2 \pi 2^n \theta)$. There can not exists any analytic $\psi \in C^\omega(\mathbb{S}^1,\mathbb{R})$ such that $\cos(2 \pi \theta) = {\psi(2\theta)}/2-\psi(\theta)$. Indeed, in this case we would find $2\pi \cos(2 \pi \theta) = \psi'(2\theta)-\psi'(\theta)$, which gives $2\pi = 2\pi \cos(0) = \psi'(0)-\psi'(0)=0$, a contradiction. Hence $\psi \in C^{1-} \setminus C^1$.
\end{remark}

Actually, all the Weierstrass maps that appear in our \say{crystal} setting are non-analytic.

\begin{lemma}
Fix a non-constant $\varphi$ of the form $\varphi(\theta) = \sum_{j=0}^J c_j \cos(2 \pi j \theta) $ for $0<J<\infty$ and $c_j \geq 0$ . Let $\lambda \in \mathbb{N} \setminus \{0,1\}$ and $\mu \in (0,1)$. Then $W(\theta) := \sum_{n=0}^\infty \mu^n \varphi(\lambda^n \theta)$ is not analytic.
\end{lemma}

\begin{proof}
Suppose that there exists an analytic $\psi \in C^\omega(\mathbb{S}^1,\mathbb{R})$ such that $\psi(\theta)=\varphi(\theta)+\mu \psi(\lambda \theta)$. Choose $k=4k'$ large enough so that $\mu \lambda^k \geq 2$. We have $\psi^{(k)}(\theta)=\varphi^{(k)}(\theta)+\mu \lambda^k \psi^{(k)}(\lambda \theta)$. Notice that $\varphi^{(k)}(\theta) = \sum_{j=0}^J c_j (2\pi j)^k \cos(2 \pi j \theta)$ has nonnegative Fourier coefficients. Consider $j_0$, the smallest non-zero integer such that $c_{j_0}>0$. Studying the Fourier coefficients yields $c_{j}(\psi^{(k)}) = 0$ for $0 \leq j<j_0$ and then $c_{j_0}(\psi^{(k)}) = c_{j_0} (2\pi j_0)^k$.
Notice then that, for $n \geq 1$, we have $c_{\lambda^n j_0}(\psi^{(k)}) = c_{\lambda^n j_0}(\varphi^{(k)})+ (\mu \lambda^k) c_{\lambda^{n-1}j_0}(\psi^{(k)}) \geq 2 c_{\lambda^{n-1}j_0}(\psi^{(k)})$
so that $c_{\lambda^n j_0}(\psi^{(k)}) \geq 2^n c_{j_0} (2 \pi j_0)^k$ which is not bounded in $n$, a contradiction.
\end{proof}


Now let us now prove Theorem \ref{thm:regWeierstrass}. To study the regularity of $X$ close to a point $x$, we introduce the centered notation:
$$X_x(t)=\sum_{n=0}^\infty \mu^n \varphi(\lambda^n(x+t))$$
and $\varphi_x(t) := \varphi(x+t)$.
Notice that family of centered functions $(X_x)_{x \in \mathbb{S}^1}$ is autosimilar:
$$X_x(t)=\sum_{n=0}^\infty \mu^n \varphi(\lambda^n(x+t)) = \varphi(x+t) + \mu \sum_{n=1}^\infty \mu^{n-1} \varphi(\lambda^n(x+t)) $$
$$ = \varphi(x+t) + \mu \sum_{n=0}^\infty \mu^{n} \varphi(\lambda^n( \lambda x+ \lambda t)) = \varphi_x(t) + \mu X_{\lambda x}(\lambda t) .$$
The fundamental relation $$ X_x(t) = \varphi_x(t) + \mu X_{\lambda x}(\lambda t)$$
will allow us to relate oscillations at large scales with oscillations at small scales. We will also use the notation, for $\beta \in (0,1)$: $\delta^\beta f(\theta) = \frac{f(\theta)-f(0)}{|\theta|^\beta} $. Clearly, $f$ is $\beta$-Hölder at $0$ iff $\delta^\beta f$ is bounded. We start by proving the regularity result, which is easier.

\begin{lemma}
Recall that $\alpha \in (0,\infty)$ is such that $\mu \lambda^\alpha=1$. If $\alpha \notin \mathbb{N}$, then $X \in C^\alpha$. If $\alpha \in \mathbb{N}$, then $X \in C^{\alpha-\varepsilon}$ for all $\varepsilon$.
\end{lemma}

\begin{proof}
Write $\alpha=k+\beta$ for some $\beta \in (0,1]$ and $k \geq 0$ an integer. Since $\mu \lambda^k = \lambda^{-\beta}<1$, the series $X(\theta) = \sum_{n=0}^\infty \mu^n \varphi(\lambda^n \theta)$ clearly defines a $C^k$ function, with $k$-th derivative $X^{(k)}(\theta) = \sum_{n=0}^\infty \lambda^{-\beta n} \varphi^{(k)}(\lambda^n \theta)$. It is also easy to see that $X^{(k)}$ is $\beta'$-Hölder regular, for any $\beta'=\beta-\varepsilon$. Indeed, taking $\delta^{\beta'}$ of the previous formula gives immediately, for any $\theta,x \in \mathbb{S}^1$ with $\theta \neq 0$:
$$\delta^{\beta'} X_x^{(k)}(\theta) = \sum_{n=0}^\infty \lambda^{-\varepsilon n} \delta^{\beta'} \varphi_x^{(k)}(\theta)  $$
which is clearly a bounded function since $\varphi_x^{(k)}$ is $\beta'$-Hölder regular (uniformly in $x$) . \\
Suppose now that $\alpha \notin \mathbb{N}$, so that $\alpha = k+\beta$ with $\beta \in (0,1)$. Let us show that $X^{(k)}$ is actually $C^\beta$. The autosimilarity relation $X^{(k)}_x(\theta) = \varphi^{(k)}_x(\theta)+\lambda^{-\beta} X_{\lambda x}^{(k)}(\lambda \theta)$ gives
$$ \delta^\beta X^{(k)}_x(\theta) = \delta^\beta \varphi^{(k)}_x(\theta)+ \delta^\beta X_{\lambda x}^{(k)}(\lambda \theta). $$
We iterate the relation until $\lambda^N \theta \simeq 1$:
$$ \delta^\beta X^{(k)}_x(\theta) = \sum_{n=0}^N \delta^\beta \varphi^{(k)}_x(\theta \lambda^n)+ \delta^\beta X_{\lambda x}^{(k)}(\lambda^N \theta)  \lesssim \sum_{n=0}^N |\theta \lambda^n|^{1-\beta} +  \| X^{(k)} \|_\infty \lesssim 1 .$$
Hence the family of functions $(\partial^\beta X_x^{(k)
})_{x \in \mathbb{S}^1}$ is uniformly bounded, and $X^{(k)}$ is $\beta$-Hölder.
\end{proof}

Now, we will begin a carefull study of the local oscillations of $X$. To do so, we need to introduce some dynamical background. 

\begin{definition}
Denote $\mathcal{A} := \{0,\dots,\lambda-1\}$ and consider the compact space $\mathcal{A}^\mathbb{Z}$, equipped with the (ultrametric) distance $d(\mathbf{a},\mathbf{b}) = \lambda^{- N(\mathbf{a},\mathbf{b})}$ where $N(\mathbf{a},\mathbf{b}) = \sup\{ n \geq 0 \ | \ \forall |j| \leq n, \ a_j = b_j \} \in \mathbb{N} \cup \{\infty\}$. Consider the full shift, defined as  
 $$ \begin{array}[t]{lrcl}
\sigma : & \mathcal{A}^\mathbb{Z} & \longrightarrow & \mathcal{A}^\mathbb{Z} \\
    & (a_n)_{n \in \mathbb{Z}}  & \longmapsto &  (a_{n+1})_{n \in \mathbb{Z}}  \end{array} $$
Notice that $\sigma$ is Lipchitz. The multiplication by $\lambda \in \mathbb{N} \setminus \{0,1\}$, denoted $$ \begin{array}[t]{lrcl}
m : & \mathbb{S}^1 & \longrightarrow & \mathbb{S}^1 \\
    & x \text{ mod } 1  & \longmapsto &  \lambda x \text{ mod } 1  \end{array} $$
is an expanding map with $\lambda$ inverse branches, that we denote, for $a \in \mathcal{A}$, $$ \begin{array}[t]{lrcl}
g_a : & [0,1) \subset \mathbb{S}^1 & \longrightarrow & [\frac{a}{\lambda},\frac{a+1}{\lambda}) \subset \mathbb{S}^1 \\
    & x  & \longmapsto &  \frac{a+x}{\lambda}  \end{array} $$
We further denote $g_{a_1 \dots a_n} := g_{a_1} \dots g_{a_n} : [0,1) \rightarrow \mathbb{S}^1$. It is well known that $m$ is a factor of $\sigma$. Indeed, the surjective projection $$\begin{array}[t]{lrcl}
 \pi : & \mathcal{A}^\mathbb{Z} & \longrightarrow & \mathbb{S}^1 \\
    & (a_n)_{n \in \mathbb{Z}}  & \longmapsto & \underset{n \rightarrow +\infty}{\lim} g_{a_1\dots a_n}(x_0)  \end{array} $$ 
is independent of $x_0$, Lipschitz, and a direct computation yields $\pi \circ \sigma = m \circ \pi$. Notice also that, writing $\mathbf{a} = (a_n)_{n \in \mathbb{Z}}$ and $x=\pi(\mathbf{a})$, we have $\pi(\sigma^{-k} \mathbf{a}) = g_{a_{-k} \dots a_{-1}}(x)$. Finally, define the family of functions $(X_\mathbf{a})_{\mathbf{a} \in \mathcal{A}^\mathbb{Z}}$ by $X_{\mathbf{a}}(\theta) := X_{\pi(\mathbf{a})}(\theta) = X(\pi(\mathbf{a})+\theta)$ and denote similarly $\varphi_\mathbf{a}(\theta) := \varphi_{\pi(\mathbf{a})}(\theta) = \varphi(\pi(\mathbf{a})+\theta)$.
\end{definition}


We now turn into studying the local oscillations of $X$ when $X$ is not analytic. We study in detail the local oscillations of $X$. The fact that $X$ is not $C^{\alpha+}$ (or $C^\alpha$ if $\alpha \in \mathbb{N}$) will be a consequence of this careful study. The idea is the following. For $x \in \mathbb{S}^1$ and $\alpha \in (0,\infty) \setminus \mathbb{N}$, a function $f$ is $C^\alpha$ at $x$ if there exists a polynomial $P$ of order $\lfloor \alpha \rfloor$ such that $|f(x+\theta)-P(\theta)| \leq C |\theta|^\alpha$. It is then natural to study the distance of $X$ with polynomials of degree $\lfloor \alpha \rfloor$ as we zoom into a point $x$. To help with this program, we will approximate $(X_{\mathbf{a}})_{\mathbf{a} \in \mathcal{A}^\mathbb{Z}}$ with a family of function $(Z_{\mathbf{a}})_{\mathbf{a} \in \mathcal{A}^\mathbb{Z}}$  that have more suited self-similar properties.

\begin{theorem}
Let $N$ be the smallest integer such that $\mu \lambda^N > 1$. That is, if $\alpha \notin \mathbb{N}$, $N:= \lceil \alpha \rceil $, and if $\alpha \in \mathbb{N}$, $N = \alpha+1$.
There exists a family of functions from $\mathbb{R}$ to $\mathbb{R}$, $(Y_{\mathbf{a}})_{\mathbf{a} \in \mathcal{A}^\mathbb{Z}},(Z_{\mathbf{a}})_{\mathbf{a} \in \mathcal{A}^\mathbb{Z}},(P_{\mathbf{a}})_{\mathbf{a} \in \mathcal{A}^\mathbb{Z}}$, such that :
\begin{itemize}
\item $X_{\mathbf{a}}=Y_{\mathbf{a}}+Z_{\mathbf{a}}$
\item $Y_{\mathbf{a}}$ is analytic and $Y_{\mathbf{a}}(0)=Y_{\mathbf{a}}'(0)=\dots = Y_{\mathbf{a}}^{(N-1)}(0)$
\item $P_{\mathbf{a}}$ is a polynomial of degree $\leq N-1$
\item $Z_{\mathbf{a}}(\theta) = P_{\mathbf{a}}(\theta) + \mu Z_{\mathbf{a}}(\lambda \theta)$
\end{itemize}
Furthermore, the families $(X_{\mathbf{a}})_{\mathbf{a} \in \mathcal{A}^\mathbb{Z}},(Y_{\mathbf{a}})_{\mathbf{a} \in \mathcal{A}^\mathbb{Z}},(Z_{\mathbf{a}})_{\mathbf{a} \in \mathcal{A}^\mathbb{Z}},(P_{\mathbf{a}})_{\mathbf{a} \in \mathcal{A}^\mathbb{Z}}$ are continuous in $\mathbf{a}$ into the space of continuous/analytic/continuous/polynomials of degree N-1.
\end{theorem}

\begin{proof}
For every $\mathbf{a} \in \mathcal{A}^\mathbb{Z}$, define the analytic map $Y_\mathbf{a}$ as the only function such that $Y_{\mathbf{a}}(0)=Y_{\mathbf{a}}'(0)=\dots = Y_{\mathbf{a}}^{(N-1)}(0) = 0$ and $$Y_{\mathbf{a}}^{(N)}(\theta) := -\sum_{n=1}^\infty \frac{\varphi_{ \sigma^{-n} (\mathbf{a})}^{(N)}(\theta \lambda^{-n}) }{(\mu \lambda^N)^n} .$$ This series clearly defines an analytic map in $\theta \in \mathbb{R}$. Inspecting the series yield $$ \forall \mathbf{a} \in \mathcal{A}^\mathbb{Z}, \forall \theta \in \mathbb{R}, \quad  Y_{\mathbf{a}}^{(N)}(\theta) = - \frac{\varphi^{(N)}_{ \sigma^{-1}(\mathbf{a})}(\theta \lambda^{-1})}{\mu \lambda^N} + \frac{Y^{(N)}_{ \sigma^{-1}(\mathbf{a})}(\theta \lambda^{-1})}{\mu \lambda^N},$$
which can be rewritten as
$$ \forall \mathbf{a} \in \mathcal{A}^\mathbb{Z}, \forall \theta \in \mathbb{R}, \quad  Y_{\mathbf{a}}^{(N)}(\theta) =  {\varphi^{(N)}_{ \mathbf{a}}(\theta)} + \mu \lambda^N {Y^{(N)}_{ \sigma(\mathbf{a})}(\lambda \theta)}.$$
Integrating $N$-times, taking into account that $Y_\mathbf{a}$ must vanish at zero at order $N-1$, yields
$$ \forall \mathbf{a} \in \mathcal{A}^\mathbb{Z}, \forall \theta \in \mathbb{R}, \quad  Y_{\mathbf{a}}(\theta) =  {\varphi_{ \mathbf{a}}(\theta)} - P_\mathbf{a}(\theta) + \mu {Y_{ \sigma(\mathbf{a})}(\lambda \theta)} ,$$
where $P_\mathbf{a}(\theta)$ is the Taylor expansion of $\varphi$ at $\pi(\mathbf{a})$ of order $N-1$: $$ P_\mathbf{a}(\theta) = \sum_{j=0}^{N-1}  \frac{\varphi^{(j)}(\pi(\mathbf{a}))}{j!} \theta^j .$$
Setting $Z_\mathbf{a}(\theta) := X_\mathbf{a}-Y_\mathbf{a}$ finally yields $$ Z_{\mathbf{a}}(\theta) = P_\mathbf{a}(\theta) + \mu Z_\mathbf{a}(\lambda \theta), $$
which concludes the proof.
\end{proof}

We will quantify the oscillations of $X$ by considering the distance of $Z$ with the space of polynomials of degree $N$.

\begin{definition}
Define the distance between $Z_\mathbf{a}$ and the space of polynomials of degree less than $N$-1 as $$ \begin{array}[t]{lrcl}
\mathcal{D} : & \mathcal{A}^\mathbb{Z} & \longrightarrow & [0,\infty) \\
    & \mathbf{a}  & \longmapsto &  \underset{P \in \mathbb{R}_{N-1}[X]}{\inf} \| Z_\mathbf{a} - P \|_{L^\infty([-1,1])}. \end{array} $$
This function is continuous on the compact space $\mathcal{A}^\mathbb{Z}$, since the vector space $\mathbb{R}_{N-1}[X]$ is a finite dimensional subspace of $L^\infty([-1,1])$.
\end{definition}

\begin{lemma}
If $X$ is not analytic, there exists $\kappa_0>0$ such that for all $\mathbf{a} \in \mathcal{A}^\mathbb{Z}$,  $\mathcal{D}(\mathbf{a})\geq\kappa_0$.
\end{lemma}

\begin{proof}
The function $\mathcal{D}$ can not vanish on $\mathcal{A}^\mathbb{Z}$.
Indeed, if there was some $\mathbf{a} \in \mathcal{A}^\mathbb{Z}$ such that $\mathcal{D}(\mathbf{a})=0$, then it would mean that there exists a polynomial $P$ such that $Z_\mathbf{a}=P$ on $[-1,1]$. But then, since $X_\mathbf{a}=Y_\mathbf{a}+Z_\mathbf{a}$, it would follow that $X_\mathbf{a}$ is analytic on $[-1,1]$. Since $X_\mathbf{a}(\theta)=X(\pi(\mathbf{a})+\theta)$ and since $X$ is $1$-periodic, it would follow that $X$ is analytic on $\mathbb{R}$, which is false by hypothesis.
It then follow that $\mathcal{D}>0$ on the compact space $\mathcal{A}^\mathbb{Z}$.
This concludes the proof since $\mathcal{D}$ is continuous.
\end{proof}

Since our goal is to study the local oscillations of $X_\mathbf{a}$ near $0$, and not \emph{at} zero, we need to work a little bit more to ensure that this sup norm is not realized too close from zero. This will allow to meaningfully relate oscillation at different scales later, the sup norm being understood as giving us oscillations at a macroscopic scale.

\begin{lemma}
Suppose that $X$ is not analytic. There exists $\kappa >0$ such that, for all $\mathbf{a} \in \mathcal{A}^\mathbb{Z}$, we have:
$$ \forall  P \in \mathbb{R}_{N-1}[X], \ \exists \theta_0 \in [-1,1] \setminus [-2\kappa,2\kappa], \ \forall \theta \in [\theta_0-\kappa,\theta_0+\kappa], \quad |Z_\mathbf{a}(\theta)-P(\theta)| \geq \kappa. $$
\end{lemma}

\begin{proof}
We saw earlier that there exists $\gamma \in (0,1)$ such that the maps $X_\mathbf{a}$ are all at least $\gamma$-Hölder, and this uniformly in $\mathbf{a}$. By construction, the same holds true for $Z_\mathbf{a}$. Fix one such $\gamma \in (0,1)$, and set $R := 1+ \sup_{\mathbf{a} \in \mathcal{A}^\mathbb{Z}} \|Z_\mathbf{a}\|_{C^\gamma([-1,1])} \geq 1$. Set, for $i \in \{1,\dots,N\}$: $I_i := [1/2+i/N,1/2+i/N+1/2N] \subset [1/2,1]$.   \\
Since $\mathbb{R}_{N-1}[X]$ is a finite dimensional real vector space, every norms are equivalent. It follows that there exists $C \geq 1$ such that $$ \forall i \in  \{ 1, \dots, N \}, \  C^{-1} \|P\|_{C^\alpha([-1,1])} \leq \|P\|_{L^\infty([-1,1])} \leq C \|P\|_{L^\infty(I_i)}.$$
Denote by $\kappa_0$ the constant given by the previous lemma, set $M := 2CR$ and set $\kappa := \frac{\kappa_0}{N} (6R C^2 )^{-1/\gamma}$, and let us prove the desired property for this choice of $\kappa$. Choose some $\mathbf{a} \in \mathcal{A}^\mathbb{Z}$, and choose some $P \in \mathbb{R}_{N-1}[X]$. If $\|P\|_{L^\infty([-1,1])} \geq M$, then since $\|P\|_{L^\infty(I_i)} \geq C^{-1} \|P\|_{L^\infty([-1,1])}$, there exists for each $i$ a number $t_i \in I_i$ such that $P(t_i) \geq C^{-1}M = 2R$. There is $N$ of these numbers, since $P$ has degree at most $N-1$, there must exists $i_0 \in \{1,\dots,N-1\}$ such that $P$ is monotonous on $[t_i,t_{i+1}]$ (which is an interval of length $\geq \kappa$.) Hence, since $\|Z_\mathbf{a}\|_{L^\infty([-1,1])} \leq R$:
$$\forall t \in [t_i,t_{i+1}], \ |P(t)-Z_\mathbf{a}(t_i)| \geq 2R-\|Z_\mathbf{a}\|_{L^\infty([1/2,1])} \geq R \geq \kappa. $$

Now suppose that $\|P\|_{L^\infty([-1,1])}\leq M$. Let $t_0 \in [-1,1]$ be a number realizing the sup norm of $\mathcal{D}(\mathbf{a}) \geq \kappa_0$. Then for every $t \in [t_0-3\kappa,t_0+3\kappa]$, 
$$ |P(t)-Z_\mathbf{a}(t)| = |P(t_0)-Z_\mathbf{a}(t_0)| - \|P-Z_\mathbf{a}\|_{C^{\gamma}} (3\kappa)^\gamma $$ $$\geq \kappa_0-(CM+R)(3\kappa)^\gamma \geq \kappa_0/2 \geq \kappa, $$
which concludes the proof.
\end{proof}


To study the local oscillations of $X$, we introduce the zoomed and rescaled functions, for $\mathbf{a} \in \mathcal{A}^\mathbb{Z}$ and $n \geq 0$:
$$ X_\mathbf{a}^{\langle n \rangle}(\theta) := \mu^{-n} X_{\mathbf{a}}(\theta \lambda^{-n}) \quad , \quad Z_\mathbf{a}^{\langle n \rangle}(\theta) := \mu^{-n} Z_{\mathbf{a}}(\theta \lambda^{-n}). $$

\begin{lemma}
Suppose that $X$ is not analytic. There exists $\kappa >0$ such that the following holds. For all $\mathbf{a} \in \mathcal{A}^\mathbb{Z}$, for all $n \geq 0$:
$$ \forall P \in \mathbb{R}_{N-1}[X], \ \exists t_0 \in [-1,1] \setminus [-2\kappa,2\kappa], \ \forall t \in [t_0-\kappa,t_0+\kappa], \quad |Z_\mathbf{a}^{\langle n \rangle}(t)-P(t)| \geq \kappa. $$
\end{lemma}

\begin{proof}
The autosimilarity relation
$$ Z_\mathbf{a}(\theta) = P_\mathbf{a}(\theta)+ \mu Z_{\sigma(\mathbf{a})}(\lambda \theta) $$
can be rewriten as $$ Z_\mathbf{a}^{\langle 1 \rangle}(\theta) = \mu^{-1} Z_{\mathbf{a}}(\theta \lambda^{-1}) = \mu^{-1} P_\mathbf{a}(\theta \lambda^{-1})+ Z_{\sigma(\mathbf{a})}(\theta) $$
so that $ Z_\mathbf{a}^{\langle 1 \rangle}(\theta) = Z_{\sigma(\mathbf{a})} \text{ mod } \mathbb{R}_{N-1}[X]. $
Iterating yields, for all $n \geq 0$
$$ Z_\mathbf{a}^{\langle n \rangle} = Z_{\sigma^{n}(\mathbf{a})} \text{ mod } \mathbb{R}_{N-1}[X] ,$$
which gives the desired property, by the previous lemma applied to $Z_{\sigma^{n}(\mathbf{a})}$.
\end{proof}


\begin{lemma}\label{lem:osc}
Suppose that $X$ is not analytic. There exists $n_0 \in \mathbb{N}$ and $\kappa >0$ such that the following holds. For all $\mathbf{a} \in \mathcal{A}^\mathbb{Z}$, for all $n \geq n_0$:
$$ \forall P \in \mathbb{R}_{N-1}[X], \ \exists t_0 \in [-1,1] \setminus [-2\kappa,2\kappa], \forall t \in [t_0-\kappa,t_0+\kappa], \quad |X_\mathbf{a}^{\langle n \rangle}(t)-P(t_0)| \geq \kappa. $$
\end{lemma}


\begin{proof}
Suppose that $X$ is not analytic and denote by $\kappa>0$ the constant given by the previous lemma. Now recall that $X_\mathbf{a}=Y_\mathbf{a}+Z_\mathbf{a}$, with $Y_\mathbf{a}(\theta)=\mathcal{O}(\theta^N)$. It follows that
$$ X_\mathbf{a}^{\langle n \rangle}(\theta) = \mu^{-n} X_\mathbf{a}(\theta \lambda^{-n}) = \mu^{-n} Y_\mathbf{a}(\theta \lambda^{-n}) + \mu^{-n} Z_\mathbf{a}^{\langle n \rangle}(\theta \lambda^{-n}) = Z_\mathbf{a}^{\langle n \rangle}(\theta) + \mathcal{O}( (\mu \lambda^N)^{-n} ), $$
which implies the desired result for $\kappa$ replaced by $\kappa/2$, as soon as $\mathcal{O}((\mu \lambda^N)^{-n_0}) \leq \kappa/2$.
\end{proof}

The previous result is the core statement about the local oscillations of Weierstrass functions. A slightly weakened version gives:

\begin{theorem}{\label{thm:osc}}
Suppose that $X$ is not analytic. There exists $n_0 \geq 1$ and $\kappa>0$ such that the following holds. For all $x \in \mathbb{S}^1$, for all $n \geq n_0$:
$$ \forall P \in \mathbb{R}_{N-1}[X], \ \exists \theta_n \in \lambda^{-n}([-1,1] \setminus [-\kappa,\kappa]), \quad |X(x+\theta_n)-P(\theta_n)| \geq \kappa \mu^n$$
Notice that $\mu^n = \lambda^{-\alpha n} \simeq \theta_n^{\alpha}$.
\end{theorem}

Let us define a space of functions regular at $x$, for $\gamma \in (0,\infty)$, by:
$$ C^{o(\gamma)}_x(\mathbb{S}^1,\mathbb{R}) := \{ f \in C^0(\mathbb{S}^1,\mathbb{R}) \ | \ \exists P \in \mathbb{R}_{\lfloor \gamma \rfloor}[X], \ |X(x+t)-P(t)| = o(t^\gamma) \}. $$
For example, $C^{o(1)}_x$ is the space of differentiable functions at $x$.

\begin{lemma}
Recall that $\alpha \in (0,\infty)$ is defined by the relation $\mu \lambda^\alpha=1$. \\ 
If there exists $x \in \mathbb{S}^1$ such that $X \in  C^{o(\alpha)}_x(\mathbb{S}^1,\mathbb{R})$, then $X$ is actually analytic. 
\end{lemma}

\begin{proof}
Suppose that $X \in \bigcup_{x \in \mathbb{S}^1} C^\alpha_x$. This means that there exists $x$ and $P \in \mathbb{R}_{\lfloor \alpha \rfloor}[X] \subset \mathbb{R}_{N-1}[X]$ such that $|X(x+t)-P(t)| = o(t^\alpha)$. This is directly contradicting the conclusion of \ref{thm:osc}, which construct a sequence of points $t_n \rightarrow 0$ such that $|
X(x+t_n)-P(t_n)| \geq \kappa t_n^\alpha$.
\end{proof}

\begin{remark}
As an example, recall that the function $X(\theta) = \sum_{n=0}^\infty 2^{-n}{\cos(2\pi 2^n \theta)}$ is not analytic. We have $\alpha=1$, so it follows that $X \in C^{1-}$, but is not differentiable at any point (not even at $0$, where the formal derivative in the sense of distributions $\sum_n \sin(2 \pi 2^n \theta)$ seem to vanish).
\end{remark}



\subsection{Frostman regularity of the occupation measure and dispersion}

We will show now that the occupation measure is always Frostman.

\begin{theorem}\label{thm:c1+alpharegularity}
If $X$ is not analytic, then there exists $\gamma(\varphi,\lambda,\mu)>0$ such that for any  interval $I \subset \mathbb{R}$ small enough, for all $P \in \mathbb{R}_{N-1}[X]$ with $N \geq 1$ the first integer such that $\mu \lambda^N >1$, we have 
$$\lambda( \theta \in {\mathbb{S}^1}  , \ |W(\theta)-P(\theta)| \leq \sigma) \lesssim \sigma^\gamma ,$$
uniformly in $P$.
\end{theorem}

\begin{remark}
In particular, for any non-analytic Weierstrass map, $W_* \lambda$ never have atoms.
\end{remark}

\begin{proof}
Suppose that $X$ is not analytic. Denote $n_0$ and $\kappa$ the constants given by Lemma \ref{lem:osc}. Fix $n_1 \geq n_0$ and choose $q\geq 1$ large enough so that $\lambda^{-q}\leq \kappa/4$. Then, let $P \in \mathbb{R}[X]$ of degree $\leq N-1$, and let $\sigma \in (0,\kappa)$. Set $n(\sigma)$ such that $\kappa\mu^{n(\sigma)(n_1+q)}\geq \sigma$. We cut the circle $\mathbb{S}^1$ into $\lambda^{n_1+q}$ equal parts segments:
$$
\mathbb{S}^1=\bigcup_{\mathbf{a}_1\in\mathcal{A}^{n_1+q}}I_{\mathbf{a}_1}.
$$
Lemma \ref{lem:osc} gives an interval of length $2\kappa$ in the rescaled coordinates on which $|X(\theta)-P(\theta)|\geq\kappa$. Since $\lambda^{-q}\leq\kappa/4$, this interval contains at least one cylinder $I_{\widehat{\mathbf a}_1}$ of generation $n_1+q$. Hence
$$
\{\theta\in\mathbb{S}^1:\ |X(\theta)-P(\theta)|\leq\sigma\}
\subset
\bigcup_{\substack{\mathbf a_1\in\mathcal A^{n_1+q}\\ \mathbf a_1\neq\widehat{\mathbf a}_1}}
I_{\mathbf a_1}.
$$
Now, we iterate the argument on each $I_{\mathbf a_1}$, that we cut into $\lambda^{n_1+q}$ equal parts:
$$
I_{\mathbf a_1}=\bigcup_{\mathbf a_2\in\mathcal A^{n_1+q}}I_{\mathbf a_1\mathbf a_2}.
$$
Lemma \ref{lem:osc}, applied at the suitable scale, gives for each $\mathbf a_1$ an interval of relative length $2\kappa$ on which $|X(\theta)-P(\theta)|\geq\kappa\mu^{n_1+q}>\sigma$. By the choice of $q$, this interval contains a whole descendant cylinder $I_{\mathbf a_1\widehat{\mathbf a}_2}$ for some $\widehat{\mathbf a}_2=\widehat{\mathbf a}_2(\mathbf a_1)\in\mathcal A^{n_1+q}$. It follows that
$$
\{\theta\in\mathbb{S}^1:\ |X(\theta)-P(\theta)|\leq\sigma\}
\subset
\bigcup_{\substack{\mathbf a_1\in\mathcal A^{n_1+q}\\ \mathbf a_1\neq\widehat{\mathbf a}_1}}
\ \bigcup_{\substack{\mathbf a_2\in\mathcal A^{n_1+q}\\ \mathbf a_2\neq\widehat{\mathbf a}_2(\mathbf a_1)}}
I_{\mathbf a_1\mathbf a_2}.
$$
This argument can be iterated $n(\sigma)$ times, yielding a nested union:
$$
\{\theta\in\mathbb{S}^1:\ |X(\theta)-P(\theta)|\leq\sigma\}
\subset
\bigcup_{\mathbf a_1\neq\widehat{\mathbf a}_1}
\bigcup_{\mathbf a_2\neq\widehat{\mathbf a}_2(\mathbf a_1)}
\dots
\bigcup_{\mathbf a_{n(\sigma)}\neq\widehat{\mathbf a}_{n(\sigma)}(\mathbf a_1\dots\mathbf a_{n(\sigma)-1})}
I_{\mathbf a_1\dots\mathbf a_{n(\sigma)}},
$$
where each $\mathbf a_j\in\mathcal A^{n_1+q}$. After $n(\sigma)$ steps there are at most $(\lambda^{n_1+q}-1)^{n(\sigma)}$ surviving cylinders, each of length $\lambda^{-(n_1+q)n(\sigma)}$. We can thus directly bound the measure by
$$
\lambda\Big(\theta\in\mathbb{S}^1,\ |X(\theta)-P(\theta)|\leq\sigma\Big)
\leq
\sum_{\mathbf a_1\neq\widehat{\mathbf a}_1}\dots
\sum_{\mathbf a_{n(\sigma)}\neq\widehat{\mathbf a}_{n(\sigma)}(\mathbf a_1\dots\mathbf a_{n(\sigma)-1})}
\lambda(I_{\mathbf a_1\dots\mathbf a_{n(\sigma)}})
$$
$$
\leq
(1-\lambda^{-(n_1+q)})^{n(\sigma)}
\simeq
\sigma^{\frac{\ln(1-\lambda^{-(n_1+q)})}{(n_1+q)\ln(\mu)}}.
$$
\end{proof}




Using the $C^{1+}$ Van Der Corput Lemma, we recover the dispersion result in the range $\mu \lambda < 1$.

\begin{corollary}
Let $X$ be a Weierstrass function, suppose that $\mu \lambda < 1$. Then there exists $\rho \in (0,1)$ such that
$$ \sup_{n \in \mathbb{Z}} \Big| \int_0^1 e^{2 i \pi t X(\theta)} e^{-2 i \pi n \theta} d\theta \Big| \lesssim |t|^{-\rho}, \quad t \to \infty.$$
\end{corollary}

\subsection{Density of the derivatives of the partial sums in the critical regime}

In the critical case where $\mu \lambda = 1$, we know that $X \in C^{1-} \setminus C^1$. In this case, instead of studying the derivative of $X$, we can study the derivative of partial sums. The goal of this section is to prove the following result.

\begin{proposition}\label{prop:density}
Let $\varphi \in C^\omega(\mathbb{S}^1,\mathbb{R})$, analytic, 1-periodic and non-constant. Suppose that there exists no analytic map such that $\varphi = \frac{1}{\lambda} \psi(\lambda \theta)-\psi(\theta)$. Then, for almost every $\theta \in \mathbb{S}^1$, the set 
$ \Big\{ \sum_{n=0}^N \varphi'(\lambda^n \theta) \ | \ N \in \mathbb{N} \Big\} $
is dense in $\mathbb{R}$.
\end{proposition}

This implies:

\begin{corollary}
Suppose $\mu\lambda=1$ and suppose $X$ is not analytic. Then, for
a.e. $\theta\in\mathbb S^1$,$ \left\{\sum_{n=0}^N \varphi'(\lambda^n\theta):N\geq0\right\}$
is dense in $\mathbb R$.
\end{corollary}

To do so, we need to recall some results from infinite ergodic theory. Fix some analytic map $\phi \in C^\omega(\mathbb{S}^1,\mathbb{R})$, and  consider the skew-product dynamical system $T_\phi:\mathbb{S}^1 \times \mathbb{R} \rightarrow \mathbb{S}^1 \times \mathbb{R} $ given by $T_\phi(\theta,y)=(\lambda \theta,y+\phi(x))$. This leaves invariant the infinite-mass measure $m := \lambda_\mathbb{S}^1 \otimes \lambda_{\mathbb{R}}$.

\begin{definition}[\cite{Gu89}]
We say that $\phi$ is not strictly aperiodic if there exists a Hölder function $\psi \in C^0(\mathbb{S}^1,\mathbb{C}^*)$ and $(\alpha,\lambda) \in \mathbb{R} \times \mathbb{R}^*$ such that $e^{i \lambda \phi(\theta)} = e^{i \alpha} \frac{\psi(\lambda \theta)}{\psi(\theta)} $.
\end{definition}

\begin{lemma}[\cite{Gu89}, Section 2.3]\label{lem:ap-erg}
If $\phi$ is strictly aperiodic and $\int_{\mathbb{S}^1} \phi =0$, then $(T_\phi,\mathbb{S}^1 \times \mathbb{R},m)$ is ergodic.
\end{lemma}

\begin{comment}
\begin{proof}
See \url{https://www.cambridge.org/core/services/aop-cambridge-core/content/view/2D210905909FEAA898D0538C83D76FE2/S0143385700005083a.pdf/proprietes-ergodiques-en-mesure-infinie-de-certains-systemes-dynamiques-fibres.pdf} , Corollary 3 in Section 2.3.
\end{proof}
\end{comment}

\begin{lemma}[\cite{Aa97}, Corollary 8.1.5]\label{lem:cons}
If $\int_{\mathbb{S}^1} \phi = 0$, then $T_\phi$ is conservative.
\end{lemma}



\begin{lemma}[\cite{Aa97}, Corollary 1.2.3]\label{lem:dens}
Suppose that $T_\phi$ is conservative and ergodic. Then for any continuous function $f:\mathbb{S}^1 \times \mathbb{R} \rightarrow \mathbb{R}$, $$ \overline{\{ f(T_\phi^n(\theta,y)) \ , \ n \geq 0\} } = \text{spt}(f_* m). $$
\end{lemma}

With that, we are ready to prove Proposition \ref{prop:density}.

\begin{proof}
Set $\phi = \varphi'$. Obviously $\int_{\mathbb{S}^1} \phi = 0$. Furthermore, $\phi$ is strictly aperiodic. Because if it was not, then $\phi$ would be cohomologous to a constant, but since $\int \phi =0$, it would actually be cohomologous to zero. Integrating the cohomology relation would give that $X$ is $C^1$, which is not possible by hypothesis, by Theorem \ref{thm:regWeierstrass}. By Lemma \ref{lem:ap-erg} and \ref{lem:cons}, the dynamical system $(T_\phi,\mathbb{S}^1 \times \mathbb{R}, m)$ is ergodic and conservative. By Lemma \ref{lem:dens}, choosing $f(\theta,y)=y$, we have $$ \overline{\{ \sum_{0\leq k \leq n} \phi(\lambda^k \theta) , \ , \ n \geq 0 \}} =  \overline{\{ f(T^n(\theta,y)) \ , \ n \geq 0\} } = \text{spt}(f_* m) = \mathbb{R}. $$
\end{proof}


\begin{comment}

As a consequence, we find the following corollary. \textbf{Actually, this should also give pure singularity for $(X e_m)_* (d\theta)$.}

\begin{corollary}
Suppose $\mu \lambda=1$ and suppose $X$ is not analytic. Then $X_* d\theta$ is purely singular.
\end{corollary}

\begin{proof}
If $X$ is not analytic then there exists no Hölder $\psi \in C^0(\mathbb{S}^1,\mathbb{R})$ such that $\varphi' = \psi \circ f - \psi$. It follows from \ref{prop:density} that for a.e. $\theta \in \mathbb{S}^1$, the set \{ $\sum_n \varphi'(\lambda^n \theta) , \ , \ n \geq 0$  \} is dense in $\mathbb{R}$. \\


\textbf{Hence the pushforward measure is purely singular by some argument. ?}

\end{proof}

\end{comment}



\section{On the large lacunarity limit} \label{ap:large-lacunarity}

Fix $0<\mu<1$, and let $\varphi\in C^0(\mathbb S^1,\mathbb R)$. Suppose that there exists an interval $I$ on which $\varphi$ is analytic and non-constant. For example, one can choose $\varphi$ to be constant equal to $1$ for $1/10<x<4/10$, constant equal to $-1$ for $6/10<x<9/10$, and piecewise affine on the remaining parts of the circle so that the resulting function is continuous. For each integer $\lambda\geq2$, recall that $W_{\mu,\lambda}(\theta)=\sum_{n=0}^\infty\mu^n\varphi(\lambda^n\theta)$. Since increasing the lacunarity produces arbitrarily fast Fourier decay and dispersion, it is natural to ask what the limiting object is as $\lambda\to\infty$. We show that in this limit the terms $\varphi(\lambda^n\theta)$ behave asymptotically as independent random variables. More precisely, the limiting measure below is the law of the random series $\sum_{n=0}^\infty\mu^nX_n$, where $(X_n)_{n\geq0}$ are i.i.d.\ with common distribution $\varphi_*d\theta$. Thus, in the large-lacunarity limit, the deterministic Weierstrass series behaves like a random series with i.i.d.\ coefficients. For functions $\varphi$ taking predominantly the values $\pm1$, as in the example above, this limit is closely related to a Bernoulli convolution.

\begin{theorem}
For any $0 < \mu < 1$ the weak limit
\[
m_\mu:=\lim_{\lambda\to\infty}(W_{\mu,\lambda})_*(d\theta)
\]
exists and is the distribution of $\sum_{n=0}^\infty\mu^nX_n$, where the $X_n$ are i.i.d.\ with distribution $\varphi_*d\theta$. Moreover, $|\widehat{m_\mu}(\xi)|\lesssim|\xi|^{-\rho}$ as $|\xi|\to\infty$ for some $\rho>0$.
\end{theorem}

\begin{proof}
Up to a subsequence, a weak limit $m_\mu$ exists. Let us show that the limit is unique. We have, for a fixed $t \in \mathbb{R}$:
$$ \widehat{W_* (d\theta)}(t) = \int_{\mathbb{S}^1} e^{i t W_{\mu,\lambda}(\theta)} d\theta $$
$$ = \int_{\mathbb{S}^1} e^{i t \varphi(\theta)} e^{it \mu W_{\mu,\lambda}(\lambda \theta)} d\theta $$
$$ = \int_{\mathbb{S}^1} \mathcal{L}_\lambda( e^{i t \varphi})(\theta) e^{it \mu W_{\mu,\lambda}(\theta)} d\theta. $$
Now, $e^{i t \varphi}$ is continuous. If we denote by $\varepsilon_t(\cdot)$ a modulus of continuity, we can write
$$ \Big| \mathcal{L}_\lambda(e^{it \varphi})(\theta) - \int_{\mathbb{S}^1} e^{i t \varphi} d\theta \Big| \leq \varepsilon_t(1/\lambda) .$$
Hence
$$ \widehat{W_*(d\theta)}(t) = \int_{\mathbb{S}^1} e^{it \varphi} \int_{\mathbb{S}^1} e^{it\mu W_{\mu,\lambda}} + O(\varepsilon_t(1/\lambda)) $$
$$ = F(t) \widehat{W_* (d\theta)}(\mu t) + O(\varepsilon_t(1/\lambda)), $$
where $F(t) := \int_{\mathbb{S}^1} e^{i t \varphi}$. Taking the limit of $\lambda \rightarrow +\infty$ along our subsequence yields the relation
$$ \widehat{m_\mu}(t) = F(t)\widehat{m_\mu}(t\mu) $$
so that
$$ \widehat{m_\mu}(t) = \prod_{n=0}^\infty F(t \mu^n). $$
If $(X_n)_{n\geq0}$ are i.i.d.\ with distribution $\varphi_*d\theta$, then $F(t)=\mathbb E(e^{itX_0})$, and by independence
$$ \mathbb E\left(e^{it\sum_{n=0}^\infty\mu^nX_n}\right)=\prod_{n=0}^\infty F(t\mu^n). $$
Hence $m_\mu$ is precisely the distribution of $\sum_{n=0}^\infty\mu^nX_n$. This also proves that the limit is unique, so the weak limit indeed exists.

Now let us bound $F$. We have
$$ F(t) = \int_{\mathbb{S}^1} e^{it \varphi(\theta)} d\theta = F_1(t) + F_2(t) $$
where $F_1(t) = \int_{I} e^{it \varphi}$, where $\varphi$ is analytic and non-constant, and $F_2(t) = \int_{\mathbb{S}^1 \setminus I} e^{it \varphi}$ is the remaining part. We have the bounds $F_1(t) \lesssim t^{-\rho}$ for some $\rho$, and $|F_2(t)| \leq |\mathbb{S}^1 \setminus I| := 1-|I|$. In particular, if $t$ is large enough $t \geq T$, we can write $|F(t)| \leq 1-|I|/2$. Then, letting $N(t)$ be the largest integer such that $t \mu^{N(t)} \geq T$, we find
$$ |\widehat{m_\mu}(t)| \leq \prod_{n=1}^{N(t)} |F(t \mu^n)| \leq (1-|I|/2)^{N(t)} \lesssim t^{- |\ln(1-|I|/2)/\ln(\mu)|} .$$
\end{proof}

\begin{thebibliography}{10}

\bibitem{Aa97}
J.~Aaronson,
\emph{An Introduction to Infinite Ergodic Theory},
Mathematical Surveys and Monographs, Vol.~50,
American Mathematical Society, Providence, RI, 1997.

\bibitem{AlgomBenOvadiaHertzShannon}
A.~Algom, S.~Ben Ovadia, F.~Rodriguez Hertz and M.~Shannon,
\emph{Fourier decay and non-decay for pseudo-affine self-conformal measures},
arXiv:2607.22001, 2026.

\bibitem{ACWW}
A.~Algom, Y.~Chang, M.~Wu and Y.-L.~Wu,
\emph{Van der Corput and metric theorems for geometric progressions for self-similar measures},
Math. Ann. \textbf{393} (2025), 183--214.

\bibitem{AlgomHertzWangLog}
A.~Algom, F.~Rodriguez Hertz and Z.~Wang,
\emph{Logarithmic Fourier decay for self-conformal measures},
J. Lond. Math. Soc. (2) \textbf{106} (2022), no.~2, 1628--1661.

\bibitem{AlgomHertzWang}
A.~Algom, F.~Rodriguez Hertz and Z.~Wang,
\emph{Polynomial Fourier decay and a cocycle version of Dolgopyat's method for self-conformal measures},
arXiv:2306.01275.

\bibitem{AndersonHousworthPitt}
J.~M.~Anderson, E.~A.~Housworth and L.~D.~Pitt,
\emph{The spectral theory of multiplication operators and the recurrence properties for nondifferentiable functions in the Zygmund class},
Mathematika \textbf{39} (1992), 136--151.

\bibitem{AvilaJitomirskaya}
A.~Avila and S.~Jitomirskaya,
\emph{Almost localization and almost reducibility},
J. Eur. Math. Soc. \textbf{12} (2010), no.~1, 93--131.

\bibitem{BakerBanaji}
S.~Baker and A.~Banaji,
\emph{Polynomial Fourier decay for fractal measures and their pushforwards},
Math. Ann. \textbf{392} (2025), no.~1, 209--261.

\bibitem{BakerKhalilSahlsten}
S.~Baker, O.~Khalil and T.~Sahlsten,
\emph{Fourier decay from $L^2$-flattening},
arXiv:2407.16699.

\bibitem{BakerSahlsten}
S.~Baker and T.~Sahlsten,
\emph{Spectral gaps and Fourier dimension for self-conformal sets with overlaps},
arXiv:2306.01389.

\bibitem{BanajiYu}
A.~Banaji and H.~Yu,
\textit{Fourier transform of nonlinear images of self-similar measures: quantitative aspects},
to appear in Peking Math. J., arXiv:2503.07508.

\bibitem{BanajiYu2}
A.~Banaji and H.~Yu,
\textit{Fourier transform of nonlinear images of self-similar measures: qualitative aspects},
arXiv:2606.09743, 2026.

\bibitem{BZ20}
D.~Bambusi and Z.~Zhao, \emph{Dispersive estimate for quasi-periodic Schr\"odinger
operators on 1-d lattices}, Adv. Math. \textbf{366} (2020), 107071.

\bibitem{BaranskiBaranyRomanowska}
K.~Bara\'nski, B.~B\'ar\'any and J.~Romanowska,
\emph{On the dimension of the graph of the classical Weierstrass function},
Adv. Math. \textbf{265} (2014), 32--59.

\bibitem{Bary}
N.~K.~Bary,
\emph{A Treatise on Trigonometric Series},
Vols.~I--II,
Pergamon Press, Oxford, 1964.

\bibitem{BM}
A.~Beurling and P.~Malliavin,
\emph{On Fourier transforms of measures with compact support},
Acta Math. \textbf{107} (1962), 291--309.

\bibitem{BD}
J.~Bourgain and S.~Dyatlov,
\emph{Fourier dimension and spectral gaps for hyperbolic surfaces},
Geom. Funct. Anal. \textbf{27} (2017), no.~4, 744--771.

\bibitem{BoutetdeMonvelSabri}
A.~Boutet de Monvel and M.~Sabri,
\emph{Ballistic transport in periodic and random media},
in \emph{From Complex Analysis to Operator Theory: A Panorama},
Oper. Theory Adv. Appl., Vol.~291,
Birkh\"auser/Springer, Cham, 2023, 163--216.

\bibitem{Buc}
Z.~Buczolich,
\emph{Occupation measure and level sets of the Weierstrass--Cellerier function},
in \emph{Recent Developments in Fractals and Related Fields},
Appl. Numer. Harmon. Anal.,
Birkh\"auser, Boston, 2010, 3--18.

\bibitem{BKM}
Z.~Buczolich, A.~K\"aenm\"aki and B.~Maga,
\emph{Level sets of prevalent Weierstrass functions},
arXiv:2507.15591.

\bibitem{CaceresNizama}
M.~O.~C\'aceres and M.~Nizama,
\emph{The quantum L\'evy walk},
J. Phys. A: Math. Theor. \textbf{43} (2010), no.~45, 455306.

\bibitem{ChenLiSuomala}
C.~Chen, B.~Li and V.~Suomala,
\emph{Fourier dimension of Mandelbrot multiplicative cascades},
Comm. Math. Phys. \textbf{406} (2025), no.~8, 182.

\bibitem{DamanikFillmanYoung}
D.~Damanik, J.~Fillman and G.~Young,
\emph{Optimal dispersion for discrete periodic Schr\"odinger operators},
arXiv:2505.14475.

\bibitem{DamanikGorodetski}
D.~Damanik and A.~Gorodetski,
\emph{Spectral and quantum dynamical properties of the weakly coupled Fibonacci Hamiltonian},
Comm. Math. Phys. \textbf{305} (2011), no.~1, 221--277.

\bibitem{DamanikGorodetskiYessen}
D.~Damanik, A.~Gorodetski and W.~Yessen,
\emph{The Fibonacci Hamiltonian},
Invent. Math. \textbf{206} (2016), no.~3, 629--692.

\bibitem{DamanikLukicYessen}
D.~Damanik, M.~Luki\'c and W.~Yessen,
\emph{Quantum dynamics of periodic and limit-periodic Jacobi and block Jacobi matrices with applications to some quantum many-body problems},
Comm. Math. Phys. \textbf{337} (2015), no.~3, 1535--1561.

\bibitem{DamanikMalinovitchYoung}
D.~Damanik, T.~Malinovitch and G.~Young,
\emph{What is ballistic transport?},
J. Spectr. Theory, published online, 2024,
DOI: 10.4171/JST/532.

\bibitem{DamanikTcheremchantsev}
D.~Damanik and S.~Tcheremchantsev,
\emph{Upper bounds in quantum dynamics},
J. Amer. Math. Soc. \textbf{20} (2007), no.~3, 799--827.

\bibitem{Do98}
D.~Dolgopyat,
\emph{On decay of correlations in Anosov flows},
Ann. of Math. (2) \textbf{147} (1998), no.~2, 357--390.

\bibitem{FillmanPeriodic}
J.~Fillman,
\emph{Ballistic transport for periodic Jacobi operators on $\mathbb Z^d$},
in \emph{From Operator Theory to Orthogonal Polynomials, Combinatorics, and Number Theory},
Oper. Theory Adv. Appl., Vol.~285,
Birkh\"auser/Springer, Cham, 2021, 57--68.

\bibitem{Firsova}
N.~Firsova,
\emph{On the time decay of a wave packet in a one-dimensional finite band periodic lattice},
J. Math. Phys. \textbf{37} (1996), 1171--1181.

\bibitem{FoucheMukeru}
W.~L.~Fouch\'e and S.~Mukeru,
\emph{On local times of Martin-L\"of random Brownian motion},
Theoret. Comput. Sci. \textbf{979} (2023), 114199.

\bibitem{GarbanVargas}
C.~Garban and V.~Vargas,
\emph{Harmonic analysis of Gaussian multiplicative chaos on the circle},
Probab. Theory Related Fields, published online, 2026,
DOI: 10.1007/s00440-026-01497-7.

\bibitem{GemanHorowitz}
D.~Geman and J.~Horowitz,
\emph{Occupation densities},
Ann. Probab. \textbf{8} (1980), no.~1, 1--67.

\bibitem{Gu89}
Y.~Guivarc'h,
\emph{Propri\'et\'es ergodiques, en mesure infinie, de certains syst\`emes dynamiques fibr\'es},
Ergodic Theory Dynam. Systems \textbf{9} (1989), no.~3, 433--453.

\bibitem{HarangLing}
F.~A.~Harang and C.~Ling,
\emph{Regularity of local times associated with Volterra--L\'evy processes
and path-wise regularization of stochastic differential equations},
J. Theoret. Probab. \textbf{35} (2022), 1706--1735.

\bibitem{HaroPuig}
A.~Haro and J.~Puig,
\emph{A Thouless formula and Aubry duality for long-range Schr\"odinger skew-products},
Nonlinearity \textbf{26} (2013), no.~5, 1163--1187.

\bibitem{Hochman}
M.~Hochman,
\emph{On self-similar sets with overlaps and inverse theorems for entropy},
Ann. of Math. (2) \textbf{180} (2014), no.~2, 773--822.

\bibitem{HughesShlesingerMontroll}
B.~D.~Hughes, M.~F.~Shlesinger and E.~W.~Montroll,
\emph{Random walks with self-similar clusters},
Proc. Natl. Acad. Sci. USA \textbf{78} (1981), no.~6, 3287--3291.

\bibitem{Hunt}
B.~R.~Hunt,
\emph{The Hausdorff dimension of graphs of Weierstrass functions},
Proc. Amer. Math. Soc. \textbf{126} (1998), no.~3, 791--800.

\bibitem{ImkellerEtAl}
P.~Imkeller, O.~Menoukeu Pamen, G.~dos Reis and A.~R\'eveillac,
\emph{Rough Weierstrass functions and dynamical systems: the smoothness of the SBR measure},
Pure Appl. Funct. Anal. \textbf{7} (2022), no.~4, 1273--1305.

\bibitem{ImkellerPamen}
P.~Imkeller and O.~Menoukeu Pamen,
\emph{Takagi type functions and dynamical systems: the smoothness of the SBR measure and the existence and smoothness of local time},
arXiv:2107.07185.

\bibitem{JordanSahlsten}
T.~Jordan and T.~Sahlsten,
\emph{Fourier transforms of Gibbs measures for the Gauss map},
Math. Ann. \textbf{364} (2016), no.~3--4, 983--1023.

\bibitem{Kahane}
J.-P.~Kahane,
\emph{Some Random Series of Functions},
2nd ed.,
Cambridge Studies in Advanced Mathematics, Vol.~5,
Cambridge University Press, Cambridge, 1985.

\bibitem{KaufmanCF}
R.~Kaufman,
\emph{Continued fractions and Fourier transforms},
Mathematika \textbf{27} (1980), no.~2, 262--267.

\bibitem{Kaufman}
R.~Kaufman,
\emph{On Bernoulli convolutions},
in \emph{Conference in Modern Analysis and Probability},
Contemp. Math., Vol.~26,
American Mathematical Society, Providence, RI, 1984, 217--222.

\bibitem{Keller}
G.~Keller,
\emph{An elementary proof for the dimension of the graph of the classical Weierstrass function},
Ann. Inst. Henri Poincar\'e Probab. Stat. \textbf{53} (2017), no.~1, 169--181.

\bibitem{KernerPostSabriTaufer}
J.~Kerner, O.~Post, M.~Sabri and M.~T\"aufer,
\emph{The curious spectra and dynamics of non-locally finite crystals},
Comm. Math. Phys. \textbf{406} (2025), no.~7, 169.

\bibitem{KillipSimon}
R.~Killip and B.~Simon,
\emph{Sum rules for Jacobi matrices and their applications to spectral theory},
Ann. of Math. (2) \textbf{158} (2003), no.~1, 253--321.

\bibitem{Kinney}
J.~R.~Kinney,
\emph{Note on a singular function of Minkowski},
Proc. Amer. Math. Soc. \textbf{11} (1960), no.~5, 788--794.

\bibitem{Leclerc}
G.~Leclerc,
\emph{On oscillatory integrals with H\"older phases},
Discrete Contin. Dyn. Syst. \textbf{44} (2024), no.~1, 263--280.

\bibitem{Leclerc2}
G.~Leclerc,
\emph{Fourier decay of equilibrium states and the Fibonacci Hamiltonian},
arXiv:2507.23731.

\bibitem{LeclercBunched}
G.~Leclerc,
\emph{Fourier decay of equilibrium states for bunched attractors},
Nonlinearity \textbf{37} (2024), no.~9, 095020.

\bibitem{LeclercHyperbolic}
G.~Leclerc,
\emph{Fourier decay of equilibrium states on hyperbolic surfaces},
arXiv:2307.10755.

\bibitem{LeclercJulia}
G.~Leclerc,
\emph{Julia sets of hyperbolic rational maps have positive Fourier dimension},
Comm. Math. Phys. \textbf{397} (2023), no.~2, 503--546.

\bibitem{Ledrappier}
F.~Ledrappier,
\emph{On the dimension of some graphs},
Contemp. Math. \textbf{135} (1992), 285--293.

\bibitem{LequenSahlsten}
F.~Lequen and T.~Sahlsten,
\emph{Quantitative Fourier decay for Patterson--Sullivan measures of dimension larger than $1/2$},
arXiv:2607.18010.

\bibitem{Li}
J.~Li,
\emph{Decrease of Fourier coefficients of stationary measures},
Math. Ann. \textbf{372} (2018), no.~3--4, 1189--1238.

\bibitem{LiENS}
J.~Li,
\emph{Fourier decay, renewal theorem and spectral gaps for random walks on split semisimple Lie groups},
Ann. Sci. \'Ec. Norm. Sup\'er. (4) \textbf{55} (2022), no.~6, 1613--1686.

\bibitem{LiSahlstenAffine}
J.~Li and T.~Sahlsten,
\emph{Fourier transform of self-affine measures},
Adv. Math. \textbf{374} (2020), 107349.

\bibitem{LiSahlsten}
J.~Li and T.~Sahlsten,
\emph{Trigonometric series and self-similar sets},
J. Eur. Math. Soc. \textbf{24} (2022), no.~1, 341--368.

\bibitem{LindenstraussVarju}
E.~Lindenstrauss and P.~P.~Varj\'u,
\emph{Random walks in the group of Euclidean isometries and self-similar measures},
Duke Math. J. \textbf{165} (2016), no.~6, 1061--1127.

\bibitem{ManticaTotik}
G.~Mantica and V.~Totik,
\emph{Regularity of Minkowski's question mark measure, its inverse and a class of IFS invariant measures},
J. Lond. Math. Soc. (2) \textbf{99} (2019), 707--732.

\bibitem{MiZhao1}
Y.~Mi and Z.~Zhao,
\emph{Dispersive estimate for two-periodic discrete one-dimensional Schr\"odinger operator},
J. Math. Anal. Appl. \textbf{485} (2020), no.~1, 123768.

\bibitem{MiZhao2}
Y.~Mi and Z.~Zhao,
\emph{Dispersive estimates for periodic discrete one-dimensional Schr\"odinger operators},
Proc. Amer. Math. Soc. \textbf{150} (2022), no.~1, 267--277.

\bibitem{MosqueraShmerkin}
C.~A.~Mosquera and P.~S.~Shmerkin,
\emph{Self-similar measures: asymptotic bounds for the dimension and Fourier decay of smooth images},
Ann. Acad. Sci. Fenn. Math. \textbf{43} (2018), 823--834.

\bibitem{Na05}
F.~Naud,
\emph{Expanding maps on Cantor sets and analytic continuation of zeta functions},
Ann. Sci. \'Ec. Norm. Sup\'er. (4) \textbf{38} (2005), no.~1, 116--153.

\bibitem{PP90}
W.~Parry and M.~Pollicott,
\emph{Zeta functions and the periodic orbit structure of hyperbolic dynamics},
Ast\'erisque \textbf{187--188} (1990).

\bibitem{QueffelecRamare}
M.~Queff\'elec and O.~Ramar\'e,
\emph{Analyse de Fourier des fractions continues \`a quotients restreints},
Enseign. Math. (2) \textbf{49} (2003), no.~3--4, 335--356.

\bibitem{Rapaport}
A.~Rapaport,
\emph{On the Rajchman property for self-similar measures on $\mathbb{R}^d$},
Adv. Math. \textbf{403} (2022), 108375.

\bibitem{Rezakhanlou}
F.~Rezakhanlou,
\emph{The packing measure of the graphs and level sets of certain continuous functions},
Math. Proc. Cambridge Philos. Soc. \textbf{104} (1988), no.~2, 347--360.

\bibitem{Romanowska}
J.~Romanowska,
\emph{Measure and Hausdorff dimension of randomized Weierstrass-type functions},
Nonlinearity \textbf{27} (2014), no.~4, 787--801.

\bibitem{Sahlsten}
T.~Sahlsten,
\emph{Fourier transforms and iterated function systems},
in \emph{Recent Developments in Fractals and Related Fields},
Trends in Mathematics, Vol.~F319,
Birkh\"auser, Cham, 2025, 297--346.

\bibitem{SahlstenStevens}
T.~Sahlsten and C.~Stevens,
\emph{Fourier transform and expanding maps on Cantor sets},
Amer. J. Math. \textbf{146} (2024), no.~4, 945--982.

\bibitem{Shen}
W.~Shen,
\emph{Hausdorff dimension of the graphs of the classical Weierstrass functions},
Math. Z. \textbf{289} (2018), no.~1--2, 223--266.

\bibitem{Shmerkin}
P.~Shmerkin,
\emph{On the exceptional set for absolute continuity of Bernoulli convolutions},
Geom. Funct. Anal. \textbf{24} (2014), no.~3, 946--958.

\bibitem{SiSC7}
B.~Simon,
\emph{Operators with singular continuous spectrum, VII. Examples with borderline time decay},
Comm. Math. Phys. \textbf{176} (1996), 713--722.

\bibitem{Solomyak}
B.~Solomyak,
\emph{On the random series $\sum \pm \lambda^n$ (an Erd\H{o}s problem)},
Ann. of Math. (2) \textbf{142} (1995), no.~3, 611--625.

\bibitem{SteinShakarchiComplex}
E.~M.~Stein and R.~Shakarchi,
\emph{Complex Analysis},
Princeton Lectures in Analysis, Vol.~II,
Princeton University Press, Princeton, NJ, 2003.

\bibitem{Tsujii}
M.~Tsujii,
\emph{Fat solenoidal attractors},
Nonlinearity \textbf{14} (2001), no.~5, 1011--1027.

\bibitem{TZ23}
M.~Tsujii and Z.~Zhang,
\emph{Smooth mixing Anosov flows in dimension three are exponentially mixing},
Ann. of Math. (2) \textbf{197} (2023), no.~1, 65--158.

\bibitem{Varju}
P.~P.~Varj\'u,
\emph{Absolute continuity of Bernoulli convolutions for algebraic parameters},
J. Amer. Math. Soc. \textbf{32} (2019), no.~2, 351--397.

\bibitem{Zyg32}
A.~Zygmund,
\emph{On lacunary trigonometric series},
Trans. Amer. Math. Soc. \textbf{34} (1932), 435--446.

\end{thebibliography}



\end{document}
